% Sensibilisation sur la nécessité de la mitose pour le maintien de l'identité biologique des organismes — SVT 1ère D — Cours signé OnBuch+
% Programme officiel MINESEC. Compiler avec : tectonic ND01-mitose-identite-biologique.tex
\newcommand{\DOCMATIERE}{SVT}
\newcommand{\DOCNIVEAU}{1ère D}
\newcommand{\DOCMODULE}{Module 1 : Le Monde Vivant}
\newcommand{\DOCLECON}{ND01}
\newcommand{\DOCTITRE}{Sensibilisation sur la nécessité de la mitose pour le maintien de l'identité biologique des organismes}
% OnBuch+ — préambule « cours signés » ENRICHI (compilé avec tectonic / XeLaTeX)
% Système visuel « Pop » d'OnBuch+ : fond crème (jamais blanc), bordures encre
% épaisses, ombres « dures » décalées, titres Archivo Black en capitales.
% Version MATHÉMATIQUES : graphes (pgfplots/tikz), boîtes pédagogiques
% (définition, propriété, méthode, attention, le savais-tu, exercice résolu,
% exercice, TP, bilan), page de garde et signature OnBuch+ ancrée en bas.
%
% Macros attendues dans main.tex AVANT \input{preamble} :
%   \DOCMATIERE  \DOCNIVEAU  \DOCMODULE  \DOCLECON (numéro)  \DOCTITRE
\documentclass[11pt]{article}

\usepackage{fontspec}
\usepackage[french]{babel}
\usepackage[a4paper,margin=2cm,top=3.4cm,bottom=2.8cm,headheight=1.4cm,footskip=1.3cm]{geometry}
\usepackage{amsmath,amssymb,amsfonts}
\usepackage{mathtools} % \xmapsto, \coloneqq... souvent utilisés par les modèles
\usepackage{cancel} % simplifications barrées (bilans de forces)
% \abs{z} (module, valeur absolue) et \norm{u} (norme), utilisés par les modèles
\providecommand{\abs}{}\let\abs\relax\DeclarePairedDelimiter{\abs}{\lvert}{\rvert}
\providecommand{\norm}{}\let\norm\relax\DeclarePairedDelimiter{\norm}{\lVert}{\rVert}
\usepackage[table]{xcolor} % [table] : fournit aussi \rowcolors (alternance de lignes)
\usepackage{pagecolor}
\usepackage{tikz}
\usetikzlibrary{arrows.meta,positioning,calc,shapes.geometric,shapes.misc,shapes.arrows,decorations.pathmorphing,decorations.pathreplacing,decorations.markings,patterns,shadows,mindmap,trees,fit,backgrounds,matrix,intersections,angles,quotes}
\usepackage{pgfplots}
\usepackage[version=4]{mhchem} % équations nucléaires \ce{^{235}_{92}U}
\usepackage[european,siunitx]{circuitikz} % schémas électriques
\pgfplotsset{compat=1.18}
\usepgfplotslibrary{fillbetween,groupplots} % aires entre courbes (calcul intégral)
\usepackage{enumitem}
\usepackage{fancyhdr}
% [most] charge toutes les bibliothèques tcolorbox (listings, minted, raster,
% documentation...) et gonfle inutilement la mémoire TeX sur un document long
% (~300 boîtes) : on ne charge que ce qui est réellement utilisé.
\usepackage[skins,breakable]{tcolorbox}
\usepackage{graphicx}
\usepackage{colortbl}
\usepackage{booktabs}
\usepackage{tabularx}
\usepackage{diagbox} % case d'en-tête diagonale des tableaux à double entrée
% Colonnes extensibles centrée / gauche / droite (C, L, R), souvent utilisées par les modèles
\newcolumntype{C}{>{\centering\arraybackslash}X}
\newcolumntype{L}{>{\raggedright\arraybackslash}X}
\newcolumntype{R}{>{\raggedleft\arraybackslash}X}
\usepackage{array}
\usepackage{multicol}
\usepackage{siunitx}
% parse-numbers=false : accepte aussi \SI{2,5 \times 10^{-3}}{...}
\sisetup{locale=FR,output-decimal-marker={,},group-minimum-digits=4,parse-numbers=false}
\DeclareSIUnit\ohm{\ensuremath{\symup{\Omega}}} % U+2126 absent de la police
\DeclareSIUnit\division{div} % réglages d'oscilloscope (ms/div, V/div)
\DeclareSIUnit\tr{tr} % tours (vitesse de rotation, ex. \SI{1500}{\tr\per\minute})
\DeclareSIUnit\ch{ch} % cheval-vapeur (puissance moteur, unité usuelle hors SI)
\DeclareSIUnit\franccfa{FCFA} % franc CFA (coûts, contexte camerounais)
\DeclareSIUnit\dioptre{\ensuremath{\delta}} % vergence d'une lentille (dioptrie)
\usepackage{qrcode}
% \degree : commande fréquemment utilisée par les modèles pour un angle ou une
% température en dehors de \SI{}{} (non fournie par les paquets chargés).
\providecommand{\degree}{^{\circ}}
% \qed (fin de démonstration), utilisé par les modèles sans amsthm
\providecommand{\qed}{\hfill$\square$}
% \barre : double barre des valeurs interdites dans un tableau de variations (tkz-tab)
\providecommand{\barre}{\ensuremath{\|}}
% environnement proof (amsthm non chargé) utilisé par les modèles
\ifdefined\proof\else\newenvironment{proof}[1][Démonstration]{\par\noindent\textit{#1.}\ }{\qed\par}\fi
\usepackage{lastpage}
\usepackage{hyperref}
\hypersetup{hidelinks}

% ============================================================
% Palette « Pop » OnBuch+ — mêmes hex que la classe P (plus_ui.dart)
% ============================================================
\definecolor{poporange}{HTML}{F59321}
\definecolor{popdark}{HTML}{E07A0C}
\definecolor{popcream}{HTML}{FFF6E8}
\definecolor{popink}{HTML}{1C1714}
\definecolor{poppurple}{HTML}{7A5AE0}
\definecolor{popgreen}{HTML}{2E9E6A}
\definecolor{poppink}{HTML}{E0457B}
\definecolor{popblue}{HTML}{2F7DE1}
\definecolor{popgold}{HTML}{C98A12}
\definecolor{popmuted}{HTML}{8A7F76}
\definecolor{popline}{HTML}{EADFD2}
\definecolor{popgray}{HTML}{8A7F76}
\colorlet{popgrey}{popgray}
% Alias tolérants (le modèle nomme parfois les couleurs différemment)
\colorlet{popwhite}{white}
\colorlet{popblack}{popink}
\colorlet{popred}{poppink}
\colorlet{popyellow}{popgold}
% Teintes claires pour fonds de tableaux / zones
\colorlet{popblueL}{popblue!10!white}
\colorlet{poporangeL}{poporange!14!white}
\colorlet{popgreenL}{popgreen!12!white}
\colorlet{poppurpleL}{poppurple!12!white}
\colorlet{poppinkL}{poppink!12!white}
\colorlet{popgoldL}{popgold!14!white}

\pagecolor{popcream}

% ============================================================
% Typographie
% ============================================================
\defaultfontfeatures{Ligatures=TeX}
\setmainfont{PlusJakartaSans-Medium.ttf}[Path=../../../fonts/,BoldFont=PlusJakartaSans-Bold.ttf]
% Maths en sans-serif (Fira Math) avec chiffres et lettres droites de Plus
% Jakarta, pour que les formules se fondent dans le texte.
\usepackage[math-style=ISO,bold-style=ISO]{unicode-math}
\setmathfont{FiraMath-Regular.otf}[Path=../../../fonts/]
\setmathfont{PlusJakartaSans-Medium.ttf}[Path=../../../fonts/,range={up/num,up/{Latin,latin},\mathup}]
\usepackage{icomma} % virgule décimale sans espace en mode math
% Caractères mathématiques Unicode absents des polices de texte (Plus Jakarta,
% Archivo Black) : rendus par leur équivalent mathématique, en texte comme en
% formule — sinon ils s'affichent en carré vide (ex. « Arithmétique dans ℤ »).
\usepackage{newunicodechar}
\newunicodechar{ℝ}{\ensuremath{\mathbb{R}}}
\newunicodechar{ℤ}{\ensuremath{\mathbb{Z}}}
\newunicodechar{ℂ}{\ensuremath{\mathbb{C}}}
\newunicodechar{ℕ}{\ensuremath{\mathbb{N}}}
\newunicodechar{ℚ}{\ensuremath{\mathbb{Q}}}
\newunicodechar{𝔻}{\ensuremath{\mathbb{D}}}
\newunicodechar{α}{\ensuremath{\alpha}}
\newunicodechar{β}{\ensuremath{\beta}}
\newunicodechar{γ}{\ensuremath{\gamma}}
\newunicodechar{δ}{\ensuremath{\delta}}
\newunicodechar{ε}{\ensuremath{\varepsilon}}
\newunicodechar{θ}{\ensuremath{\theta}}
\newunicodechar{λ}{\ensuremath{\lambda}}
\newunicodechar{μ}{\ensuremath{\mu}}
\newunicodechar{σ}{\ensuremath{\sigma}}
\newunicodechar{τ}{\ensuremath{\tau}}
\newunicodechar{φ}{\ensuremath{\varphi}}
\newunicodechar{ψ}{\ensuremath{\psi}}
\newunicodechar{ω}{\ensuremath{\omega}}
\newunicodechar{Σ}{\ensuremath{\Sigma}}
% majuscules produites par \MakeUppercase dans les titres (α-aminés -> Α)
\newunicodechar{Α}{\ensuremath{\alpha}}
\newunicodechar{Β}{\ensuremath{\beta}}
\newunicodechar{Γ}{\ensuremath{\Gamma}}
\newunicodechar{Δ}{\ensuremath{\Delta}}
\newunicodechar{Ε}{\ensuremath{\varepsilon}}
\newunicodechar{Θ}{\ensuremath{\Theta}}
\newunicodechar{Λ}{\ensuremath{\Lambda}}
\newunicodechar{Μ}{\ensuremath{\mu}}
\newunicodechar{Τ}{\ensuremath{\tau}}
\newunicodechar{Φ}{\ensuremath{\Phi}}
\newunicodechar{Ψ}{\ensuremath{\Psi}}
\newunicodechar{Ω}{\ensuremath{\Omega}}
\newunicodechar{↦}{\ensuremath{\mapsto}}
\newunicodechar{∈}{\ensuremath{\in}}
\newunicodechar{∉}{\ensuremath{\notin}}
\newunicodechar{∧}{\ensuremath{\wedge}}
\newunicodechar{⇔}{\ensuremath{\Leftrightarrow}}
\newunicodechar{⟺}{\ensuremath{\Longleftrightarrow}}
\newunicodechar{⇒}{\ensuremath{\Rightarrow}}
\newunicodechar{⟹}{\ensuremath{\Longrightarrow}}
\newunicodechar{∘}{\ensuremath{\circ}}
\newunicodechar{∪}{\ensuremath{\cup}}
\newunicodechar{∩}{\ensuremath{\cap}}
\newunicodechar{≡}{\ensuremath{\equiv}}
\newunicodechar{⊂}{\ensuremath{\subset}}
\newunicodechar{⊥}{\ensuremath{\perp}}
\newunicodechar{∃}{\ensuremath{\exists}}
\newunicodechar{∀}{\ensuremath{\forall}}
\newunicodechar{∛}{\ensuremath{\sqrt[3]{\,}}}
\newunicodechar{∜}{\ensuremath{\sqrt[4]{\,}}}
\newunicodechar{⃗}{\ensuremath{{}^{\to}}}
\newunicodechar{ⁿ}{\ifmmode^{n}\else\textsuperscript{\ensuremath{n}}\fi}
\newunicodechar{ˣ}{\ifmmode^{x}\else\textsuperscript{\ensuremath{x}}\fi}
\newunicodechar{ᵇ}{\ifmmode^{b}\else\textsuperscript{\ensuremath{b}}\fi}
\newunicodechar{ʳ}{\ifmmode^{r}\else\textsuperscript{\ensuremath{r}}\fi}
\newunicodechar{ᵘ}{\ifmmode^{u}\else\textsuperscript{\ensuremath{u}}\fi}
\newunicodechar{ʸ}{\ifmmode^{y}\else\textsuperscript{\ensuremath{y}}\fi}
\newunicodechar{ᵃ}{\ifmmode^{a}\else\textsuperscript{\ensuremath{a}}\fi}
\newunicodechar{ᵉ}{\ifmmode^{e}\else\textsuperscript{\ensuremath{e}}\fi}
\newunicodechar{ᵏ}{\ifmmode^{k}\else\textsuperscript{\ensuremath{k}}\fi}
\newunicodechar{ᵖ}{\ifmmode^{p}\else\textsuperscript{\ensuremath{p}}\fi}
\newunicodechar{ᴺ}{\ifmmode^{N}\else\textsuperscript{\ensuremath{N}}\fi}
\newunicodechar{ᶜ}{\ifmmode^{c}\else\textsuperscript{\ensuremath{c}}\fi}
\newunicodechar{ʰ}{\ifmmode^{h}\else\textsuperscript{\ensuremath{h}}\fi}
\newunicodechar{ᵠ}{\ifmmode^{q}\else\textsuperscript{\ensuremath{q}}\fi}
\newunicodechar{ₙ}{\ifmmode_{n}\else\textsubscript{\ensuremath{n}}\fi}
\newunicodechar{ₐ}{\ifmmode_{a}\else\textsubscript{\ensuremath{a}}\fi}
\newunicodechar{ᵢ}{\ifmmode_{i}\else\textsubscript{\ensuremath{i}}\fi}
\newunicodechar{ᵣ}{\ifmmode_{r}\else\textsubscript{\ensuremath{r}}\fi}
\newunicodechar{ₖ}{\ifmmode_{k}\else\textsubscript{\ensuremath{k}}\fi}
\newunicodechar{ᵦ}{\ifmmode_{\beta}\else\textsubscript{\ensuremath{\beta}}\fi}
\newunicodechar{ₓ}{\ifmmode_{x}\else\textsubscript{\ensuremath{x}}\fi}
\newfontfamily\popdisplay{ArchivoBlack-Regular.ttf}[Path=../../../fonts/]
\newfontfamily\poplogo{SpaceGrotesk-Bold.ttf}[Path=../../../fonts/]
\newfontfamily\popsemi{PlusJakartaSans-SemiBold.ttf}[Path=../../../fonts/]
\newfontfamily\popxbold{PlusJakartaSans-ExtraBold.ttf}[Path=../../../fonts/]
\newfontfamily\popxboldls{PlusJakartaSans-ExtraBold.ttf}[Path=../../../fonts/,LetterSpace=3.0]

\setlist[enumerate]{leftmargin=1.7em,itemsep=5pt,topsep=4pt}
\setlist[itemize]{leftmargin=1.7em,itemsep=4pt,topsep=4pt,label={\color{poporange}\textbullet}}
\setlength{\parindent}{0pt}
\setlength{\parskip}{5pt}
\renewcommand{\arraystretch}{1.25}

% pgfplots : style Pop par défaut
\pgfplotsset{
  every axis/.append style={axis line style={popink,line width=1pt},tick style={popink},
    grid style={popline,line width=0.5pt},label style={font=\small},tick label style={font=\footnotesize}},
  popcurve/.style={poporange,line width=1.8pt,no markers,smooth},
}
% Même style disponible dans un \draw TikZ simple (hors pgfplots)
\tikzset{popcurve/.style={poporange,line width=1.8pt}}
\tikzset{popgrid/.style={popline,very thin}}
\colorlet{popcurve}{poporange} % le nom du style est parfois utilisé comme couleur

% ============================================================
% Pill / squiggle
% ============================================================
\newcommand{\poppill}[3]{%
  \tikz[baseline=-0.55ex]\node[draw=popink,line width=1.2pt,rounded corners=2.6pt,%
    fill=#2,inner xsep=6.5pt,inner ysep=3pt]{%
    \popxboldls\fontsize{7}{7}\selectfont\color{#3}\MakeUppercase{#1}};%
}
\newcommand{\popsquiggle}[1]{%
  \tikz[baseline=-0.4ex]{
    \draw[#1,line width=2.6pt,line cap=round]
      plot [smooth] coordinates {(0,0.09) (0.34,0.01) (0.68,0.21) (1.02,0.01) (1.36,0.10)};
  }%
}
% Mot-clé surligné (marqueur orange sous le texte)
\newcommand{\cle}[1]{{\popxbold\color{popdark}#1}}

% ============================================================
% En-tête / pied
% ============================================================
\pagestyle{fancy}
\fancyhf{}
\renewcommand{\headrulewidth}{0pt}
\renewcommand{\footrulewidth}{0pt}
\lhead{\raisebox{-0.15ex}{\poplogo\fontsize{13}{13}\selectfont\color{popink}OnBuch\color{poporange}+}}
\rhead{\poppill{\DOCMATIERE\ \textbullet\ \DOCNIVEAU\ \textbullet\ Leçon \DOCLECON}{white}{popink}}
\lfoot{\popxbold\fontsize{7.6}{7.6}\selectfont\color{popmuted}\MakeUppercase{Cours signé OnBuch+}\ \textbullet\ {\popsemi\DOCTITRE}}
\rfoot{\tikz[baseline=-0.6ex]\node[rounded corners=8.5pt,fill=popink,minimum height=17pt,minimum width=17pt,inner xsep=5pt,inner sep=0]{\popxbold\fontsize{8}{8}\selectfont\color{popcream}\thepage\,/\,\pageref*{LastPage}};}

% ============================================================
% Titres
% ============================================================
\newcommand{\chaptitre}[2]{%
  \begin{tcolorbox}[enhanced,colback=poporange,colframe=popink,boxrule=2.6pt,arc=11pt,
      drop shadow=popink,left=15pt,right=15pt,top=12pt,bottom=11pt,before skip=3pt,after skip=18pt]
    \poppill{#2}{popink}{popcream}\par\vspace{9pt}
    {\raggedright\hyphenpenalty=10000\exhyphenpenalty=10000\popdisplay\fontsize{22}{24}\selectfont\color{popcream}\MakeUppercase{#1}\par}
  \end{tcolorbox}
}
% Section numérotée : pastille orange avec le numéro + titre Archivo
\newcounter{popsec}
\newcommand{\coursec}[1]{%
  \stepcounter{popsec}\setcounter{popsub}{0}%
  \par\vspace{16pt}\noindent
  \begin{minipage}{\linewidth}
  \tikz[baseline=-0.7ex]\node[draw=popink,line width=1.6pt,fill=poporange,rounded corners=4pt,minimum size=22pt,inner sep=0]{\popdisplay\fontsize{12}{12}\selectfont\color{popcream}\thepopsec};%
  \hspace{8pt}{\popdisplay\fontsize{15}{17}\selectfont\color{popink}\MakeUppercase{#1}}\par
  \vspace{4pt}\noindent{\color{poporange}\rule{36pt}{3.6pt}}
  \end{minipage}\par\nopagebreak\vspace{8pt}
}
\newcounter{popsub}
\newcommand{\courssub}[1]{%
  \stepcounter{popsub}%
  \par\vspace{9pt}\noindent{\popxbold\fontsize{11.5}{13}\selectfont\color{popink}{\color{poporange}\thepopsec.\thepopsub}\hspace{6pt}#1}\par\nopagebreak\vspace{4pt}
}

% ============================================================
% Boîtes pédagogiques — même recette : fond blanc, bordure épaisse colorée,
% ombre dure encre, badge plein. Titre optionnel [#1].
% ============================================================
\tcbset{popbase/.style={breakable,enhanced,colback=white,boxrule=2.3pt,arc=9pt,
  shadow={3pt}{-3pt}{0pt}{popink},left=14pt,right=14pt,top=11pt,bottom=11pt,before skip=11pt,after skip=15pt,
  fontupper=\popsemi\fontsize{10.6}{14.5}\selectfont\color{popink},
  fontlower=\popsemi\fontsize{10.6}{14.5}\selectfont\color{popink}}}
% Coche dessinée (le glyphe ✓ n'existe pas dans les polices du document)
\newcommand{\popcheck}[1]{\tikz[baseline=-0.5ex]\draw[#1,line width=1.6pt,line cap=round,line join=round](-0.13,0.0)--(-0.04,-0.09)--(0.13,0.1);}
\providecommand{\coche}{\popcheck{popgreen}}
\providecommand{\resultat}[1]{\par\smallskip\noindent\fcolorbox{popgreen}{popgreenL}{\parbox{\dimexpr\linewidth-2\fboxsep-2\fboxrule\relax}{\textbf{Résultat.} #1}}\par\smallskip}
\newcommand{\popboxhead}[3]{\poppill{#1}{#2}{white}\ifx\relax#3\relax\else\hspace{6pt}{\popxbold\fontsize{10.6}{12}\selectfont\color{popink}#3}\fi\par\vspace{7pt}}

\newtcolorbox{aretenir}[1][]{popbase,colframe=popgreen,before upper={\popboxhead{À retenir}{popgreen}{#1}}}
\newtcolorbox{exemplebox}[1][]{popbase,colframe=poporange,before upper={\popboxhead{Exemple}{poporange}{#1}}}
\newtcolorbox{definition}[1][]{popbase,colframe=popblue,before upper={\popboxhead{Définition}{popblue}{#1}}}
\newtcolorbox{propriete}[1][]{popbase,colframe=poppurple,before upper={\popboxhead{Propriété}{poppurple}{#1}}}
\newtcolorbox{methode}[1][]{popbase,colframe=popgold,before upper={\popboxhead{Méthode}{popgold}{#1}}}
\newtcolorbox{attention}[1][]{popbase,colframe=poppink,before upper={\popboxhead{Attention}{poppink}{#1}}}
\newtcolorbox{experience}[1][]{popbase,colframe=popblue,colback=popblueL,before upper={\popboxhead{Expérience}{popblue}{#1}}}
% « Le savais-tu ? » : carte violette pleine (contexte, culture, Cameroun)
\newtcolorbox{savaistu}[1][]{popbase,colback=poppurple,colframe=popink,
  fontupper=\popsemi\fontsize{10.4}{14.2}\selectfont\color{white},
  before upper={\poppill{Le savais-tu\,?}{white}{poppurple}\ifx\relax#1\relax\else\hspace{6pt}{\popxbold\color{white}#1}\fi\par\vspace{7pt}}}
% Exercice résolu : énoncé en haut, \tcblower puis la solution sur fond crème
\newcounter{popexo}\newcounter{popexr}
\newtcolorbox{exoresolu}[1][]{popbase,colframe=poporange,colbacklower=popcream,
  segmentation style={popink,line width=1.2pt,dashed},
  before upper={\stepcounter{popexr}\popboxhead{Exercice résolu \thepopexr}{poporange}{#1}},
  before lower={\poppill{Solution}{popink}{popcream}\par\vspace{6pt}}}
% Exercice (non résolu en place) — la correction est dans la partie Corrigés
\newtcolorbox{exercice}[1][]{popbase,colframe=popink,
  before upper={\stepcounter{popexo}\popboxhead{Exercice \thepopexo}{popink}{#1}}}
% Corrigé
\newtcolorbox{corrige}[1][]{popbase,colframe=popgreen,colback=popgreenL,
  before upper={\popboxhead{Corrigé}{popgreen}{#1}}}
% Objectifs / prérequis / bilan
\newtcolorbox{objectifs}{popbase,colframe=popink,colback=poporangeL,
  before upper={\popboxhead{Objectifs de la leçon}{popink}{}}}
\newtcolorbox{prerequis}{popbase,colframe=popink,colback=popblueL,
  before upper={\popboxhead{Prérequis}{popblue}{}}}
\newtcolorbox{bilan}{popbase,colframe=popink,colback=white,boxrule=2.8pt,
  before upper={\popboxhead{Fiche bilan}{poporange}{L'essentiel à connaître pour l'examen}}}
% Figure encadrée avec légende
\newtcolorbox{popfigure}[1][]{enhanced,breakable=false,colback=white,colframe=popline,boxrule=1.4pt,arc=8pt,
  left=10pt,right=10pt,top=10pt,bottom=8pt,before skip=10pt,after skip=12pt,halign=center,
  lower separated=false,halign lower=center,
  fontlower=\popsemi\fontsize{9}{11.5}\selectfont\color{popmuted},
  #1}
\newcommand{\legende}[1]{\par\vspace{6pt}{\popsemi\fontsize{9}{11.5}\selectfont\color{popmuted}#1}}

% ============================================================
% Signature OnBuch+
% ============================================================
\newcommand{\poplogoinline}[1]{{\poplogo\fontsize{#1}{#1}\selectfont\color{popink}OnBuch\color{poporange}+}}
\newcommand{\signatureonbuch}{%
  \begin{tcolorbox}[enhanced,colback=popink,colframe=popink,boxrule=2.2pt,arc=10pt,
      drop shadow=poporange,left=14pt,right=14pt,top=10pt,bottom=10pt,before skip=0pt,after skip=0pt]
    \tikz[baseline=-0.7ex]\node[circle,fill=popgreen,draw=popcream,line width=1.3pt,minimum size=22pt,inner sep=0]{\popcheck{white}};%
    \hspace{9pt}\begin{minipage}[c]{0.62\linewidth}
      {\popxbold\fontsize{12}{13}\selectfont\color{popcream}Signé\ \textbullet\ {\poplogo OnBuch\color{poporange}+}}\par\vspace{2pt}
      {\raggedright\popsemi\fontsize{8}{10.5}\selectfont\color{popline}\MakeUppercase{Équipe pédagogique OnBuch+}\\\MakeUppercase{Conforme au programme officiel MINESEC}\par}
    \end{minipage}\hfill
    {\poplogo\fontsize{22}{22}\selectfont\color{popcream}OnBuch\color{poporange}+}
  \end{tcolorbox}%
}

% ============================================================
% Page de garde
% #1 titre · #2 matière · #3 niveau · #4 module · #5 n° leçon · #6 durée ·
% #7 description / objectifs (texte ou itemize)
% ============================================================
\newcommand{\pagegarde}[7]{%
  \thispagestyle{empty}
  \newgeometry{a4paper,margin=1.7cm,top=1.6cm,bottom=1.4cm}%
  \begin{tikzpicture}[remember picture,overlay]
    \fill[poporange!18] ([xshift=-3.2cm,yshift=-3.6cm]current page.north east) circle (4.6cm);
    \fill[poppurple!14] ([xshift=2.4cm,yshift=7.6cm]current page.south west) circle (3.8cm);
  \end{tikzpicture}%
  {\poplogo\fontsize{30}{30}\selectfont\color{popink}OnBuch\color{poporange}+}\hfill
  \raisebox{0.6ex}{\poppill{Cours signé \textbullet\ Premium}{popink}{popcream}}\par
  \vspace{1pt}\popsquiggle{poppurple}\par
  \vspace{8pt}
  \begin{tcolorbox}[enhanced,colback=poporange,colframe=popink,boxrule=2.8pt,arc=13pt,
      drop shadow=popink,left=18pt,right=18pt,top=12pt,bottom=13pt,after skip=10pt]
    \poppill{#2 \textbullet\ #3}{popink}{popcream}\hspace{5pt}\poppill{#4}{white}{popink}\par\vspace{8pt}
    {\popdisplay\fontsize{14}{14}\selectfont\color{popink}LEÇON #5}\par\vspace{4pt}
    {\raggedright\hyphenpenalty=10000\exhyphenpenalty=10000\popdisplay\fontsize{25}{27}\selectfont\color{popcream}\MakeUppercase{#1}\par}
    \vspace{8pt}
    {\popxbold\fontsize{9.5}{9.5}\selectfont\color{popink}\MakeUppercase{Durée conseillée : #6}}
  \end{tcolorbox}
  \begin{tcolorbox}[enhanced,colback=white,colframe=popink,boxrule=2.2pt,arc=10pt,
      drop shadow=popink,left=16pt,right=16pt,top=10pt,bottom=10pt,after skip=8pt]
    \poppill{Dans cette leçon}{poppurple}{white}\par\vspace{5pt}
    \popsemi\fontsize{9.3}{12.6}\selectfont\color{popink}#7
  \end{tcolorbox}
  \noindent\begin{minipage}[t]{0.487\linewidth}
    \begin{tcolorbox}[enhanced,colback=popgreen,colframe=popink,boxrule=2.2pt,arc=10pt,
        drop shadow=popink,left=11pt,right=11pt,top=10pt,bottom=10pt,height=3.1cm,valign=center]
      \tcbox[colback=white,colframe=popink,boxrule=1.3pt,arc=5pt,boxsep=0pt,nobeforeafter,
        left=3pt,right=3pt,top=3pt,bottom=3pt,valign=center]{\qrcode[height=1.9cm]{https://play.google.com/store/apps/details?id=com.luvvix.onbuch&hl=fr}}%
      \hspace{8pt}\begin{minipage}[c]{0.5\linewidth}
        {\popdisplay\fontsize{11}{12.5}\selectfont\color{white}\MakeUppercase{Télécharge OnBuch}}\par\vspace{4pt}
        {\raggedright\popsemi\fontsize{8.4}{11}\selectfont\color{white}Gratuit \textbullet\ Google Play\\Tuteur IA, annales, concours, résultats\par}
      \end{minipage}
    \end{tcolorbox}
  \end{minipage}\hfill
  \begin{minipage}[t]{0.487\linewidth}
    \begin{tcolorbox}[enhanced,colback=poppurple,colframe=popink,boxrule=2.2pt,arc=10pt,
        drop shadow=popink,left=11pt,right=11pt,top=10pt,bottom=10pt,height=3.1cm,valign=center]
      \tikz[baseline=-0.7ex]\node[circle,fill=white,draw=popink,line width=1.3pt,minimum size=1.6cm,inner sep=0]{\popdisplay\fontsize{24}{24}\selectfont\color{poppurple}+};%
      \hspace{8pt}\begin{minipage}[c]{0.6\linewidth}
        {\popdisplay\fontsize{11}{12.5}\selectfont\color{white}\MakeUppercase{Passe à OnBuch+}}\par\vspace{4pt}
        {\raggedright\popsemi\fontsize{8.4}{11}\selectfont\color{white}Tous les cours signés, Tuteur IA illimité, fiches PDF — onglet OnBuch+ de l'app\par}
      \end{minipage}
    \end{tcolorbox}
  \end{minipage}
  \vspace{8pt}
  \popcontenu
  \vfill
  \signatureonbuch
  \restoregeometry
  \newpage
}

% Rangée « Ce cours contient » (page de garde)
\newcommand{\poptile}[3]{% #1 couleur #2 gros texte #3 légende
  \begin{tcolorbox}[enhanced,colback=#1,colframe=popink,boxrule=2pt,arc=9pt,shadow={2.5pt}{-2.5pt}{0pt}{popink},
      left=6pt,right=6pt,top=9pt,bottom=9pt,halign=center,valign=center,height=1.75cm,width=\linewidth,nobeforeafter]
    {\popdisplay\fontsize{12.5}{13}\selectfont\color{white}\MakeUppercase{#2}}\par\vspace{3pt}
    {\centering\popsemi\fontsize{7.8}{9.5}\selectfont\color{white}#3\par}
  \end{tcolorbox}}
\newcommand{\popcontenu}{%
  \par\noindent{\popxboldls\fontsize{7.5}{7.5}\selectfont\color{popmuted}CE COURS SIGNÉ CONTIENT}\par\vspace{7pt}
  \noindent\begin{minipage}[t]{0.235\linewidth}\poptile{popblue}{Cours}{complet, illustré, pas à pas}\end{minipage}\hfill
  \begin{minipage}[t]{0.235\linewidth}\poptile{poporange}{Méthodes}{et exercices résolus}\end{minipage}\hfill
  \begin{minipage}[t]{0.235\linewidth}\poptile{poppink}{Exercices}{type examen + corrigés}\end{minipage}\hfill
  \begin{minipage}[t]{0.235\linewidth}\poptile{popgreen}{Fiche bilan}{l'essentiel à retenir}\end{minipage}\par}

% Sommaire
\newtcolorbox{sommaire}{enhanced,colback=white,colframe=popink,boxrule=2.2pt,arc=10pt,
  drop shadow=popink,left=18pt,right=18pt,top=13pt,bottom=13pt,before skip=6pt,after skip=20pt,
  before upper={\poppill{Sommaire}{poppurple}{white}\par\vspace{9pt}\popsemi\fontsize{10.6}{14.5}\selectfont\color{popink}}}

% Signature de fin de cours — poussée tout en bas de la dernière page
\newcommand{\findecours}{%
  \par\vfill\nopagebreak
  \begin{center}\popsquiggle{poporange}\end{center}\vspace{4pt}
  \signatureonbuch
}
\AtBeginDocument{\def\llbracket{\mathopen{[\mkern-3.5mu[}}\def\rrbracket{\mathclose{]\mkern-3.5mu]}}} % intervalles d'entiers (glyphe ⟦ absent de la police)
% Glyphes absents de Fira Math : redéfinis à partir de glyphes disponibles
\AtBeginDocument{%
  \def\setminus{\mathbin{\backslash}}\let\smallsetminus\setminus
  \def\vdots{\mathord{\vbox{\baselineskip3.2pt\lineskiplimit0pt\kern4pt\hbox{.}\hbox{.}\hbox{.}}}}%
  \def\ddots{\mathinner{\mkern1mu\raise6.5pt\hbox{.}\mkern2mu\raise3.6pt\hbox{.}\mkern2mu\raise0.7pt\hbox{.}\mkern1mu}}%
  \def\triangle{\mathord{\scalebox{0.9}{$\symup{\Delta}$}}}%
  \def\nabla{\mathord{\raisebox{\depth}{\scalebox{1}[-1]{$\symup{\Delta}$}}}}%
  \def\checkmark{\popcheck{popdark}}%
  \def\star{\mathbin{\ast}}\def\bigstar{\mathord{\ast}}%
  \def\diamond{\mathbin{\scalebox{0.6}{$\Diamond$}}}%
  \def\bigcup{\mathop{\mathchoice{\raisebox{-0.3em}{\scalebox{1.7}{$\cup$}}}{\scalebox{1.25}{$\cup$}}{\cup}{\cup}}\displaylimits}%
  \def\bigcap{\mathop{\mathchoice{\raisebox{-0.3em}{\scalebox{1.7}{$\cap$}}}{\scalebox{1.25}{$\cap$}}{\cap}{\cap}}\displaylimits}%
  \def\lesseqqgtr{\mathrel{\lessgtr}}\def\bigoplus{\mathop{\oplus}\displaylimits}%
}
\providecommand{\ssi}{si et seulement si} % abréviation fréquente
\AtBeginDocument{\providecommand{\wideparen}{\overparen}} % arc de cercle

%% FIGFIX-BEGIN v4 — correctifs globaux des figures (docs/onbuch-generation/tools/figfix)
\usepackage{adjustbox}\usepackage{etoolbox}
% toute figure TikZ/pgfplots plus large que la ligne est réduite à la largeur disponible
\BeforeBeginEnvironment{tikzpicture}{\begin{adjustbox}{max width=\linewidth}}
\AfterEndEnvironment{tikzpicture}{\end{adjustbox}}
% légendes pgfplots lisibles par-dessus les courbes
\pgfplotsset{every axis/.append style={legend style={fill=white,fill opacity=0.92,text opacity=1,draw=black!15,inner sep=3pt}}}
% étiquettes d'arêtes (syntaxe edge["..."]) avec fond blanc
\tikzset{every edge quotes/.append style={fill=white,inner sep=1.2pt}}
%% FIGFIX-END
\begin{document}

\pagegarde{Sensibilisation sur la nécessité de la mitose pour le maintien de l'identité biologique des organismes}{SVT}{1ère D}{Module 1 : Le Monde Vivant}{ND01}{16 h}{Cette leçon étudie le cycle cellulaire et la mitose, mécanisme fondamental qui assure la multiplication des cellules tout en conservant l'identité génétique des organismes. Elle montre comment la réplication de l'ADN et le comportement ordonné des chromosomes garantissent la stabilité du caryotype, et analyse les conséquences d'un dérèglement de ce processus.
\vspace{4pt}
\begin{multicols}{2}\small
\begin{itemize}
\item À la fin de cette leçon, je sais décrire les étapes du cycle cellulaire (interphase et phase M)
\item À la fin de cette leçon, je sais identifier les quatre phases de la mitose sur des photographies ou des schémas
\item À la fin de cette leçon, je sais décrire le comportement des chromosomes au cours de chaque phase de la mitose
\item À la fin de cette leçon, je sais expliquer le rôle de la réplication de l'ADN dans la conservation de l'information génétique
\item À la fin de cette leçon, je sais établir le lien entre réplication de l'ADN et conservation du caryotype d'une génération cellulaire à l'autre
\item À la fin de cette leçon, je sais interpréter des données expérimentales sur la prolifération cellulaire (courbes, comptages, marquages)
\end{itemize}\end{multicols}}

\begin{sommaire}
\begin{multicols}{2}
\begin{enumerate}
\item Le cycle cellulaire
\item La réplication de l'ADN et la conservation de l'information génétique
\item Les phases de la mitose
\item Comportement des chromosomes et conservation du caryotype
\item Dysfonctionnements de la mitose et cancérisation
\item Interprétation de données expérimentales sur la prolifération cellulaire
\item Activité d'intégration
\item Méthodes et exercices résolus
\item Exercices
\item Corrigés des exercices
\item Fiche bilan
\end{enumerate}
\end{multicols}
\end{sommaire}

\begin{objectifs}
\begin{itemize}
\item À la fin de cette leçon, je sais décrire les étapes du cycle cellulaire (interphase et phase M)
\item À la fin de cette leçon, je sais identifier les quatre phases de la mitose sur des photographies ou des schémas
\item À la fin de cette leçon, je sais décrire le comportement des chromosomes au cours de chaque phase de la mitose
\item À la fin de cette leçon, je sais expliquer le rôle de la réplication de l'ADN dans la conservation de l'information génétique
\item À la fin de cette leçon, je sais établir le lien entre réplication de l'ADN et conservation du caryotype d'une génération cellulaire à l'autre
\item À la fin de cette leçon, je sais interpréter des données expérimentales sur la prolifération cellulaire (courbes, comptages, marquages)
\item À la fin de cette leçon, je sais expliquer comment un dysfonctionnement de la mitose peut conduire à la cancérisation
\item À la fin de cette leçon, je sais justifier l'importance de la mitose dans la croissance, la réparation et la reproduction asexuée
\end{itemize}
\end{objectifs}

\begin{prerequis}
\begin{itemize}
\item La cellule : organisation générale, noyau et chromosomes (notions de Seconde)
\item L'ADN, support de l'information génétique (notion de gène, chromosome)
\item Le caryotype : nombre et morphologie des chromosomes d'une espèce (2n = 46 chez l'Homme)
\item Notions de microscopie et d'observation de coupes histologiques
\item Notion de division cellulaire vue au cycle d'observation (collège)
\end{itemize}
\end{prerequis}

\newpage

\coursec{Le cycle cellulaire}

\courssub{Situation d'accroche : une plaie qui se referme, un manioc qui repousse}

Aminatou prépare le ndolé pour toute la famille. En éminçant les feuilles, le couteau glisse et lui entaille le doigt. Quelques jours plus tard, la plaie est refermée : la peau s'est reconstituée, identique à l'ancienne, avec les mêmes caractéristiques. À des centaines de kilomètres de là, dans un champ de Bafoussam, un planteur enfouit un simple bout de tige de manioc. Quelques mois plus tard, un plant vigoureux s'est développé, génétiquement identique au plant mère : mêmes tubercules, même rendement, même goût.

Ces deux phénomènes, en apparence très différents, reposent sur le même mécanisme fondamental : la \cle{division cellulaire}. Pour réparer le doigt d'Aminatou, des cellules de la peau se sont divisées en grand nombre ; pour produire le nouveau plant de manioc, les cellules du fragment de tige se sont multipliées. Or, dans les deux cas, les cellules filles possèdent \emph{exactement} la même information génétique que les cellules dont elles sont issues.

\begin{center}
\begin{tcolorbox}[colback=poporangeL,colframe=poporange,title=\textbf{Problématique de la leçon}]
Comment les cellules qui se divisent parviennent-elles à transmettre exactement la même information génétique à leurs descendantes ? Que se passe-t-il quand ce mécanisme déraille, comme dans le cancer ? C'est toute la question de la \cle{mitose}, que nous allons étudier en commençant par le cadre qui l'organise : le \cle{cycle cellulaire}.
\end{tcolorbox}
\end{center}

\courssub{Observation et mise en évidence du cycle cellulaire}

\textbf{Observation.} Au microscope optique, si l'on observe une racine d'oignon ou de maïs (zone située juste derrière la coiffe, appelée \cle{méristème}), on constate un fait surprenant : la grande majorité des cellules présente un noyau « normal », bien coloré, sans structure particulière visible ; seule une petite minorité montre des chromosomes condensés en train de se répartir entre deux pôles. Pourtant, toutes ces cellules finissent par se diviser.

\textbf{Question.} Si la division est l'événement spectaculaire que l'on cherche, pourquoi la plupart des cellules semblent-elles « ne rien faire » ?

\textbf{Hypothèse.} La division proprement dite est un événement bref, précédé d'une longue période de préparation invisible au microscope optique.

\textbf{Données expérimentales.} Des expériences de marquage radioactif (thymidine tritiée, qui ne s'incorpore que dans l'ADN en cours de fabrication) ont montré qu'une cellule en croissance passe par des périodes bien distinctes : une période où elle fabrique de l'ADN (le marquage apparaît), encadrée par deux périodes où elle n'en fabrique pas, avant de se diviser. La mesure des durées confirme que la division ne dure qu'environ une heure, alors que la période de préparation dure souvent plus de dix heures.

\textbf{Conclusion.} La vie d'une cellule entre deux divisions n'est pas un « repos » : c'est un programme ordonné d'événements. Ce programme porte un nom : le \cle{cycle cellulaire}.

\begin{definition}[Le cycle cellulaire]
Le \cle{cycle cellulaire} est l'ensemble ordonné des événements qui se déroulent dans une cellule entre deux divisions successives. Il comprend deux grandes périodes :
\begin{itemize}
\item l'\cle{interphase}, longue période pendant laquelle la cellule croît, réalise ses activités métaboliques et \textbf{duplique son ADN} ;
\item la \cle{phase M} (ou \cle{mitose}), courte période pendant laquelle la cellule répartit équitablement son matériel génétique entre deux noyaux, puis se divise en deux cellules filles.
\end{itemize}
\end{definition}

\begin{attention}[L'interphase n'est pas une phase de repos]
Erreur très fréquente au Probatoire : écrire que « pendant l'interphase, la cellule est au repos ». C'est \textbf{faux}. L'interphase est une période d'\textbf{intense activité} : la cellule synthétise des protéines, produit de l'énergie, grandit et surtout \textbf{réplique la totalité de son ADN}. On l'appelait autrefois « phase de repos » uniquement parce qu'aucun événement n'y est visible au microscope optique. Retiens : repos \emph{apparent}, activité \emph{réelle}.
\end{attention}

\courssub{Les sous-phases de l'interphase : G1, S, G2}

L'interphase n'est pas homogène : elle se décompose en trois sous-phases successives, que l'on peut mémoriser comme les trois étapes de la préparation d'un grand voyage.

\textbf{La phase G1 (\emph{Gap 1}, premier « intervalle »).} C'est la phase de \cle{croissance cellulaire}. La cellule augmente de volume, fabrique activement des protéines (enzymes, protéines structurales), multiplie ses organites (mitochondries, ribosomes) et exerce sa fonction spécialisée (une cellule musculaire se contracte, une cellule du pancréas sécrète de l'insuline...). C'est la phase la plus variable en durée : elle peut durer quelques heures comme... toute une vie (voir « Le savais-tu »).

\textbf{La phase S (\emph{Synthèse}).} C'est l'étape décisive : la cellule \cle{réplique} (duplique) la totalité de son ADN. Chaque molécule d'ADN est copiée fidèlement, de sorte que chaque chromosome, initialement formé d'une seule molécule d'ADN (chromosome à une chromatide), devient un chromosome à \textbf{deux chromatides sœurs} identiques, réunies par le \cle{centromère}. La quantité d'ADN par noyau passe ainsi de $Q$ à $2Q$. C'est la garantie que chaque cellule fille recevra un exemplaire complet de l'information génétique.

\textbf{La phase G2 (\emph{Gap 2}, deuxième « intervalle »).} La cellule achève sa croissance et \textbf{prépare activement la division} : elle synthétise les protéines nécessaires à la mitose (notamment celles du fuseau de division), vérifie que la réplication de l'ADN s'est déroulée sans erreur et accumule les réserves d'énergie indispensables.

\begin{methode}[Mémoriser l'ordre des phases]
\begin{enumerate}
\item \textbf{G1} : la cellule \emph{grandit} et travaille (croissance).
\item \textbf{S} : la cellule \emph{synthétise} son ADN (réplication, $Q \to 2Q$).
\item \textbf{G2} : la cellule \emph{se prépare} à diviser (vérifications, synthèses).
\item \textbf{M} : la cellule \emph{mitose} : elle se divise ($2Q \to Q$ dans chaque cellule fille).
\end{enumerate}
Astuce : G = \emph{Gap} (intervalle), S = \emph{Synthèse}, M = \emph{Mitose}. L'ordre alphabétique des lettres suit l'ordre chronologique : G1, S, G2, M.
\end{methode}

\courssub{Représentation schématique du cycle cellulaire}

Le cycle cellulaire se représente classiquement par un cercle, car il se répète indéfiniment dans les cellules qui se divisent : chaque division redonne des cellules filles qui entrent à leur tour en G1.

\begin{popfigure}
\begin{center}
\begin{tikzpicture}[scale=1.0]
% Fond des secteurs
\fill[popgreenL] (0,0) -- (90:2.6) arc[start angle=90, end angle=54, radius=2.6] -- cycle;
\fill[poporangeL] (0,0) -- (54:2.6) arc[start angle=54, end angle=-126, radius=2.6] -- cycle;
\fill[popblueL] (0,0) -- (-126:2.6) arc[start angle=-126, end angle=-162, radius=2.6] -- cycle;
\fill[poppinkL] (0,0) -- (-162:2.6) arc[start angle=-162, end angle=-270, radius=2.6] -- cycle;
% Contours
\draw[popdark, thick] (0,0) circle (2.6);
\draw[popdark] (90:2.6) -- (0,0) -- (54:2.6);
\draw[popdark] (-126:2.6) -- (0,0) -- (-162:2.6);
\draw[popdark] (-270:2.6) -- (0,0);
% Flèches de sens
\draw[-{Stealth[length=3mm]}, popdark, thick] (90:2.95) arc[start angle=90, end angle=60, radius=2.95];
\draw[-{Stealth[length=3mm]}, popdark, thick] (54:2.95) arc[start angle=54, end angle=24, radius=2.95];
\draw[-{Stealth[length=3mm]}, popdark, thick] (-126:2.95) arc[start angle=-126, end angle=-156, radius=2.95];
\draw[-{Stealth[length=3mm]}, popdark, thick] (-162:2.95) arc[start angle=-162, end angle=-192, radius=2.95];
% Labels des phases
\node[font=\bfseries\large, popdark] at (72:1.7) {M};
\node[font=\small, popdark, align=center] at (72:1.15) {mitose\\ ($\approx$ 1 h)};
\node[font=\bfseries\large, popdark] at (-36:1.7) {G1};
\node[font=\small, popdark, align=center] at (-36:1.05) {croissance\\ ($\approx$ 6--8 h)};
\node[font=\bfseries\large, popdark] at (-144:1.7) {S};
\node[font=\small, popdark, align=center] at (-144:1.05) {réplication\\ de l'ADN ($\approx$ 7 h)};
\node[font=\bfseries\large, popdark] at (-198:1.7) {G2};
\node[font=\small, popdark, align=center] at (-198:1.05) {préparation\\ ($\approx$ 3--4 h)};
% Accolades textuelles
\node[font=\small\itshape, popdark, align=center] at (-90:3.6) {Interphase (longue) : G1 + S + G2};
\node[font=\small\itshape, popdark, align=center] at (90:3.5) {Phase M (courte)};
\end{tikzpicture}
\end{center}
\legende{Schéma circulaire du cycle cellulaire. Les secteurs sont proportionnés aux durées : l'interphase (G1, S, G2) occupe environ 90 \% du cycle, la phase M (mitose) environ 10 \%. Durées indicatives pour une cellule de mammifère en culture (cycle total $\approx$ 18 à 24 h).}
\end{popfigure}

\courssub{Évolution de la quantité d'ADN par noyau au cours du cycle}

Notons $Q$ la quantité d'ADN contenue dans le noyau d'une cellule en début de cycle (cellule en G1). Suivons cette quantité au cours du temps :

\begin{itemize}
\item pendant \textbf{G1}, la quantité d'ADN reste constante : $Q$ ;
\item pendant \textbf{S}, la réplication fait passer progressivement la quantité d'ADN de $Q$ à $2Q$ ;
\item pendant \textbf{G2} et toute la \textbf{mitose} (jusqu'à la séparation des deux noyaux), la quantité d'ADN reste égale à $2Q$ ;
\item à l'\textbf{issue de la mitose}, chaque cellule fille reçoit exactement la moitié : la quantité d'ADN par noyau revient à $Q$.
\end{itemize}

\begin{propriete}[Conservation de la quantité d'ADN au cours du cycle (admise)]
Au cours du cycle cellulaire, la quantité d'ADN par noyau \textbf{double pendant la phase S} (passage de $Q$ à $2Q$), puis \textbf{revient à sa valeur initiale $Q$ après la mitose}, par répartition équitable entre les deux cellules filles. Ainsi, de génération cellulaire en génération, la quantité d'ADN — et donc l'information génétique — reste \textbf{constante}. Cette propriété, admise ici, est illustrée par les données expérimentales de mesure de l'ADN (marquage radioactif, cytométrie).
\end{propriete}

\begin{popfigure}
\begin{center}
\begin{tikzpicture}
\begin{axis}[
width=0.9\linewidth, height=6.5cm,
xlabel={Temps (heures)},
ylabel={Quantité d'ADN par noyau},
xmin=0, xmax=24, ymin=0, ymax=2.6,
xtick={0,8,15,19,24},
ytick={1,2},
yticklabels={$Q$, $2Q$},
axis lines=left,
every axis x label/.style={at={(ticklabel* cs:1.02)}, anchor=west},
every axis y label/.style={at={(ticklabel* cs:1.02)}, anchor=south},
]
% Courbe quantité d'ADN
\addplot[popcurve, domain=0:8, samples=50] {1};
\addplot[popcurve, domain=8:15, samples=100] {1 + (x-8)/7};
\addplot[popcurve, domain=15:23, samples=50] {2};
\addplot[popcurve, domain=23:24, samples=20] {2 - (x-23)};
% Délimitations des phases
\draw[dashed, popmuted] (axis cs:8,0) -- (axis cs:8,2.5);
\draw[dashed, popmuted] (axis cs:15,0) -- (axis cs:15,2.5);
\draw[dashed, popmuted] (axis cs:19,0) -- (axis cs:19,2.5);
\draw[dashed, popmuted] (axis cs:23,0) -- (axis cs:23,2.5);
% Annotations des phases
\node[font=\small\bfseries, popgreen!60!popdark] at (axis cs:4,2.35) {G1};
\node[font=\small\bfseries, popblue!70!popdark] at (axis cs:11.5,2.35) {S};
\node[font=\small\bfseries, poporange!80!popdark] at (axis cs:17,2.35) {G2};
\node[font=\small\bfseries, poppink!80!popdark] at (axis cs:21,2.35) {M};
\node[font=\small, popdark, align=center] at (axis cs:11.5,1.35) {réplication\\ de l'ADN};
\node[font=\small, popdark, align=center] at (axis cs:22,0.55) {division :\\ $2Q \to Q$};
\end{axis}
\end{tikzpicture}
\end{center}
\legende{Évolution de la quantité d'ADN par noyau au cours du cycle cellulaire. Palier à $Q$ pendant G1, montée progressive de $Q$ à $2Q$ pendant la phase S (réplication), palier à $2Q$ pendant G2 et la mitose, puis chute brutale à $Q$ lors de la séparation des deux cellules filles.}
\end{popfigure}

\begin{methode}[Analyser un graphique de quantité d'ADN en fonction du temps]
Face à une courbe de quantité d'ADN par noyau au cours du temps, procède toujours en trois étapes :
\begin{enumerate}
\item \textbf{Repère le premier palier} (quantité constante $Q$) : c'est la phase \textbf{G1} — la cellule n'a pas encore répliqué son ADN.
\item \textbf{Repère la montée progressive} (de $Q$ à $2Q$) : c'est la phase \textbf{S} — la réplication est en cours. Une montée \emph{progressive} (et non brutale) traduit le fait que la copie de l'ADN prend plusieurs heures.
\item \textbf{Repère le second palier} (quantité constante $2Q$) : c'est \textbf{G2 puis la mitose}. La \textbf{chute brutale} de $2Q$ à $Q$ marque le partage de l'ADN entre les deux cellules filles, c'est-à-dire la fin de la division.
\end{enumerate}
Piège classique : ne confonds pas la montée (S) avec la mitose ! La mitose correspond à la \emph{chute} de la courbe, pas à sa montée.
\end{methode}

\begin{exoresolu}[Exploiter un graphique de quantité d'ADN]
\textbf{Énoncé.} On mesure la quantité d'ADN par noyau dans des cellules de méristème de racine de maïs en culture. On obtient les valeurs suivantes : de 0 à 7 h, la quantité est stable à $Q$ ; de 7 h à 14 h, elle augmente régulièrement jusqu'à $2Q$ ; de 14 h à 18 h, elle reste à $2Q$ ; à 19 h, elle retombe brutalement à $Q$.
\begin{enumerate}
\item Attribuer à chaque intervalle la phase du cycle cellulaire correspondante, en justifiant.
\item Calculer la durée approximative du cycle cellulaire de ces cellules.
\end{enumerate}
\tcblower
\textbf{Solution détaillée.}
\begin{enumerate}
\item \textbf{De 0 à 7 h : phase G1.} La quantité d'ADN est constante et égale à $Q$ : la cellule croît mais n'a pas encore répliqué son ADN. \textbf{De 7 h à 14 h : phase S.} L'augmentation régulière de $Q$ à $2Q$ traduit la réplication progressive de l'ADN. \textbf{De 14 h à 18 h : phase G2.} La quantité reste à $2Q$ : la réplication est achevée, la cellule prépare la division. \textbf{De 18 h à 19 h : phase M (mitose).} La chute brutale de $2Q$ à $Q$ correspond au partage équitable de l'ADN entre les deux cellules filles.
\item Le cycle complet s'étend du début de G1 (0 h) à la fin de la mitose (19 h). La durée du cycle est donc d'environ \textbf{19 heures}, dont 18 h d'interphase et 1 h de mitose : l'interphase représente environ 95 \% du cycle, ce qui confirme qu'elle est de loin la période la plus longue.
\end{enumerate}
\end{exoresolu}

\courssub{Durées du cycle selon les types cellulaires}

Toutes les cellules ne se divisent pas au même rythme. La durée du cycle cellulaire est très variable d'un type cellulaire à l'autre, et c'est presque toujours la durée de \textbf{G1} qui fait la différence.

\begin{center}
\begin{tabularx}{\linewidth}{|l|X|X|}
\hline
\rowcolor{popblueL}
\textbf{Type cellulaire} & \textbf{Durée indicative du cycle} & \textbf{Particularité} \\
\hline
Cellules du méristème racinaire (plantes) & 15 à 20 h & Division quasi continue : croissance de la racine \\
\hline
Cellules de la peau (épiderme) & 20 à 24 h & Renouvellement permanent de la couche superficielle \\
\hline
Cellules de la moelle osseuse & 20 à 30 h & Production incessante des globules sanguins \\
\hline
Cellules du foie (hépatocytes) & plusieurs mois à 1 an & Division rare, sauf en cas de lésion \\
\hline
Neurones & ne se divisent plus & Bloqués définitivement en G1 \\
\hline
\end{tabularx}
\end{center}

Les tissus à \cle{renouvellement rapide} (peau, moelle osseuse, muqueuse intestinale, méristèmes des plantes) possèdent des cellules au cycle court : ce sont eux qui permettent la cicatrisation du doigt d'Aminatou et la repousse du manioc. C'est d'ailleurs pour cette raison que les traitements anticancéreux (chimiothérapie), qui ciblent les cellules en division, provoquent des effets secondaires sur ces tissus : chute des cheveux, anémie, troubles digestifs.

\begin{savaistu}[Pourquoi une lésion de la moelle épinière est-elle si grave ?]
Certaines cellules humaines, comme les \cle{neurones}, sortent définitivement du cycle cellulaire : on dit qu'elles sont \textbf{bloquées en G1} (état parfois noté G0). Elles ne répliquent plus leur ADN et ne se divisent plus jamais. C'est pourquoi une blessure qui détruit des neurones — par exemple une section de la moelle épinière lors d'un accident de moto-taxi, fréquent sur les routes camerounaises — entraîne des séquelles souvent irréversibles : les neurones perdus ne sont pas remplacés, contrairement aux cellules de la peau. La recherche sur les \emph{cellules souches} tente précisément de contourner cette limite.
\end{savaistu}

\courssub{Les points de contrôle du cycle cellulaire}

Une cellule ne se divise pas n'importe quand ni n'importe comment. Le cycle cellulaire est surveillé par des mécanismes de contrôle internes, appelés \cle{points de contrôle} (\emph{checkpoints}), qui fonctionnent comme des « barrières de péage » : la cellule ne passe à l'étape suivante que si les conditions sont réunies.

\begin{definition}[Point de contrôle du cycle cellulaire]
Un \cle{point de contrôle} est une étape de vérification du cycle cellulaire : la cellule n'est autorisée à poursuivre le cycle que si les conditions requises sont remplies (taille suffisante, nutriments disponibles, ADN intact et correctement répliqué). Sinon, le cycle est bloqué, voire la cellule est éliminée.
\end{definition}

On retiendra qualitativement trois grands points de contrôle :

\begin{itemize}
\item \textbf{Point de contrôle G1/S} : la cellule vérifie sa taille, ses réserves et l'\textbf{intégrité de son ADN} avant de s'engager dans la réplication. C'est le contrôle le plus décisif : une fois la phase S commencée, la cellule est engagée jusqu'à la division.
\item \textbf{Point de contrôle G2/M} : la cellule vérifie que la réplication de l'ADN est \textbf{terminée et sans erreur} avant d'entrer en mitose.
\item \textbf{Point de contrôle de la mitose} : la cellule vérifie que tous les chromosomes sont correctement attachés avant de les séparer entre les deux futures cellules filles.
\end{itemize}

\begin{attention}[Le lien avec le cancer]
Ces points de contrôle sont des verrous de sécurité. Si les mécanismes de contrôle sont défaillants (par exemple à la suite de mutations), une cellule porteuse d'anomalies peut se diviser sans frein : c'est l'une des étapes de la \cle{cancérisation}, que nous étudierons dans la section 5. Retiens dès maintenant l'idée directrice : \textbf{le cancer est une maladie du cycle cellulaire}, dans laquelle les « barrières de péage » ne fonctionnent plus.
\end{attention}

\begin{aretenir}[L'essentiel de la section]
\begin{itemize}
\item Le \cle{cycle cellulaire} est la succession ordonnée d'une \textbf{interphase} et d'une \textbf{mitose} entre deux divisions cellulaires.
\item L'interphase (environ 90 \% du cycle) comprend trois sous-phases : \textbf{G1} (croissance), \textbf{S} (réplication de l'ADN), \textbf{G2} (préparation à la division). Ce n'est \textbf{pas} une phase de repos.
\item La quantité d'ADN par noyau passe de $Q$ à $2Q$ pendant la phase S, puis revient à $Q$ après la mitose : l'information génétique est ainsi \textbf{conservée} d'une génération cellulaire à l'autre.
\item La durée du cycle varie selon les cellules : courte dans les tissus à renouvellement rapide (peau, moelle osseuse, méristèmes), nulle chez les neurones bloqués en G1.
\item Des \textbf{points de contrôle} autorisent la poursuite du cycle uniquement si les conditions sont réunies ; leur défaillance peut conduire à la cancérisation.
\end{itemize}
\end{aretenir}

\begin{exercice}[Pour t'entraîner]
Des cellules de peau humaine en culture ont un cycle de 22 h, dont 1 h de mitose et 7 h de phase S. La phase G2 dure 4 h.
\begin{enumerate}
\item Calculer la durée de la phase G1.
\item Représenter ce cycle par un schéma circulaire en respectant les proportions.
\item Sur un graphique de la quantité d'ADN en fonction du temps, indiquer à quels instants la quantité d'ADN vaut exactement $2Q$.
\end{enumerate}
\end{exercice}

\begin{corrige}[Exercice : durées du cycle cellulaire]
\begin{enumerate}
\item Le cycle total vaut 22 h et se décompose en G1 + S + G2 + M. Donc :
\[ \text{G1} = 22 - (7 + 4 + 1) = 10 \text{ h}. \]
La phase G1 dure \textbf{10 heures}.
\item Le cercle (360°) représente 22 h, soit environ 16,4° par heure : G1 occupe $10 \times 16{,}4 \approx 164°$, S occupe $7 \times 16{,}4 \approx 115°$, G2 occupe $4 \times 16{,}4 \approx 65°$ et M occupe $1 \times 16{,}4 \approx 16°$. L'interphase (G1 + S + G2 = 21 h) occupe donc environ 95 \% du cercle.
\item La quantité d'ADN vaut $2Q$ de la \textbf{fin de la phase S} jusqu'à la \textbf{séparation des deux cellules filles}, c'est-à-dire pendant toute la phase G2 et la majeure partie de la mitose (palier à $2Q$ sur le graphique), soit entre l'instant $t = 17$ h et $t = 22$ h environ.
\end{enumerate}
\end{corrige}

\coursec{La réplication de l'ADN et la conservation de l'information génétique}

Dans la section précédente, nous avons découvert que la vie d'une cellule est rythmée par un \cle{cycle cellulaire} : une longue \cle{interphase} (phases G1, S et G2) pendant laquelle la cellule grandit et se prépare, suivie d'une courte \cle{mitose} au cours de laquelle elle se divise. Nous avons annoncé que la phase S était le moment d'un événement capital. Le voici : c'est pendant la phase S que la cellule \textbf{recopie intégralement son patrimoine génétique}. Sans cette recopie, la division serait un partage injuste : chaque cellule fille ne recevrait que la moitié de l'information. C'est ce mécanisme, la \cle{réplication de l'ADN}, que nous allons maintenant étudier en détail.

\courssub{Une situation concrète pour comprendre le problème}

Observe une simple éraflure au genou après une chute sur un terrain de quartier à Yaoundé. Quelques jours plus tard, la plaie s'est refermée : de nouvelles cellules de peau ont remplacé les cellules détruites. Pourtant, ces cellules nouvelles sont \emph{exactement} comme les anciennes : même couleur de peau, même fonction, même organisation. De même, quand un enfant grandit, les milliards de cellules qui s'ajoutent à son organisme portent toutes la même information génétique que la cellule œuf dont il est issu.

Posons le problème en termes simples. Toute l'information qui définit une cellule — son « mode d'emploi » — est écrite dans son \cle{ADN}. Quand une cellule se divise en deux, il faut que \textbf{chacune des deux cellules filles reçoive un exemplaire complet de ce mode d'emploi}. Or une cellule ne possède qu'un seul exemplaire de son ADN. Il est donc indispensable, \emph{avant} toute division, de fabriquer une copie. Comment la cellule s'y prend-elle, et comment être sûr que la copie est fidèle ?

\courssub{Rappel : la structure de l'ADN, support de l'information génétique}

\begin{definition}[L'ADN, molécule de l'hérédité]
L'\cle{ADN} (acide désoxyribonucléique) est une molécule formée de \textbf{deux brins} enroulés en double hélice. Chaque brin est une chaîne de \cle{nucléotides} ; chaque nucléotide comporte un sucre (désoxyribose), un groupement phosphate et une \cle{base azotée}. Il existe quatre bases : l'\textbf{adénine} (A), la \textbf{thymine} (T), la \textbf{guanine} (G) et la \textbf{cytosine} (C). Les deux brins sont unis par des liaisons hydrogène entre bases \textbf{complémentaires} : A s'associe toujours à T, et G toujours à C. L'\cle{information génétique} est portée par la \textbf{séquence} (l'ordre de succession) des nucléotides le long des brins.
\end{definition}

La conséquence de la complémentarité des bases est capitale : \textbf{connaître la séquence d'un brin, c'est connaître automatiquement celle de l'autre}. Si un brin porte la séquence A-T-G-C-C-A, le brin complémentaire porte nécessairement T-A-C-G-G-T. Les deux brins contiennent donc la \emph{même} information, écrite en quelque sorte « en miroir ». C'est cette redondance qui va permettre la copie fidèle.

\courssub{La réplication de l'ADN : une copie semi-conservative}

\begin{definition}[Réplication de l'ADN]
La \cle{réplication de l'ADN} est la duplication de la molécule d'ADN qui a lieu \textbf{avant la division cellulaire, pendant la phase S de l'interphase}. Elle se fait selon un mode dit \cle{semi-conservatif} : les deux brins de la molécule mère se séparent, et chacun sert de \textbf{matrice} (de modèle) à la synthèse d'un nouveau brin complémentaire. On obtient ainsi \textbf{deux molécules d'ADN filles identiques}, chacune constituée d'\textbf{un brin parental} (conservé) et d'\textbf{un brin néoformé} (nouvellement synthétisé).
\end{definition}

\begin{propriete}[Complémentarité des bases et fidélité de la copie (admise)]
Grâce à la complémentarité des bases (A avec T, G avec C), chaque brin parental sert de matrice à la synthèse d'un brin nouveau dont la séquence est \textbf{identique à celle du brin complémentaire d'origine}. En conséquence, les deux molécules filles sont identiques entre elles et identiques à la molécule mère : l'information génétique est recopiée \textbf{sans perte ni modification} (sauf erreurs rares, les mutations). Cette propriété est admise en Première D ; sa démonstration n'est pas exigible au Probatoire.
\end{propriete}

Déroulons le mécanisme pas à pas, comme un film ralenti :
\begin{enumerate}
\item Des \cle{enzymes} (hélicases) « dézippent » la double hélice : les liaisons hydrogène entre bases complémentaires sont rompues et les deux brins parentaux se séparent localement, formant un « œil de réplication ».
\item Sur chaque brin parental exposé, une enzyme appelée \cle{ADN polymérase} assemble des nucléotides libres présents dans le nucléoplasme, en respectant strictement la règle de complémentarité : face à un A, elle place un T ; face à un G, elle place un C.
\item La synthèse progresse le long de la molécule jusqu'à ce que la totalité de l'ADN soit dupliquée.
\item Résultat : deux molécules d'ADN filles, chacune hybride — un brin ancien, un brin neuf — et toutes deux rigoureusement identiques à la molécule mère.
\end{enumerate}

\begin{popfigure}
\begin{center}
\begin{tikzpicture}[scale=1, every node/.style={font=\small}]
% Molécule mère
\draw[line width=2.2pt, popblue] (0,0) -- (3.2,0);
\draw[line width=2.2pt, popblue] (0,0.5) -- (3.2,0.5);
\node[popblue, anchor=east] at (-0.1,0.25) {Molécule mère};
\node[popblue, anchor=west] at (3.3,0.5) {\scriptsize brin parental 1};
\node[popblue, anchor=west] at (3.3,0) {\scriptsize brin parental 2};
% Flèche séparation
\draw[-{Stealth[length=3mm]}, line width=1.2pt, popdark] (4.9,0.25) -- (6.1,0.25);
\node[popdark, align=center] at (5.5,0.85) {\scriptsize séparation des brins};
\node[popdark, align=center] at (5.5,-0.4) {\scriptsize + synthèse des brins nouveaux};
% Molécule fille 1
\draw[line width=2.2pt, popblue] (6.4,1.6) -- (9.6,1.6);
\draw[line width=2.2pt, poporange] (6.4,2.1) -- (9.6,2.1);
\node[anchor=east, align=right] at (6.3,1.85) {Molécule fille 1};
\node[popblue, anchor=west] at (9.7,1.6) {\scriptsize brin parental 1};
\node[poporange, anchor=west] at (9.7,2.1) {\scriptsize brin néoformé};
% Molécule fille 2
\draw[line width=2.2pt, popblue] (6.4,-1.6) -- (9.6,-1.6);
\draw[line width=2.2pt, poporange] (6.4,-1.1) -- (9.6,-1.1);
\node[anchor=east, align=right] at (6.3,-1.35) {Molécule fille 2};
\node[popblue, anchor=west] at (9.7,-1.6) {\scriptsize brin parental 2};
\node[poporange, anchor=west] at (9.7,-1.1) {\scriptsize brin néoformé};
% Accolades identiques
\node[popgreen, align=center] at (8,-2.6) {\textbf{Deux molécules filles identiques à la molécule mère}};
\end{tikzpicture}
\end{center}
\legende{Principe de la réplication semi-conservative. Les deux brins parentaux (bleu) se séparent et chacun sert de matrice à la synthèse d'un brin nouveau (orange). Chaque molécule fille conserve un brin parental et reçoit un brin néoformé : la copie est dite « semi-conservative ».}
\end{popfigure}

\courssub{La preuve expérimentale du mode semi-conservatif}

\begin{experience}[L'expérience de Meselson et Stahl (1958), présentée simplement]
\textbf{Problème.} Trois hypothèses étaient envisageables pour le mode de réplication : \emph{conservatif} (la molécule mère reste intacte et une molécule entièrement neuve est fabriquée), \emph{dispersif} (les molécules filles sont des mosaïques de fragments anciens et neufs) ou \emph{semi-conservatif} (un brin ancien + un brin neuf par molécule fille).

\textbf{Protocole (principe).} Meselson et Stahl cultivent des bactéries pendant plusieurs générations dans un milieu contenant de l'azote « lourd » ($^{15}$N) : tout leur ADN devient lourd. Ils transfèrent ensuite les bactéries dans un milieu à azote « léger » ($^{14}$N) et prélèvent des échantillons à chaque génération. L'ADN est séparé selon sa densité par centrifugation : l'ADN lourd migre plus bas que l'ADN léger.

\textbf{Observations.}
\begin{itemize}
\item Génération 0 : une seule bande, ADN lourd ($^{15}$N/$^{15}$N).
\item Après 1 génération : une seule bande, de densité \textbf{intermédiaire} (ADN hybride $^{15}$N/$^{14}$N).
\item Après 2 générations : \textbf{deux bandes} — une bande intermédiaire et une bande légère ($^{14}$N/$^{14}$N), en quantités égales.
\end{itemize}

\textbf{Interprétation.} L'absence de bande lourde dès la première génération élimine l'hypothèse conservative. L'apparition d'une bande légère à la deuxième génération, alors que la moitié des molécules reste hybride, élimine l'hypothèse dispersive.

\textbf{Conclusion.} Seule l'hypothèse \textbf{semi-conservative} rend compte des résultats : chaque molécule fille conserve un brin parental (lourd, marqué) et synthétise un brin nouveau (léger). Le marquage de l'ADN (ici par isotope, ailleurs par des précurseurs radioactifs comme la thymidine tritiée) est donc un outil expérimental majeur pour suivre le devenir des molécules d'ADN au cours des générations cellulaires.
\end{experience}

\courssub{Conséquence chromosomique : de 1 chromatide à 2 chromatides sœurs}

L'ADN nageant librement serait impossible à partager proprement : il est donc organisé en \cle{chromosomes}. Avant la phase S (en G1), chaque chromosome est formé d'\textbf{une seule molécule d'ADN}, c'est-à-dire d'\textbf{une seule chromatide}. Après la réplication (en G2), chaque chromosome est formé de \textbf{deux molécules d'ADN identiques}, c'est-à-dire de \textbf{deux chromatides identiques} appelées \cle{chromatides sœurs}, qui restent unies au niveau d'une région rétrécie : le \cle{centromère}.

\begin{attention}[Nombre de chromosomes ou quantité d'ADN ?]
Piège classique du Probatoire ! Après la phase S :
\begin{itemize}
\item le \textbf{nombre de chromosomes} est \textbf{inchangé} (un chromosome à 2 chromatides reste \emph{un} chromosome, car il n'a qu'\emph{un} centromère) ;
\item la \textbf{quantité d'ADN} a \textbf{doublé}, car chaque chromosome contient désormais deux molécules d'ADN.
\end{itemize}
On compte les chromosomes par leurs centromères, jamais par leurs chromatides.
\end{attention}

\begin{popfigure}
\begin{center}
\begin{tikzpicture}[scale=1.05, every node/.style={font=\small}]
% Chromosome 1 chromatide
\draw[line width=3pt, popblue, rounded corners=2pt] (0,1.4) -- (0,-1.4);
\fill[popdark] (0,0) circle (0.09);
\node[align=center] at (0,-2) {Chromosome à \textbf{1 chromatide}\\ (phase G1 : 1 molécule d'ADN)};
\draw[thin, popmuted] (0.1,0.6) -- (1.6,1.3);
\node[anchor=west, popmuted] at (1.6,1.3) {\scriptsize chromatide};
\draw[thin, popmuted] (0.1,0) -- (1.6,0.2);
\node[anchor=west, popmuted] at (1.6,0.2) {\scriptsize centromère};
% Flèche phase S
\draw[-{Stealth[length=3mm]}, line width=1.3pt, poporange] (3.1,0) -- (4.5,0);
\node[poporange, align=center] at (3.8,0.55) {\textbf{phase S}\\ \scriptsize réplication de l'ADN};
% Chromosome 2 chromatides
\draw[line width=3pt, popblue, rounded corners=2pt] (5.6,1.4) .. controls (5.6,0.35) and (5.85,0.35) .. (5.85,0) .. controls (5.85,-0.35) and (5.6,-0.35) .. (5.6,-1.4);
\draw[line width=3pt, poporange, rounded corners=2pt] (6.6,1.4) .. controls (6.6,0.35) and (6.35,0.35) .. (6.35,0) .. controls (6.35,-0.35) and (6.6,-0.35) .. (6.6,-1.4);
\fill[popdark] (6.1,0) ellipse (0.28 and 0.12);
\node[align=center] at (6.1,-2) {Chromosome à \textbf{2 chromatides sœurs}\\ (phase G2 : 2 molécules d'ADN identiques)};
\draw[thin, popmuted] (6.7,0.8) -- (8.2,1.3);
\node[anchor=west, popmuted] at (8.2,1.3) {\scriptsize chromatides sœurs};
\draw[thin, popmuted] (6.45,0) -- (8.2,0.2);
\node[anchor=west, popmuted] at (8.2,0.2) {\scriptsize centromère};
\end{tikzpicture}
\end{center}
\legende{Conséquence chromosomique de la réplication. Pendant la phase S, la molécule d'ADN de chaque chromosome est dupliquée : le chromosome passe de 1 chromatide (G1) à 2 chromatides sœurs identiques (G2), réunies par le centromère. Le nombre de chromosomes ne change pas, mais la quantité d'ADN double.}
\end{popfigure}

\begin{exemplebox}[Le cas chiffré de la cellule humaine]
Une cellule humaine en phase G1 possède \textbf{46 chromosomes à 1 chromatide}, soit une quantité d'ADN notée \cle{Q}. Pendant la phase S, la réplication double l'ADN : en phase G2, la cellule possède \textbf{46 chromosomes à 2 chromatides}, soit une quantité d'ADN égale à \cle{2Q}. À la fin de la mitose, chaque cellule fille reçoit 46 chromosomes à 1 chromatide, soit Q : la quantité d'ADN et le nombre de chromosomes sont ainsi \textbf{maintenus constants} d'une génération cellulaire à l'autre. C'est la \cle{conservation du caryotype}.
\end{exemplebox}

\courssub{Le lien fondamental : réplication fidèle et transmission intégrale de l'information}

\begin{methode}[Établir le lien « réplication $\rightarrow$ conservation du caryotype » en 3 étapes]
Face à une question du type « montrer que la mitose conserve le patrimoine génétique », raisonne toujours ainsi :
\begin{enumerate}
\item \textbf{Duplication de l'ADN (phase S)} : la réplication semi-conservative produit deux molécules d'ADN identiques pour chaque chromosome ; chaque chromosome comporte alors 2 chromatides sœurs.
\item \textbf{Deux chromatides identiques} : comme chaque chromatide sœur porte une copie fidèle de la même information, la cellule dispose de \emph{deux jeux complets et identiques} du patrimoine génétique.
\item \textbf{Séparation équitable (mitose)} : la mitose distribue une chromatide de chaque chromosome à chaque cellule fille ; chacune reçoit donc un jeu complet, identique à celui de la cellule mère.
\end{enumerate}
Conclusion : la réplication fidèle de l'ADN, suivie du partage équitable des chromatides, garantit la \textbf{transmission intégrale de l'information génétique} aux deux cellules filles, donc le \textbf{maintien de l'identité biologique} de l'organisme.
\end{methode}

\begin{aretenir}[L'essentiel de la section]
\begin{itemize}
\item L'information génétique est portée par la séquence des nucléotides de l'ADN ; les deux brins sont complémentaires (A-T, G-C).
\item La \textbf{réplication} a lieu en \textbf{phase S} de l'interphase ; elle est \textbf{semi-conservative} : chaque molécule fille = 1 brin parental + 1 brin néoformé (preuve : expérience de Meselson et Stahl).
\item Conséquence : chaque chromosome passe de \textbf{1 chromatide} à \textbf{2 chromatides sœurs} réunies par le centromère ; la quantité d'ADN double (Q $\rightarrow$ 2Q) mais le nombre de chromosomes est inchangé.
\item La fidélité de la réplication est la condition de la \textbf{conservation du caryotype} et de l'identité biologique des cellules filles.
\end{itemize}
\end{aretenir}

\begin{exoresolu}[Quantité d'ADN au cours du cycle cellulaire]
\textbf{Énoncé.} On mesure la quantité d'ADN par noyau dans une culture de cellules de peau. On trouve trois populations de cellules : population X avec \SI{6,6}{pg} d'ADN, population Y avec des valeurs comprises entre \SI{6,6}{pg} et \SI{13,2}{pg}, population Z avec \SI{13,2}{pg}. 1) Attribuer à chaque population la phase du cycle cellulaire correspondante. 2) Expliquer pourquoi la population Y présente des valeurs intermédiaires variables. 3) Quelle sera la quantité d'ADN par noyau dans les cellules issues de la division des cellules Z ? Justifier.
\tcblower
\textbf{Solution détaillée.}
\begin{enumerate}
\item La population X (\SI{6,6}{pg}) correspond aux cellules en \textbf{phase G1} : 46 chromosomes à 1 chromatide, quantité Q. La population Z (\SI{13,2}{pg} $= 2 \times$ \SI{6,6}{pg}) correspond aux cellules en \textbf{phase G2} (ou en mitose avant la séparation finale) : 46 chromosomes à 2 chromatides, quantité 2Q. La population Y, aux valeurs intermédiaires, correspond aux cellules en \textbf{phase S}.
\item En phase S, la réplication est \textbf{en cours} : selon le moment où la cellule est observée, une fraction plus ou moins grande de son ADN est déjà dupliquée. D'où des quantités d'ADN variables, toutes comprises entre Q et 2Q.
\item Les cellules Z achèvent la mitose : les chromatides sœurs de chaque chromosome se séparent et sont réparties équitablement entre les deux cellules filles. Chaque cellule fille reçoit donc 46 chromosomes à 1 chromatide, soit \SI{6,6}{pg} d'ADN (Q). La réplication fidèle suivie du partage équitable assure ainsi le maintien de la quantité d'ADN, donc de l'information génétique, d'une génération cellulaire à l'autre.
\end{enumerate}
\end{exoresolu}

\begin{savaistu}[Réplication et médecine au Cameroun]
La fidélité de la réplication n'est pas absolue : des erreurs rares (mutations) peuvent échapper aux mécanismes de réparation. Si elles touchent des gènes qui contrôlent le cycle cellulaire, la cellule peut se diviser sans frein : c'est le point de départ d'un \textbf{cancer}, que nous étudierons à la fin de cette leçon. À l'inverse, de nombreux médicaments anticancéreux et antiviraux utilisés dans les hôpitaux camerounais agissent précisément en \textbf{bloquant la réplication de l'ADN} des cellules tumorales ou des virus : comprendre la réplication, c'est comprendre le principe de ces traitements.
\end{savaistu}

\coursec{Les phases de la mitose}

Dans la section précédente, nous avons vu que pendant l'interphase, et plus précisément pendant la phase S, la molécule d'ADN de chaque chromosome est \cle{répliquée} : chaque chromosome passe ainsi de l'état « une chromatide » à l'état « deux chromatides sœurs identiques » réunies par le centromère. La cellule possède donc, à l'issue de l'interphase, un double exemplaire de toute son information génétique. Mais ce doublement ne suffit pas : il faut encore \emph{distribuer équitablement} ces deux exemplaires entre deux futures cellules filles. C'est précisément le rôle de la mitose, véritable « chorégraphie chromosomique » que nous allons maintenant décrire phase par phase.

\courssub{Qu'est-ce que la mitose ?}

\textbf{Une situation concrète pour commencer.} Lorsque tu te coupes au doigt, la plaie se referme en quelques jours : des cellules de la peau se sont multipliées pour remplacer les cellules détruites. Or la peau reconstituée est identique à la peau d'origine — mêmes cellules, mêmes caractères. De même, un bout de tige de manioc planté dans un champ du Cameroun donne un plant entier, génétiquement identique au plant mère. Comment chaque cellule fille peut-elle recevoir \emph{exactement} la même information génétique que la cellule mère ?

\begin{definition}[La mitose]
La \cle{mitose} (ou caryocinèse) est l'ensemble des phénomènes conduisant à la \textbf{répartition équitable des chromosomes} entre deux noyaux fils génétiquement identiques au noyau mère. Elle est suivie de la \cle{cytodiérèse}, division du cytoplasme, qui aboutit à la formation de deux cellules filles. La mitose se déroule classiquement en quatre phases : \textbf{prophase, métaphase, anaphase, télophase}.
\end{definition}

\begin{attention}[Ce que la mitose n'est pas]
La mitose est la division du \textbf{noyau}, pas de la cellule entière. La division du cytoplasme (cytodiérèse) est un phénomène distinct qui la suit. Au Probatoire, ne réduis jamais la mitose à « la division cellulaire » : c'est une imprécision sanctionnée.
\end{attention}

\courssub{Comment observe-t-on les phases de la mitose ?}

\begin{experience}[Observation de la mitose en microscopie]
\textbf{Matériel :} pointes de racines d'oignon ou de jeunes pousses (zones de forte multiplication cellulaire), colorant (acétocarmin ou vert de méthyle), lame et lamelle, microscope optique.

\textbf{Protocole :} on fixe et on colore les cellules de la pointe de racine, puis on les écrase entre lame et lamelle pour les étaler en une couche fine (technique de l'écrasement), et on observe au fort grossissement (objectif $\times 40$ ou $\times 100$).

\textbf{Observations :} les cellules ne présentent pas toutes le même aspect. Certaines montrent un noyau intact ; d'autres montrent des bâtonnets colorés condensés, tantôt dispersés, tantôt alignés au centre, tantôt séparés en deux lots qui migrent vers les extrémités.

\textbf{Interprétation :} chaque aspect correspond à une étape figée de la mitose. En classant les images, on reconstitue la chronologie : prophase $\rightarrow$ métaphase $\rightarrow$ anaphase $\rightarrow$ télophase. La majorité des cellules étant en interphase, on en déduit aussi que l'interphase est la phase la plus longue du cycle cellulaire.
\end{experience}

\courssub{Prophase : la condensation des chromosomes}

Au début de la prophase, la chromatine, jusque-là décondensée et invisible, se \cle{condense} progressivement : les filaments d'ADN s'enroulent et se tassent. Les \textbf{chromosomes deviennent visibles} au microscope sous forme de bâtonnets, chacun constitué de \textbf{deux chromatides sœurs} réunies au niveau du centromère. Simultanément :
\begin{itemize}
\item le \textbf{nucléole disparaît} (arrêt de la synthèse des ribosomes) ;
\item l'\textbf{enveloppe nucléaire se désorganise} et disparaît en fin de prophase (prométaphase), libérant les chromosomes dans le cytoplasme ;
\item le \textbf{fuseau achromatique} (ou fuseau mitotique) commence à se former : des faisceaux de microtubules s'organisent entre les deux centrosomes (pôles cellulaires), qui migrent vers les pôles opposés de la cellule.
\end{itemize}

\courssub{Métaphase : l'alignement équatorial}

Les chromosomes, toujours à deux chromatides, sont capturés par les \textbf{fibres du fuseau} (microtubules kinétochoriens) qui se fixent sur leurs \textbf{kinétochores}, structures protéiques situées au niveau du centromère. Sous l'action de ces fibres, les chromosomes se déplacent et viennent s'\textbf{aligner sur un plan équatorial} de la cellule, appelé \cle{plaque équatoriale} (ou plaque métaphasique). C'est le stade où les chromosomes sont les plus condensés et les plus faciles à observer : c'est donc en métaphase que l'on réalise les \textbf{caryotypes}.

\courssub{Anaphase : la séparation des chromatides sœurs}

C'est l'étape décisive. Les centromères se divisent et les \textbf{chromatides sœurs de chaque chromosome se séparent}. Sous l'action des fibres du fuseau qui se raccourcissent (dépolymérisation des microtubules au niveau des kinétochores), chaque chromatide — devenue un chromosome à une chromatide — \textbf{migre vers un pôle opposé} de la cellule. À la fin de l'anaphase, chaque pôle a reçu un \textbf{lot complet et identique} de chromosomes : c'est ici que s'effectue la répartition équitable de l'information génétique.

\begin{attention}[L'erreur classique du Probatoire]
C'est en \textbf{anaphase} (et non en métaphase) que les chromatides sœurs se séparent. En métaphase, elles sont encore \emph{unies} par le centromère. Critère de reconnaissance : si l'on voit des chromosomes simples (en V ou en J) qui migrent vers les pôles, c'est l'anaphase ; si l'on voit des chromosomes doubles (en X) alignés au centre, c'est la métaphase.
\end{attention}

\courssub{Télophase : la naissance de deux noyaux fils}

Les chromosomes parvenus aux pôles se \textbf{décondensent} progressivement et redeviennent invisibles (chromatine). Autour de chaque lot de chromosomes, une \textbf{enveloppe nucléaire se reforme} à partir de vésicules du réticulum endoplasmique, et le nucléole réapparaît (reprise de la transcription de l'ARNr). Le fuseau disparaît. La cellule contient alors \textbf{deux noyaux fils}, chacun possédant le même nombre de chromosomes et la même information génétique que le noyau mère.

\courssub{La cytodiérèse : le partage du cytoplasme}

La cytodiérèse partage le cytoplasme et les organites entre les deux noyaux fils, donnant naissance à deux cellules filles. Le mécanisme diffère selon le type cellulaire :
\begin{itemize}
\item chez la \textbf{cellule animale}, un anneau de filaments d'actine et de myosine (anneau contractile) étrangle la cellule en son équateur : c'est le \cle{sillon de division}, qui se resserre jusqu'à séparer les deux cellules filles ;
\item chez la \textbf{cellule végétale}, la paroi rigide (cellulosique) empêche tout étranglement : des vésicules issues de l'appareil de Golgi s'accumulent au centre (phragmoplaste) et fusionnent pour former une \cle{plaque cellulaire}, qui s'étend du centre vers la périphérie et devient la nouvelle paroi (lamelle moyenne) séparant les deux cellules filles.
\end{itemize}

\begin{popfigure}
\begin{center}
\begin{tikzpicture}[scale=0.55, every node/.style={font=\scriptsize}, >=Stealth]
% Styles
\tikzset{
  chromo/.style={poporange, very thick, line cap=round},
  spindle/.style={popgreen!70!black, thin},
  centro/.style={popdark, fill=popdark},
  nucenv/.style={popblue, thick},
  nucenvdash/.style={popmuted, dashed, thick},
  labelptr/.style={popdark, thin, ->, shorten >=1pt, shorten <=1pt},
  annote/.style={popdark, font=\tiny, align=center}
}

% PROPHASE (xshift=0)
\begin{scope}[xshift=0cm]
\draw[nucenvdash] (0,0) circle (1.6); % Enveloppe en disparition
\draw[chromo] (-0.6,0.5) -- (-0.3,0.8) -- (0,0.5); \draw[chromo] (-0.6,0.5) -- (-0.3,0.2) -- (0,0.5);
\draw[chromo] (0.7,-0.6) -- (1.0,-0.3) -- (1.3,-0.6); \draw[chromo] (0.7,-0.6) -- (1.0,-0.9) -- (1.3,-0.6);
\draw[chromo] (-1.0,-0.7) -- (-0.7,-0.4) -- (-0.4,-0.7); \draw[chromo] (-1.0,-0.7) -- (-0.7,-1.0) -- (-0.4,-0.7);
\draw[chromo] (0.4,0.9) -- (0.7,1.2) -- (1.0,0.9); \draw[chromo] (0.4,0.9) -- (0.7,0.6) -- (1.0,0.9);
\fill[centro] (-1.9,0) circle (0.07); \fill[centro] (1.9,0) circle (0.07);
\draw[spindle] (-1.9,0) -- (-1.0,0.2); \draw[spindle] (-1.9,0) -- (-1.0,-0.2);
\draw[spindle] (1.9,0) -- (1.0,0.2); \draw[spindle] (1.9,0) -- (1.0,-0.2);
\node at (0,-2.3) {\textbf{Prophase}};
\node[annote, align=center] at (2.5, 1.5){Chromosomes\\se condensent};
\node[annote, align=center] at (2.5, -1.5){Centrosomes\\migrent};
\end{scope}

% METAPHASE (xshift=3.8)
\begin{scope}[xshift=3.8cm]
\draw[nucenv] (-2.2,-1.5) rectangle (2.2,1.5); % Limite cellule
\foreach \y in {-0.5, 0.5}{
  \draw[chromo] (-0.8,\y) -- (-0.5,\y+0.3) -- (-0.2,\y); \draw[chromo] (-0.8,\y) -- (-0.5,\y-0.3) -- (-0.2,\y);
  \draw[chromo] (0.4,\y) -- (0.7,\y+0.3) -- (1.0,\y); \draw[chromo] (0.4,\y) -- (0.7,\y-0.3) -- (1.0,\y);
}
\fill[centro] (-2.2,0) circle (0.07); \fill[centro] (2.2,0) circle (0.07);
\foreach \x in {-0.5, 0.7}{
  \draw[spindle] (-2.2,0) -- (\x,0.15); \draw[spindle] (-2.2,0) -- (\x,-0.15);
  \draw[spindle] (2.2,0) -- (\x,0.15); \draw[spindle] (2.2,0) -- (\x,-0.15);
}
\draw[popmuted, dashed, thick] (-1.3,0.15) -- (1.3,0.15) node[midway, above=2pt, font=\tiny, popmuted] {Plaque équatoriale};
\node at (0,-2.3) {\textbf{Métaphase}};
% Annotations Métaphase
\draw[labelptr] (1.2, 0.8) -- (2.0, 1.3) node[annote, anchor=west, align=center]{Chromosome\\à 2 chromatides};
\draw[labelptr] (0.7, 0.5) -- (1.5, 0.8) node[annote, anchor=west] {Centromère};
\draw[labelptr] (-2.2, 0) -- (-2.8, 0.5) node[annote, anchor=east, align=center]{Centrosome\\(pôle)};
\draw[labelptr] (-1.5, 0) -- (-2.2, -0.5) node[annote, anchor=north, align=center]{Fibres du\\fuseau};
\end{scope}

% ANAPHASE (xshift=7.6)
\begin{scope}[xshift=7.6cm]
\draw[nucenv] (-2.5,-1.5) rectangle (2.5,1.5);
\foreach \s in {1,-1}{
  % Lot gauche (\s=-1) et droit (\s=1) - chromosomes simples (V)
  \draw[chromo] (\s*0.6,0.4) -- (\s*0.9,0.1);
  \draw[chromo] (\s*0.6,0.4) -- (\s*0.9,0.7);
  \draw[chromo] (\s*0.6,-0.4) -- (\s*0.9,-0.1);
  \draw[chromo] (\s*0.6,-0.4) -- (\s*0.9,-0.7);
  \draw[chromo] (\s*1.4,0.4) -- (\s*1.7,0.1);
  \draw[chromo] (\s*1.4,0.4) -- (\s*1.7,0.7);
  \draw[chromo] (\s*1.4,-0.4) -- (\s*1.7,-0.1);
  \draw[chromo] (\s*1.4,-0.4) -- (\s*1.7,-0.7);
}
\fill[centro] (-2.5,0) circle (0.07); \fill[centro] (2.5,0) circle (0.07);
\draw[labelptr] (-0.5, 1.0) -- (-1.5, 1.3) node[annote, anchor=east, align=center]{Migration\\vers les pôles};
\draw[labelptr] (0.5, 1.0) -- (1.5, 1.3) node[annote, anchor=west, align=center]{Chromosomes\\simples (1 chromatide)};
\node at (0,-2.3) {\textbf{Anaphase}};
\end{scope}

% TELOPHASE (xshift=11.4)
\begin{scope}[xshift=11.4cm]
\draw[nucenv] (-2.5,-1.5) rectangle (2.5,1.5);
\draw[nucenv] (-1.0,0) circle (0.7);
\draw[nucenv] (1.0,0) circle (0.7);
\foreach \s in {1,-1}{
  \draw[chromo, thin] (\s*1.0,0.3) arc (30:150:0.3);
  \draw[chromo, thin] (\s*1.0,-0.3) arc (-30:-150:0.3);
  \fill[popdark] (\s*1.0, 0.3) circle (0.04); % Nucléole
}
\node at (0,-2.3) {\textbf{Télophase}};
\node[annote, align=center] at (0, 1.8){Décondensation\\Réformation enveloppes};
\end{scope}

% CYTODIERESE (xshift=15.0)
\begin{scope}[xshift=15.0cm]
\draw[nucenv] (-1.0,0) circle (1.1);
\draw[nucenv] (1.0,0) circle (1.1);
\draw[nucenv] (-1.0,0) circle (0.45);
\draw[nucenv] (1.0,0) circle (0.45);
\node at (0,-2.3) {\textbf{Cytodiérèse}};
\node[annote, align=center] at (0, 1.8){Deux cellules\\filles distinctes};
\end{scope}
\end{tikzpicture}
\end{center}
\legende{Les quatre phases de la mitose et la cytodiérèse dans une cellule animale à $2n = 4$ chromosomes. En orange : les chromosomes ; points noirs : centrosomes (pôles) ; traits verts : fibres du fuseau (microtubules) ; traits bleus : enveloppe nucléaire / membrane cellulaire ; traits gris pointillés : plaque équatoriale. En prophase les chromosomes se condensent et l'enveloppe se fragmente ; en métaphase les chromosomes (doubles) s'alignent sur la plaque équatoriale ; en anaphase les chromatides sœurs (devenues chromosomes simples) migrent vers les pôles ; en télophase deux noyaux fils se reforment ; la cytodiérèse achève la séparation en deux cellules filles.}
\end{popfigure}

\begin{popfigure}
\begin{center}
\begin{tikzpicture}[scale=0.85, every node/.style={font=\scriptsize}, >=Stealth]
% Cellule animale : sillon
\begin{scope}[xshift=0cm]
\draw[popblue, very thick] plot[smooth cycle, tension=1.2] coordinates {(0,1.4) (1.2,1.0) (1.4,0) (1.2,-1.0) (0,-1.4) (-1.2,-1.0) (-1.4,0) (-1.2,1.0)};
\draw[popblue, very thick] (-1.4,0) .. controls (-0.5,0.3) and (0.5,0.3) .. (1.4,0);
\draw[popblue, very thick] (-1.4,0) .. controls (-0.5,-0.3) and (0.5,-0.3) .. (1.4,0);
\draw[popmuted] (-0.7,0) circle (0.45);
\draw[popmuted] (0.7,0) circle (0.45);
\draw[poppink, thick, ->] (-2.1,0.7) -- (-1.45,0.2);
\draw[poppink, thick, ->] (2.1,0.7) -- (1.45,0.2);
\node[poppink] at (0,1.9) {\textbf{Cellule animale : sillon de division}};
\node[popmuted, align=center] at (0,-1.9){Étranglement équatorial de la membrane\\par anneau contractile (actine/myosine)};
\end{scope}
% Cellule végétale : plaque cellulaire
\begin{scope}[xshift=7.5cm]
\draw[popgreen!60!black, very thick, rounded corners=2pt] (-1.7,-1.4) rectangle (1.7,1.4); % Paroi mère
\draw[popblue, thick, rounded corners=2pt] (-1.6,-1.3) rectangle (1.6,1.3); % Membrane
\draw[popmuted] (-0.9,0) circle (0.45);
\draw[popmuted] (0.9,0) circle (0.45);
\draw[poppink, very thick, dashed] (-0.4,0) -- (0.4,0);
\draw[poppink, thick, ->] (-0.4,0) -- (-1.65,0);
\draw[poppink, thick, ->] (0.4,0) -- (1.65,0);
\node[poppink] at (0,1.9) {\textbf{Cellule végétale : plaque cellulaire}};
\node[popmuted, align=center] at (0,-1.9){Vésicules de Golgi $\rightarrow$ Phragmoplaste $\rightarrow$\\Plaque cellulaire (lamelle moyenne)};
\end{scope}
\end{tikzpicture}
\end{center}
\legende{Comparaison de la cytodiérèse. À gauche, chez la cellule animale (ex. épithélium), la membrane plasmique s'étrangle en un sillon de division (flèches roses) grâce à un anneau contractile. À droite, chez la cellule végétale (ex. méristème de racine de maïs), entourée d'une paroi rigide (vert foncé), des vésicules golgiennes fusionnent au plan équatorial pour former la plaque cellulaire (pointillés roses) qui s'étend du centre vers la périphérie.}
\end{popfigure}

\courssub{Comment identifier une phase sur une photographie ou un schéma ?}

\begin{methode}[Les trois critères d'identification]
Face à une photographie de microscopie ou un schéma, observe successivement trois critères :
\begin{enumerate}
\item \textbf{L'enveloppe nucléaire} : présente (interphase ou télophase) ou absente (métaphase, anaphase) ; en cours de fragmentation (prophase).
\item \textbf{La position des chromosomes} : dispersés dans le noyau (prophase), alignés sur la plaque équatoriale (métaphase), en deux lots migrant vers les pôles (anaphase), regroupés aux pôles (télophase).
\item \textbf{L'état des chromatides} : chromosomes à deux chromatides (prophase, métaphase) ou chromatides séparées devenues chromosomes simples (anaphase, télophase).
\end{enumerate}
La combinaison des trois réponses désigne la phase sans ambiguïté.
\end{methode}

\begin{aretenir}[Moyen mnémotechnique : P-M-A-T]
Retiens l'ordre des phases grâce à la suite \textbf{P-M-A-T} :
\begin{itemize}
\item \textbf{P}rophase : les chromosomes se \emph{condensent} (deviennent visibles), l'enveloppe nucléaire se fragmente ;
\item \textbf{M}étaphase : les chromosomes s'alignent au \emph{milieu} (plaque équatoriale), centromères sur le plan équatorial ;
\item \textbf{A}naphase : les chromatides s'\emph{écartent} vers les pôles opposés (séparation des centromères) ;
\item \textbf{T}élophase : \emph{deux noyaux} fils se reforment (décondensation, réapparition nucléoles).
\end{itemize}
\end{aretenir}

\begin{exoresolu}[Identifier une phase et justifier]
Sur une photographie de cellule de pointe de racine d'oignon (\textit{Allium cepa}, $2n = 16$), on observe des bâtonnets colorés en forme de X alignés sur le plan médian de la cellule ; aucune enveloppe nucléaire n'est visible. (a) Identifier la phase. (b) Justifier la réponse à l'aide des trois critères. (c) Combien de chromatides la cellule contient-elle à ce stade ? (d) Combien de chromosomes possédera chaque noyau fils ?
\tcblower
\textbf{(a)} Il s'agit de la \textbf{métaphase}.

\textbf{(b)} Application des trois critères :
\begin{itemize}
\item enveloppe nucléaire : \emph{absente} $\rightarrow$ prophase achevée, métaphase ou anaphase ;
\item position des chromosomes : \emph{alignés sur la plaque équatoriale} $\rightarrow$ métaphase ;
\item état des chromatides : les chromosomes en X sont \emph{à deux chromatides}, encore unies par le centromère $\rightarrow$ ce n'est pas l'anaphase.
\end{itemize}

\textbf{(c)} Chacun des 16 chromosomes comporte 2 chromatides sœurs : la cellule contient donc $16 \times 2 = 32$ \textbf{chromatides}, réparties en 16 chromosomes doubles (repliés).

\textbf{(d)} À l'issue de la mitose, chaque noyau fils reçoit 16 chromosomes simples (une chromatide chacun), soit le même caryotype ($2n=16$) que la cellule mère.
\end{exoresolu}

\courssub{Bilan : l'équipartition de l'information génétique}

\begin{aretenir}[Bilan de la mitose]
À l'issue de la mitose et de la cytodiérèse, une cellule mère à $2n$ chromosomes donne \textbf{deux cellules filles} possédant chacune :
\begin{itemize}
\item le \textbf{même nombre de chromosomes} ($2n$) que la cellule mère ;
\item la \textbf{même information génétique}, grâce à la réplication de l'ADN en interphase (phase S) puis à la séparation équitable des chromatides sœurs en anaphase.
\end{itemize}
La mitose assure ainsi la \cle{conservation du caryotype} et l'identité biologique des cellules au fil des générations cellulaires : c'est le fondement de la croissance, du renouvellement des tissus (cicatrisation, hématopoïèse) et de la reproduction asexuée (bouturage du manioc, du bananier, marcottage).
\end{aretenir}

\begin{savaistu}[Une heure pour diviser un noyau]
Chez une cellule humaine en culture (fibroblaste), la mitose ne dure qu'environ \textbf{1 heure} (l'anaphase, elle, n'excède pas 10 à 30 minutes), alors que l'interphase peut durer \textbf{plus de 20 heures}. C'est pourquoi, sur un cliché de tissu en multiplication, l'immense majorité des cellules observées sont en interphase : les phases de la mitose, bien que spectaculaires, ne représentent qu'une fraction minime (environ 5\%) de la durée du cycle cellulaire. Certaines cellules très actives, comme celles de la moelle osseuse rouge (hématopoïèse) qui produisent chaque jour environ 200 milliards de globules rouges, ou les cellules de la base de l'épiderme, réalisent cet exploit en continu, tout au long de notre vie.
\end{savaistu}

\begin{exercice}[À toi de jouer]
Une cellule animale possède $2n = 8$ chromosomes. (1) Combien de chromosomes et de chromatides comporte-t-elle en métaphase ? (2) Combien de chromosomes migrent vers chaque pôle en anaphase ? (3) Quel est le nombre de chromosomes de chaque cellule fille ? (4) Citer une différence fondamentale entre la cytodiérèse de cette cellule et celle d'une cellule de racine de maïs ($2n = 20$).
\end{exercice}

\begin{corrige}[Exercice]
(1) En métaphase : 8 chromosomes, chacun constitué de 2 chromatides sœurs, soit 16 chromatides au total. (2) En anaphase, les chromatides sœurs se séparent au niveau du centromère : 8 chromosomes simples (à une chromatide) migrent vers chaque pôle. (3) Chaque cellule fille possède 8 chromosomes (à une chromatide initialement), identiques à ceux de la cellule mère ($2n=8$). (4) Chez la cellule animale, la cytodiérèse se fait par étranglement de la membrane plasmique (sillon de division) grâce à un anneau contractile actine/myosine ; chez la cellule végétale de maïs, la paroi cellulosique rigide empêche l'étranglement : la cytodiérèse se fait par formation d'une plaque cellulaire (fusion de vésicules golgiennes au phragmoplaste) qui s'étend radialement pour devenir la nouvelle paroi séparatrice (lamelle moyenne).
\end{corrige}

\coursec{Comportement des chromosomes et conservation du caryotype}

Dans la section précédente, nous avons décrit le déroulement de la mitose : prophase, métaphase, anaphase et télophase se succèdent comme les actes d'une pièce bien réglée. Mais une question fondamentale demeure : \emph{que deviennent exactement les chromosomes au cours de cette pièce ?} Combien y en a-t-il à chaque instant ? Et surtout, comment deux cellules filles peuvent-elles hériter d'un patrimoine génétique rigoureusement identique à celui de la cellule mère ? C'est à ces questions quantitatives que répond cette section. Nous allons suivre, chromosome par chromosome et chromatide par chromatide, le devenir du matériel génétique.

\courssub{Observer pour comprendre : le problème de la cicatrisation}

Lorsqu'un enfant de Yaoundé se blesse au genou en jouant au football, la plaie se referme en quelques jours : de nouvelles cellules épidermiques remplacent les cellules détruites. Or ces nouvelles cellules fonctionnent exactement comme les anciennes. Comment est-ce possible ? Si chaque division cellulaire « diluait » ou modifiait le patrimoine génétique, les cellules de réparation seraient différentes des cellules d'origine, et la cicatrisation produirait un tissu défectueux.

\begin{itemize}
\item \textbf{Observation :} les cellules d'un même tissu, avant et après des milliers de divisions, présentent le même nombre et la même forme de chromosomes (le même \cle{caryotype}).
\item \textbf{Hypothèse :} la mitose comporte un mécanisme qui garantit que chaque cellule fille reçoit un exemplaire complet et fidèle du matériel génétique.
\item \textbf{Vérification :} suivons quantitativement les chromosomes et l'ADN au cours du cycle cellulaire.
\end{itemize}

\courssub{Suivi quantitatif des chromosomes et de l'ADN au cours du cycle}

Pour raisonner clairement, il faut distinguer trois grandeurs que l'on confond trop souvent :

\begin{definition}[Les trois grandeurs à compter]
\begin{itemize}
\item Le \cle{nombre de chromosomes} : on compte les \textbf{centromères}. Un chromosome à deux chromatides ne compte que pour \textbf{un} chromosome, car il ne possède qu'un seul centromère.
\item Le \cle{nombre de chromatides} : on compte les \textbf{brins d'ADN} (molécules d'ADN double brin). Après la réplication, chaque chromosome comporte deux chromatides sœurs identiques.
\item La \cle{quantité d'ADN} (notée $Q$) : quantité d'ADN d'une cellule en phase G1, avant réplication. Après la phase S, elle vaut $2Q$.
\end{itemize}
\end{definition}

\begin{methode}[Construire le tableau d'évolution sans se tromper]
\begin{enumerate}
\item \textbf{Compter d'abord les centromères} : cela donne le nombre de chromosomes. Tant que les chromatides sœurs restent unies par leur centromère, il n'y a qu'un chromosome.
\item \textbf{Compter ensuite les brins d'ADN} : chaque chromatide est une molécule d'ADN double brin. En G1 : 1 chromatide par chromosome ; de la fin de la phase S à la métaphase : 2 chromatides par chromosome.
\item \textbf{Repérer le moment du basculement} : c'est la \textbf{séparation des centromères en anaphase} qui double instantanément le nombre de chromosomes (chaque chromatide devient un chromosome à part entière) sans changer la quantité totale d'ADN dans la cellule.
\item \textbf{Terminer par la cytodiérèse} : le partage du cytoplasme divise par deux, dans chaque cellule fille, le nombre de chromosomes et la quantité d'ADN.
\end{enumerate}
\end{methode}

Appliquons cette méthode à une cellule mère à $2n = 4$ chromosomes (deux paires de chromosomes homologues), dont la quantité d'ADN en G1 vaut $Q$.

\begin{popfigure}
\begin{center}
\begin{tabular}{|l|c|c|c|}
\hline
\rowcolor{popblueL} \textbf{Étape du cycle} & \textbf{Chromosomes} & \textbf{Chromatides} & \textbf{Quantité d'ADN} \\
\hline
Interphase G1 & 4 (à 1 chromatide) & 4 & $Q$ \\
\hline
Fin de la phase S / G2 & 4 (à 2 chromatides) & 8 & $2Q$ \\
\hline
Prophase & 4 (à 2 chromatides) & 8 & $2Q$ \\
\hline
Métaphase & 4 (à 2 chromatides) & 8 & $2Q$ \\
\hline
Anaphase & 8 (à 1 chromatide) & 8 & $2Q$ \\
\hline
Télophase (avant cytodiérèse) & 8 répartis en 2 lots de 4 & 8 & $2Q$ \\
\hline
\rowcolor{popgreenL} Chaque cellule fille & 4 (à 1 chromatide) & 4 & $Q$ \\
\hline
\end{tabular}
\end{center}
\legende{Évolution du nombre de chromosomes, du nombre de chromatides et de la quantité d'ADN au cours du cycle cellulaire, pour une cellule mère à $2n = 4$. On remarque le doublement transitoire du nombre de chromosomes en anaphase (séparation des centromères) et le retour à l'état initial dans chaque cellule fille.}
\end{popfigure}

Deux enseignements majeurs se dégagent de ce tableau. D'une part, le \textbf{doublement de la quantité d'ADN} a lieu une seule fois, pendant la phase S de l'interphase : c'est la \cle{réplication}. D'autre part, le \textbf{doublement du nombre de chromosomes} en anaphase n'est qu'apparent et temporaire : il correspond à la séparation des chromatides sœurs, et il est immédiatement compensé par la répartition en deux cellules filles.

\courssub{Le moment clé : la séparation des chromatides sœurs en anaphase}

Toute la fidélité de la mitose repose sur deux événements complémentaires : la réplication, qui fabrique deux copies identiques de chaque chromosome, et l'\cle{équipartition}, qui distribue une copie à chaque cellule fille. Cette équipartition est réalisée en anaphase.

Au cours de la métaphase, chaque chromosome (formé de deux chromatides sœurs) se place sur la plaque équatoriale, son centromère étant accroché aux fibres du fuseau venant des deux pôles opposés. Au début de l'anaphase, les \textbf{centromères se scindent} : les chromatides sœurs se séparent et migrent vers les pôles opposés, tractées par le raccourcissement des fibres du fuseau. Chaque chromatide devient alors un \textbf{chromosome à une chromatide}.

\begin{propriete}[Conservation du caryotype par la mitose — démonstration guidée en classe]
Toute cellule issue d'une mitose possède le \textbf{même caryotype} (même nombre $2n$ de chromosomes, même morphologie chromosomique) que la cellule mère. La démonstration repose sur deux piliers :
\begin{enumerate}
\item \textbf{Réplication (phase S) :} chaque molécule d'ADN est dupliquée selon le mécanisme semi-conservatif ; chaque chromosome devient un chromosome à deux chromatides sœurs \emph{génétiquement identiques}. On passe de $2n$ chromosomes à 1 chromatide ($Q$) à $2n$ chromosomes à 2 chromatides ($2Q$).
\item \textbf{Équipartition (anaphase) :} la séparation des centromères envoie une et une seule chromatide de chaque chromosome vers chaque pôle. Chaque lot polaire contient donc $2n$ chromosomes à 1 chromatide, soit la totalité de l'information génétique en un seul exemplaire de chaque sorte.
\end{enumerate}
Après la télophase et la cytodiérèse, chaque cellule fille reçoit $2n$ chromosomes à 1 chromatide et une quantité d'ADN $Q$ : son caryotype est \textbf{identique} à celui de la cellule mère.
\end{propriete}

\courssub{Démonstration guidée : suivi d'une cellule mère à $2n = 4$}

Suivons concrètement le devenir de chaque chromatide dans une cellule mère possédant deux paires de chromosomes : une paire de grands chromosomes (un d'origine paternelle, un d'origine maternelle) et une paire de petits chromosomes.

\begin{popfigure}
\begin{center}
\begin{tikzpicture}[scale=0.9, every node/.style={font=\small}]
% ---- Cellule mère G1 ----
\draw[popblue, thick] (0,0) ellipse (2.1 and 1.4);
\draw[popdark, thick] (0,0) ellipse (1.5 and 0.95);
% grand chromosome 1 chromatide
\draw[line width=2.2pt, poporange] (-0.9,0.35) -- (-0.9,-0.35);
\draw[line width=2.2pt, poppink] (-0.45,0.35) -- (-0.45,-0.35);
% petit chromosome
\draw[line width=2.2pt, poporange] (0.45,0.18) -- (0.45,-0.18);
\draw[line width=2.2pt, poppink] (0.9,0.18) -- (0.9,-0.18);
\node[below] at (0,-1.55) {\textbf{Cellule mère en G1} : $2n=4$, 1 chromatide, $Q$};
% flèche réplication
\draw[-{Stealth[length=3mm]}, very thick, popgreen] (2.4,0) -- (3.6,0) node[midway, above, align=center]{Phase S\\ (réplication)};
% ---- Après S / métaphase ----
\draw[popblue, thick] (6.2,0) ellipse (2.1 and 1.4);
\draw[popdark, thick] (6.2,0) ellipse (1.5 and 0.95);
% grands chromosomes X
\draw[line width=2.2pt, poporange] (5.35,0.4) -- (5.55,-0.4);
\draw[line width=2.2pt, poporange] (5.55,0.4) -- (5.35,-0.4);
\draw[line width=2.2pt, poppink] (5.95,0.4) -- (6.15,-0.4);
\draw[line width=2.2pt, poppink] (6.15,0.4) -- (5.95,-0.4);
% petits chromosomes X
\draw[line width=2.2pt, poporange] (6.6,0.22) -- (6.75,-0.22);
\draw[line width=2.2pt, poporange] (6.75,0.22) -- (6.6,-0.22);
\draw[line width=2.2pt, poppink] (7.05,0.22) -- (7.2,-0.22);
\draw[line width=2.2pt, poppink] (7.2,0.22) -- (7.05,-0.22);
\node[below] at (6.2,-1.55) {\textbf{Métaphase} : $2n=4$, 2 chromatides, $2Q$};
% flèche anaphase
\draw[-{Stealth[length=3mm]}, very thick, popgreen] (6.2,-1.9) -- (6.2,-2.9) node[midway, right, align=left]{Anaphase : séparation\\ des chromatides sœurs};
% ---- Anaphase ----
\draw[popblue, thick] (6.2,-4.6) ellipse (2.3 and 1.3);
% lot gauche
\draw[line width=2.2pt, poporange] (4.9,-4.3) -- (4.9,-4.9);
\draw[line width=2.2pt, poppink] (5.35,-4.3) -- (5.35,-4.9);
\draw[line width=2.2pt, poporange] (5.8,-4.45) -- (5.8,-4.75);
\draw[line width=2.2pt, poppink] (6.05,-4.45) -- (6.05,-4.75);
% lot droit
\draw[line width=2.2pt, poporange] (6.45,-4.45) -- (6.45,-4.75);
\draw[line width=2.2pt, poppink] (6.7,-4.45) -- (6.7,-4.75);
\draw[line width=2.2pt, poporange] (7.15,-4.3) -- (7.15,-4.9);
\draw[line width=2.2pt, poppink] (7.6,-4.3) -- (7.6,-4.9);
\draw[-{Stealth}, popmuted] (5.4,-5.1) -- (4.6,-5.6);
\draw[-{Stealth}, popmuted] (7.1,-5.1) -- (7.9,-5.6);
\node[below] at (6.2,-6.1) {\textbf{Anaphase} : 8 chromosomes à 1 chromatide en migration, $2Q$};
% flèches vers cellules filles
\draw[-{Stealth[length=3mm]}, very thick, popgreen] (4.9,-6.4) -- (2.6,-7.4);
\draw[-{Stealth[length=3mm]}, very thick, popgreen] (7.5,-6.4) -- (9.8,-7.4);
% ---- Cellules filles ----
\draw[popblue, thick] (1.6,-8.6) ellipse (1.7 and 1.15);
\draw[popdark, thick] (1.6,-8.6) ellipse (1.2 and 0.75);
\draw[line width=2.2pt, poporange] (0.95,-8.35) -- (0.95,-8.85);
\draw[line width=2.2pt, poppink] (1.35,-8.35) -- (1.35,-8.85);
\draw[line width=2.2pt, poporange] (1.85,-8.45) -- (1.85,-8.75);
\draw[line width=2.2pt, poppink] (2.2,-8.45) -- (2.2,-8.75);
\node[below, align=center] at (1.6,-9.9) {\textbf{Cellule fille 1}\\ $2n=4$, 1 chromatide, $Q$};
\draw[popblue, thick] (10.8,-8.6) ellipse (1.7 and 1.15);
\draw[popdark, thick] (10.8,-8.6) ellipse (1.2 and 0.75);
\draw[line width=2.2pt, poporange] (10.15,-8.35) -- (10.15,-8.85);
\draw[line width=2.2pt, poppink] (10.55,-8.35) -- (10.55,-8.85);
\draw[line width=2.2pt, poporange] (11.05,-8.45) -- (11.05,-8.75);
\draw[line width=2.2pt, poppink] (11.4,-8.45) -- (11.4,-8.75);
\node[below, align=center] at (10.8,-9.9) {\textbf{Cellule fille 2}\\ $2n=4$, 1 chromatide, $Q$};
\end{tikzpicture}
\end{center}
\legende{Devenir des chromosomes d'une cellule mère à $2n = 4$ au cours de la mitose. En orange et en rose : les deux chromosomes homologues de chaque paire (grand et petit). Après réplication, chaque chromosome comporte deux chromatides sœurs (forme en X). En anaphase, les chromatides sœurs se séparent : chaque pôle reçoit un exemplaire de chaque chromosome. Les deux cellules filles obtiennent chacune $2n = 4$ chromosomes à une chromatide, identiques à ceux de la cellule mère.}
\end{popfigure}

Raisonnons étape par étape :
\begin{enumerate}
\item \textbf{En G1}, la cellule possède 4 chromosomes simples (1 chromatide chacun) : 2 grands et 2 petits.
\item \textbf{Après la phase S}, chacun des 4 chromosomes s'est répliqué en 2 chromatides sœurs identiques, réunies par un centromère. Il y a toujours 4 centromères, donc \textbf{4 chromosomes}, mais 8 molécules d'ADN.
\item \textbf{En anaphase}, les 4 centromères se divisent ; les 8 chromatides migrent par groupes de 4 vers les deux pôles. Chaque lot polaire contient obligatoirement 1 grand chromosome orange, 1 grand rose, 1 petit orange, 1 petit rose : c'est la disposition métaphasique sur la plaque équatoriale qui garantit cette régularité.
\item \textbf{Après la cytodiérèse}, chaque cellule fille possède $2n = 4$ chromosomes à 1 chromatide, de mêmes tailles et de mêmes formes que ceux de la cellule mère : le \textbf{caryotype est conservé}.
\end{enumerate}

\begin{exemplebox}[Le cas de la cellule humaine : $2n = 46$]
Chez l'Homme, $2n = 46$ chromosomes (23 paires). Une cellule humaine en \textbf{métaphase} contient 46 chromosomes, chacun formé de 2 chromatides sœurs, soit \textbf{92 molécules d'ADN} (quantité $2Q$). En anaphase, les 92 chromatides se séparent : 46 migrent vers chaque pôle. Après la mitose, on obtient \textbf{deux cellules filles à 46 chromosomes à 1 chromatide} (46 molécules d'ADN, quantité $Q$), génétiquement identiques à la cellule mère. C'est ainsi que les milliards de cellules de ton organisme, toutes issues d'une unique cellule-œuf, portent le même caryotype.
\end{exemplebox}

\courssub{Conservation du caryotype et rôles biologiques de la mitose}

La conservation du caryotype n'est pas une simple curiosité de laboratoire : elle fonde les grands rôles biologiques de la mitose.

\begin{aretenir}[Les rôles biologiques de la mitose]
\begin{itemize}
\item \textbf{La croissance} : un organisme pluricellulaire grandit par augmentation du nombre de ses cellules, toutes issues de mitoses successives à partir de la cellule-œuf. Chaque cellule nouvelle reçoit le patrimoine génétique complet.
\item \textbf{Le renouvellement et la réparation des tissus} : les cellules usées ou détruites (peau, sang, muqueuses) sont remplacées par des cellules identiques grâce aux mitoses des cellules souches.
\item \textbf{La reproduction asexuée} : chez de nombreux végétaux, un fragment de l'organisme parent régénère un individu entier par mitoses. C'est le principe du \textbf{bouturage du manioc}, largement pratiqué par les agriculteurs camerounais : une tige de manioc mise en terre produit un plant génétiquement identique au plant mère, ce qui conserve les qualités de la variété (rendement, goût, résistance). Le même principe s'applique aux boutures de plantes ornementales (hibiscus, rosier).
\end{itemize}
\end{aretenir}

\begin{savaistu}[Un épiderme neuf chaque mois]
L'épiderme humain se renouvelle entièrement en environ \textbf{28 jours} grâce aux mitoses incessantes des cellules de la couche basale, la couche la plus profonde de l'épiderme. Chaque jour, des milliards de cellules superficielles mortes desquament (c'est la fine poussière blanche sur la peau sèche) et sont remplacées par des cellules filles identiques. Sans la fidélité de la mitose, ce renouvellement permanent produirait des cellules défectueuses et la peau perdrait son rôle de barrière protectrice.
\end{savaistu}

\begin{attention}[La mitose ne réduit pas le nombre de chromosomes]
Erreur classique au Probatoire : écrire que la mitose « divise par deux » le nombre de chromosomes. C'est faux ! La mitose \textbf{conserve} le nombre de chromosomes : une cellule mère $2n$ donne deux cellules filles $2n$. Le doublement transitoire observé en anaphase (passage de $2n$ à $4n$ chromosomes à 1 chromatide dans la cellule entière) est immédiatement suivi du partage en deux cellules. La division du stock de chromosomes par deux, qui fait passer de $2n$ à $n$, est le propre de la \textbf{méiose}, que nous étudierons ultérieurement. Ne confonds jamais ces deux mécanismes.
\end{attention}

\courssub{Exercice résolu et conclusion}

\begin{exoresolu}[Suivi quantitatif chez une cellule animale]
\textbf{Énoncé.} Une cellule animale possède $2n = 8$ chromosomes et une quantité d'ADN égale à $Q = \SI{6}{\pico\gram}$ en phase G1. (1) Donner le nombre de chromosomes, le nombre de chromatides et la masse d'ADN de cette cellule en métaphase, puis en anaphase. (2) Mêmes questions pour chaque cellule fille issue de la mitose. (3) Expliquer pourquoi ce résultat illustre la conservation du caryotype.
\tcblower
\textbf{Solution détaillée.}
\begin{enumerate}
\item \textbf{En métaphase} : la phase S a eu lieu, chaque chromosome est double. On compte les centromères : \textbf{8 chromosomes}, chacun à 2 chromatides, soit \textbf{16 chromatides}. La masse d'ADN a doublé : $2Q = \SI{12}{\pico\gram}$. \textbf{En anaphase} : les centromères se sont divisés, chaque chromatide devient un chromosome : \textbf{16 chromosomes à 1 chromatide} (16 chromatides) dans la cellule entière ; la masse d'ADN reste $2Q = \SI{12}{\pico\gram}$, car aucune synthèse ni perte d'ADN n'a eu lieu.
\item \textbf{Chaque cellule fille} reçoit la moitié du lot anaphasique : \textbf{8 chromosomes à 1 chromatide} (8 chromatides) et une masse d'ADN $Q = \SI{6}{\pico\gram}$.
\item La cellule mère en G1 et chaque cellule fille présentent le même nombre de chromosomes ($2n = 8$), la même quantité d'ADN ($Q$) et, puisque les chromatides sœurs étaient des copies identiques issues de la réplication, la même information génétique : le \textbf{caryotype est intégralement conservé} d'une génération cellulaire à l'autre.
\end{enumerate}
\end{exoresolu}

En définitive, tout se résume en une formule que tu dois retenir : \textbf{réplication + équipartition = stabilité}. La réplication de l'ADN en phase S fabrique deux copies fidèles de chaque chromosome ; la mitose, par la mise en place métaphasique puis la séparation des chromatides sœurs en anaphase, distribue une copie de chaque chromosome à chaque cellule fille. Le caryotype — nombre et morphologie des chromosomes — est ainsi transmis sans altération. C'est cette fidélité qui permet la croissance d'un organisme, la cicatrisation d'une plaie, le renouvellement de l'épiderme et le bouturage du manioc, en garantissant le \textbf{maintien de l'identité biologique} des cellules et des organismes. Nous verrons dans la section suivante ce qui se produit lorsque ce mécanisme si précis se dérègle : la cancérisation.

\coursec{Dysfonctionnements de la mitose et cancérisation}

Nous venons d'étudier comment la mitose, précédée d'une réplication fidèle de l'ADN, assure la transmission intégrale du patrimoine génétique d'une cellule mère à ses deux filles. Ce processus est d'une précision remarquable : chaque chromosome est dupliqué, aligné, puis séparé avec une exactitude qui garantit la stabilité du caryotype. Pourtant, dans un organisme vivant, des milliards de divisions ont lieu chaque jour. Que se passe-t-il lorsque la « machine » se dérègle ? Observons une situation de santé publique majeure au Cameroun.

\courssub{Du contrôle physiologique à la prolifération anarchique}

\begin{savaistu}[Le cancer au Cameroun : un défi de santé publique]
Au Cameroun, selon les données du Registre des Cancers de Yaoundé et de l'Organisation Mondiale de la Santé, les cancers les plus fréquents sont le cancer du col de l'utérus (lié au papillomavirus humain - HPV), le cancer du sein, le cancer du foie (lié aux hépatites B/C et aux aflatoxines) et le cancer de la prostate. Le tabagisme, l'alcoolisation, l'exposition aux pesticides dans l'agriculture et les infections chroniques sont des facteurs de risque majeurs. La prise en charge reste difficile du fait du coût des traitements (chimiothérapie, radiothérapie au Centre Hospitalier Universitaire de Yaoundé ou au Centre de Cancérologie de Douala) et du diagnostic souvent tardif. La prévention (vaccination HPV, dépistage, lutte antitabac) est donc une priorité nationale.
\end{savaistu}

Dans un tissu sain, la division cellulaire n'est pas un processus continu et incontrôlé. Elle répond à des besoins précis : croissance de l'organisme, renouvellement des cellules usées (épiderme, muqueuse intestinale, globules rouges), réparation d'une blessure. Ce contrôle est assuré par des \cle{points de contrôle} (ou \emph{checkpoints}) au cours du cycle cellulaire (notamment en $G_1$, $G_2$ et $M$) et par des signaux extracellulaires (\cle{facteurs de croissance}, hormones, contact entre cellules).

\begin{experience}[Observation microscopique : tissu sain vs tissu tumoral]
\textbf{Matériel :} Lames minces de peau saine (épiderme) et de tumeur maligne (ex. carcinome épidermoïde), coloration Hématoxyline-Éosine (HE), microscope optique.
\textbf{Protocole :} Observer à fort grossissement (×400) la morphologie des noyaux, le rapport noyau/cytoplasme, la présence de mitoses, l'organisation tissulaire.
\textbf{Observations :}
\begin{itemize}
    \item \textbf{Tissu sain :} Cellules organisées en couches régulières. Noyaux uniformes, peu de figures de mitose visibles (seules les cellules basales se divisent). Respect de la polarité cellulaire.
    \item \textbf{Tissu tumoral :} Architecture désorganisée, perte de la polarité. Noyaux volumineux, pléomorphes (formes et tailles variables), chromatine dense. Nombreuses figures de mitose, souvent \textbf{anormales} (fuseau multipolaire, chromosomes retardataires). Rapport noyau/cytoplasme augmenté.
\end{itemize}
\textbf{Interprétation :} La perte de l'organisation tissulaire et la présence de mitoses atypiques traduisent un échappement aux contrôles normaux du cycle cellulaire. Les cellules prolifèrent sans respecter les signaux d'arrêt (inhibition par contact) ni les points de contrôle de l'ADN.
\textbf{Conclusion :} La cancérisation se manifeste au niveau cellulaire par une perte de la régulation du cycle cellulaire entraînant une prolifération anarchique et la production de cellules aux caractéristiques morphologiques anormales.
\end{experience}

\begin{definition}[Cancer et cellule cancéreuse]
Un \textbf{cancer} est une maladie caractérisée par la \textbf{prolifération incontrôlée} de cellules dont le cycle cellulaire a échappé aux mécanismes de régulation physiologique. Une \textbf{cellule cancéreuse} est une cellule qui a acquis la capacité de se diviser de manière illimitée et autonome, indépendamment des signaux de croissance ou d'arrêt normaux, et qui peut envahir les tissus voisins et coloniser des organes distants (métastases).
\end{definition}

\courssub{L'origine moléculaire : l'accumulation de mutations}

Comment une cellule normale devient-elle cancéreuse ? Ce n'est pas un événement unique, mais un processus évolutif à l'échelle cellulaire : la \cle{cancérisation}. Elle résulte de l'accumulation successive de \cle{mutations} dans des gènes clés qui contrôlent le cycle cellulaire, la réparation de l'ADN ou la mort cellulaire programmée (apoptose).

\begin{propriete}[Base génétique de la cancérisation - Admise]
La transformation d'une cellule saine en cellule cancéreuse nécessite l'accumulation de plusieurs mutations (généralement 3 à 7) affectant deux grandes classes de gènes :
\begin{itemize}
    \item Les \textbf{proto-oncogènes} : gènes normaux favorisant la division (codant pour facteurs de croissance, récepteurs, protéines de transduction du signal, cyclines). Leur mutation (activation constitutive) les transforme en \textbf{oncogènes} $\rightarrow$ « accélérateur bloqué ».
    \item Les \textbf{gènes suppresseurs de tumeurs} (ex. \emph{TP53}, \emph{Rb}) : gènes freinant la division ou induisant l'apoptose en cas d'anomalie. Leur inactivation (perte de fonction des deux allèles) $\rightarrow$ « freins coupés ».
\end{itemize}
L'instabilité génomique (défaut de réparation de l'ADN, dysfonctionnement de la mitose) accélère cette accumulation.
\end{propriete}

\begin{attention}[Piège : « Un gène = un cancer » ?]
Non. Un cancer résulte \textbf{toujours} de l'accumulation de plusieurs mutations dans une même lignée cellulaire. C'est pourquoi le cancer est une maladie de l'âge (le temps permet l'accumulation) et pourquoi l'exposition prolongée à des mutagènes augmente le risque. Une seule mutation dans un oncogène ou un gène suppresseur ne suffit pas à déclencher un cancer.
\end{attention}

\courssub{De la tumeur aux métastases : l'évolution clinique}

La prolifération anarchique d'une cellule transformée donne naissance à un clone de cellules identiques : la \cle{tumeur} (ou néoplasme). On distingue deux stades majeurs :

\begin{center}
\begin{tabularx}{\linewidth}{|l|X|X|}
\hline
\rowcolor{popblueL} \textbf{Critère} & \textbf{Tumeur bénigne} & \textbf{Tumeur maligne (Cancer)} \\
\hline
Croissance & Lente, expansive, encapsulée & Rapide, infiltrative, non encapsulée \\
\hline
Limites & Nettes, repoussant les tissus sains & Floues, envahissant les tissus voisins \\
\hline
Métastases & \textbf{Absentes} & \textbf{Présentes} (voie lymphatique, sanguine) \\
\hline
Cellularité & Proche du tissu d'origine (bien différenciée) & Anormale (mal différenciée, atypies nucléaires, mitoses bizarres) \\
\hline
Risque vital & Généralement faible (sauf compression organe vital) & Élevé (envahissement, cachexie, défaillance d'organes) \\
\hline
Exemple & Adénome, fibrome, lipome & Carcinome (épithélium), sarcome (conjonctif), leucémie, lymphome \\
\hline
\end{tabularx}
\end{center}

\begin{aretenir}[Notion qualitative de métastase]
Une \textbf{métastase} est une tumeur secondaire, située à distance de la tumeur primitive, issue de cellules cancéreuses qui se sont détachées, ont migré (via la lymphe ou le sang), ont envahi un nouvel organe et s'y sont multipliées. C'est la métastase, et non la tumeur primitive, qui est la cause principale de décès par cancer. Seule une tumeur \textbf{maligne} métastase.
\end{aretenir}

\begin{popfigure}
\begin{tikzpicture}[scale=0.9, every node/.style={font=\small}]
% Tissu sain (gauche)
\begin{scope}[xshift=-4cm]
\draw[popblue, thick, fill=popblue!5] (-3,-2.5) rectangle (3,2.5);
\draw[popdark, thick] (-3,0) -- (3,0) node[midway, above, sloped] {Membrane basale};
% Couche basale : cellules cubiques, noyaux ronds
\foreach \x in {-2.5,-1.5,-0.5,0.5,1.5,2.5} {
    \draw[popblue, fill=popblue!20] (\x,-0.2) rectangle ++(0.8,0.4);
    \draw[popdark, fill=popdark] (\x+0.1,0) circle (0.12); % Noyau
}
% Couches supérieures : cellules aplaties, noyaux petits
\foreach \x in {-2.5,-1.5,-0.5,0.5,1.5,2.5} {
    \draw[popblue, fill=popblue!10] (\x,0.3) rectangle ++(0.8,0.6);
    \draw[popdark, fill=popdark] (\x+0.1,0.6) circle (0.06);
}
\foreach \x in {-2.5,-1.5,-0.5,0.5,1.5,2.5} {
    \draw[popblue, fill=popblue!10] (\x,1.0) rectangle ++(0.8,0.5);
    \draw[popdark, fill=popdark] (\x+0.1,1.2) circle (0.04);
}
% Une mitose normale en bas
\draw[popgreen, fill=popgreen!20] (-0.5,-0.8) rectangle ++(0.8,0.4);
\draw[popgreen, thick] (-0.1,-0.4) -- (0.3,-0.4); % Plan équatorial
\draw[popgreen, thick] (-0.1,-0.6) -- (-0.1,-0.2); % Fuseau
\node[popdark, align=center] at (0,-1.5) {Tissu sain :\\ organisation polarisée,\\ mitoses rares et normales};
\end{scope}

% Tumeur (droite)
\begin{scope}[xshift=4cm]
\draw[poporange, thick, fill=poporange!5] (-3,-2.5) rectangle (3,2.5);
% Membrane basale rompue
\draw[poporange, thick, dashed] (-3,0) -- (-1,0) -- (-0.5,0.3) -- (0.5,-0.3) -- (1,0) -- (3,0) node[midway, above, sloped] {Membrane basale rompue};
% Cellules désorganisées, noyaux gros, mitoses multiples
\foreach \i in {1,...,30} {
    \pgfmathsetmacro{\rx}{rand*2.5}
    \pgfmathsetmacro{\ry}{rand*2.0}
    \pgfmathsetmacro{\rs}{0.3+rand*0.3}
    \pgfmathsetmacro{\rn}{0.08+rand*0.06}
    \draw[poporange, fill=poporange!20] (\rx,\ry) circle (\rs);
    \draw[popdark, fill=popdark] (\rx,\ry) circle (\rn);
    % Quelques mitoses anormales (multipolaires)
    \ifnum\i<4
        \draw[poppink, thick] (\rx,\ry) -- ++(0.3,0.3) (\rx,\ry) -- ++(-0.3,0.3) (\rx,\ry) -- ++(0.3,-0.3) (\rx,\ry) -- ++(-0.3,-0.3);
    \fi
}
% Cellules dans le tissu conjonctif (invasion)
\foreach \i in {1,...,10} {
    \pgfmathsetmacro{\rx}{-2.5+rand*5}
    \pgfmathsetmacro{\ry}{-2.2+rand*0.5}
    \pgfmathsetmacro{\rs}{0.2+rand*0.15}
    \pgfmathsetmacro{\rn}{0.06+rand*0.04}
    \draw[poporange, fill=poporange!30] (\rx,\ry) circle (\rs);
    \draw[popdark, fill=popdark] (\rx,\ry) circle (\rn);
}
\node[popdark, align=center] at (0,-1.5) {Tumeur maligne :\\ architecture perdue,\\ noyaux pléomorphes,\\ mitoses anormales,\\ invasion};
\end{scope}
\end{tikzpicture}
\legende{Comparaison schématique d'un épithélium sain (gauche) et d'un carcinome épidermoïde (droite). En haut : membrane basale intacte vs rompue. Au centre : cellules polarisées, noyaux uniformes, mitose rare et bipolaire vs cellules désorganisées, noyaux volumineux et pléomorphes (pléomorphisme), mitoses multipolaires anormales. En bas : respect de la frontière vs invasion du tissu conjonctif sous-jacent.}
\end{popfigure}

\courssub{Facteurs favorisants : les mutagènes cancérogènes}

Les agents qui augmentent la fréquence des mutations (mutagènes) sont des \cle{cancérogènes}. Au Cameroun, l'exposition est multiple :

\begin{itemize}
    \item \textbf{Rayonnements ionisants} (UV solaires intenses au Nord, radiographies répétées sans protection, radon dans certaines zones granitiques) et \textbf{non ionisants} (UV solaires $\rightarrow$ mélanomes, carcinomes cutanés).
    \item \textbf{Substances chimiques} : \textbf{Goudrons du tabac} (benzopyrène $\rightarrow$ cancers du poumon, larynx, vessie, œsophage) ; \textbf{Aflatoxines} (toxines de moisissures \emph{Aspergillus} sur maïs, arachides mal séchés $\rightarrow$ cancer du foie, synergie avec hépatite B) ; \textbf{Pesticides} (agriculteurs non protégés) ; \textbf{Amiante} (bâtiments anciens).
    \item \textbf{Agent biologiques} : Virus \textbf{HPV} (types 16, 18 $\rightarrow$ cancer col utérus, anus, oropharynx) ; Virus \textbf{HBV/HCV} (hépatites B/C $\rightarrow$ cirrhose puis hépatocarcinome) ; \emph{Helicobacter pylori} (cancer gastrique) ; Virus \textbf{EBV} (lymphome de Burkitt, zone équatoriale humide) ; Schistosomes (cancer vessie).
    \item \textbf{Prédispositions génétiques} : Mutations constitutionnelles (germinales) de \emph{BRCA1/2} (sein, ovaire), \emph{TP53} (syndrome de Li-Fraumeni), \emph{APC} (polypose adénomateuse familiale). Représentent 5 à 10\% des cancers.
\end{itemize}

\begin{savaistu}[Le cancer du col de l'utérus : une tragédie évitable au Cameroun]
C'est le 1er cancer féminin au Cameroun (environ 2000 nouveaux cas/an, $>$1000 décès/an). Causé par une infection persistante à HPV à haut risque (transmission sexuelle). La vaccination des jeunes filles (9-14 ans) avant le début de l'activité sexuelle, introduite dans le Programme Élargi de Vaccination (PEV) depuis 2020, et le dépistage par test HPV ou inspection visuelle (IVA/VILI) chez les femmes de 30-49 ans sont les piliers de l'élimination visée par l'OMS pour 2030. Le dépistage précoce permet de traiter les lésions précancéreuses (CIN) par cryothérapie ou LEEP, évitant le cancer invasif.
\end{savaistu}

\courssub{Principes thérapeutiques : cibler la division rapide}

Puisque la caractéristique fondamentale de la cellule cancéreuse est sa division rapide et incontrôlée, les traitements conventionnels visent à bloquer le cycle cellulaire ou à tuer les cellules en division.

\begin{methode}[Principes de la chimiothérapie et radiothérapie]
\begin{enumerate}
    \item \textbf{Chimiothérapie (systémique) :} Administration de molécules \textbf{cytotoxiques} (par voie veineuse ou orale) qui circulent dans tout l'organisme.
        \begin{itemize}
            \item \textbf{Agents alkylants} (Cyclophosphamide) : créent des ponts entre brins d'ADN $\rightarrow$ blocage réplication/mitose.
            \item \textbf{Antimétabolites} (5-Fluorouracile, Méthotrexate) : miment les bases azotées ou l'acide folique $\rightarrow$ inhibent la synthèse d'ADN (phase S).
            \item \textbf{Alcaloïdes de la pervenche de Madagascar (\emph{Catharanthus roseus})} : Vincristine, Vinblastine. Bloquent la polymérisation de la tubuline $\rightarrow$ \textbf{blocage du fuseau mitotique en métaphase}.
            \item \textbf{Taxanes} (Paclitaxel) : stabilisent les microtubules $\rightarrow$ empêchent la dépolymérisation nécessaire à l'anaphase.
            \item \textbf{Inhibiteurs de topoisomérases} (Doxorubicine, Étoposide) : bloquent le dénouement de l'ADN $\rightarrow$ cassures.
        \end{itemize}
    \item \textbf{Radiothérapie (locale) :} Faisceaux de rayons X (photons haute énergie) ou d'électrons dirigés vers la tumeur. L'énergie ionise l'eau cellulaire $\rightarrow$ radicaux libres $\rightarrow$ \textbf{cassures double brin de l'ADN}. Les cellules en mitose (chromosomes condensés) et en phases $G_2/M$ sont les plus radiosensibles.
    \item \textbf{Thérapies ciblées / Immunothérapie (modernes) :} Inhibiteurs de tyrosine kinases (ex. Imatinib pour leucémie myéloïde chronique), anticorps monoclonaux (anti-HER2, anti-VEGF, anti-PD1/PD-L1). Plus sélectifs, moins d'effets sur cellules saines.
\end{enumerate}
\end{methode}

\begin{attention}[Effets secondaires : le prix de la non-sélectivité]
La chimiothérapie et la radiothérapie classiques touchent \textbf{toutes} les cellules à renouvellement rapide : \textbf{moelle osseuse} (aplasie $\rightarrow$ anémie, infections, saignements), \textbf{muqueuse digestive} (nausées, vomissements, diarrhées, mucites), \textbf{follicules pileux} (alopécie/chute des cheveux), \textbf{gamétogenèse} (infertilité transitoire ou définitive). C'est la limite majeure de ces traitements « cytotoxiques » non spécifiques. La prise en charge des effets secondaires (soins de support : antiémétiques, facteurs de croissance hématopoïétiques, nutrition, psychologie) est indissociable du traitement anticancéreux au Cameroun.
\end{attention}

\courssub{Interprétation de données expérimentales : courbes de prolifération}

\begin{methode}[Analyser une courbe de croissance cellulaire]
\begin{enumerate}
    \item \textbf{Identifier les axes :} Abscisse = Temps (heures, jours) ; Ordonnée = Nombre de cellules (échelle linéaire ou logarithmique).
    \item \textbf{Repérer les phases :} Latence ($G_1$ long), croissance exponentielle (pente constante en échelle log), plateau (confluence, manque nutriments, inhibition par contact).
    \item \textbf{Comparer les pentes :} Pente plus raide = temps de doublement plus court = prolifération plus rapide.
    \item \textbf{Détecter l'anomalie :} Absence de plateau (croissance illimitée), reprise de croissance après traitement, courbe ne passant pas par une phase de latence.
    \item \textbf{Conclure :} Relier la forme de la courbe au statut du cycle cellulaire (contrôlé vs dérégulé).
\end{enumerate}
\end{methode}

\begin{exoresolu}[Analyse comparative de cinétiques de croissance]
\textbf{Énoncé :} On réalise une mise en culture de deux lignées cellulaires humaines : des fibroblastes normaux (lignée WI-38) et des cellules de carcinome du col utérin (lignée HeLa). On compte le nombre de cellules vivantes chaque jour pendant 10 jours. Résultats (échelle logarithmique) :
\begin{center}
\begin{tabular}{|c|c|c|}\hline
\rowcolor{popblueL} \textbf{Jour} & \textbf{WI-38 (normales)} & \textbf{HeLa (cancéreuses)} \\ \hline
0 & $1 \times 10^4$ & $1 \times 10^4$ \\ \hline
2 & $2 \times 10^4$ & $4 \times 10^4$ \\ \hline
4 & $4 \times 10^4$ & $16 \times 10^4$ \\ \hline
6 & $8 \times 10^4$ & $64 \times 10^4$ \\ \hline
8 & $8 \times 10^4$ (plateau) & $256 \times 10^4$ \\ \hline
10 & $7 \times 10^4$ (décès) & $1000 \times 10^4$ \\ \hline
\end{tabular}
\end{center}
\textbf{Questions :}
\begin{enumerate}
    \item Représenter graphiquement les deux courbes en échelle semi-logarithmique (linéaire en abscisse, log en ordonnée).
    \item Calculer le temps de doublement (Td) de chaque lignée pendant la phase exponentielle.
    \item Interpréter les différences observées en lien avec le cycle cellulaire et la cancérisation.
\end{enumerate}

\textbf{Solution détaillée :}
\begin{enumerate}
    \item \textbf{Graphique :} Voir figure ci-dessous. En échelle semi-log, une croissance exponentielle donne une \textbf{droite}. La pente de cette droite est proportionnelle à la vitesse de division.
    \item \textbf{Calcul du temps de doublement (Td) :}
        \begin{itemize}
            \item \textbf{WI-38 :} De J0 à J6, le nombre passe de $10^4$ à $8 \times 10^4$ (3 doubles : $1 \to 2 \to 4 \to 8$). 3 doubles en 6 jours $\rightarrow$ Td = 2 jours. \textbf{Mais} plateau atteint dès J6-J8 (inhibition par contact, sénescence réplicative : limite de Hayflick $\approx$ 50 divisions).
            \item \textbf{HeLa :} De J0 à J10, passage de $10^4$ à $10^7$ ($1000 \times$). $\log_2(1000) \approx 9,97 \approx 10$ doubles en 10 jours $\rightarrow$ \textbf{Td $\approx$ 1 jour}. Pas de plateau observé sur la période (immortalité cellulaire : expression de la télomérase, inactivation de p53/Rb).
        \end{itemize}
    \item \textbf{Interprétation :}
        \begin{itemize}
            \item La lignée \textbf{HeLa croît 2 fois plus vite} (Td = 1 j vs 2 j) et \textbf{indéfiniment} (pas de plateau), traduisant la perte des points de contrôle ($G_1/S$ : Rb, p53 ; sénescence : télomérase) et l'autonomie vis-à-vis des facteurs de croissance.
            \item La lignée \textbf{WI-38} respecte l'inhibition par contact (arrêt en $G_1$ à confluence) et subit la sénescence réplicative (raccourcissement télomères $\rightarrow$ arrêt définitif).
            \item Cela illustre la \textbf{propriété} : mutations accumulées $\rightarrow$ oncogènes activés (cycline D, Myc) + suppresseurs inactivés (p53, Rb) + télomérase réactivée $\rightarrow$ prolifération autonome, illimitée, immortelle = phénotype cancéreux.
        \end{itemize}
\end{enumerate}
\end{exoresolu}

\begin{popfigure}
\begin{tikzpicture}
\begin{semilogyaxis}[
    width=\linewidth, height=10cm,
    xlabel={Temps (jours)}, ylabel={Nombre de cellules (échelle log)},
    xmin=0, xmax=10.5, ymin=5e3, ymax=2e7,
    grid=major, grid style={popline},
    legend pos=north west,
    legend cell align=left,
    tick label style={font=\small},
    label style={font=\small},
    xtick={0,2,4,6,8,10},
    ytick={1e4,1e5,1e6,1e7},
    yticklabels={$10^4$,$10^5$,$10^6$,$10^7$}
]
\addplot[popcurve, popblue, very thick, mark=square*, mark size=3] coordinates {
    (0,10000) (2,20000) (4,40000) (6,80000) (8,80000) (10,70000)
};
\addlegendentry{Fibroblastes normaux (WI-38) : plateau, sénescence}

\addplot[popcurve, poporange, very thick, mark=triangle*, mark size=3] coordinates {
    (0,10000) (2,40000) (4,160000) (6,640000) (8,2560000) (10,10000000)
};
\addlegendentry{Cellules HeLa (cancer col utérus) : exponentielle continue, immortalité}
\end{semilogyaxis}
\end{tikzpicture}
\legende{Cinétiques de croissance en culture : comparaison entre cellules normales (WI-38, bleu) et cellules cancéreuses (HeLa, orange) en échelle semi-logarithmique. Les cellules normales montrent une phase exponentielle (pente constante) suivie d'un plateau (arrêt de la division par inhibition de contact/sénescence). Les cellules HeLa conservent une croissance exponentielle ininterrompue (pente plus raide = temps de doublement plus court) sans plateau, illustrant l'immortalité cellulaire et l'échappement aux contrôles du cycle.}
\end{popfigure}

\courssub{Prévention et sensibilisation : agir en amont}

La lutte contre le cancer ne se limite pas au traitement. La prévention primaire (éviter l'apparition) et secondaire (dépister tôt) sont les plus efficaces et rentables, surtout dans un contexte de ressources limitées comme au Cameroun.

\begin{aretenir}[Piliers de la prévention des cancers au Cameroun]
\begin{enumerate}
    \item \textbf{Lutte contre le tabagisme :} Application stricte de la loi 2006/018 (interdiction fumer lieux publics, paquet neutre, taxation). Éducation jeunesse. \textit{1er facteur évitable de cancer (poumon, vessie, ORL, pancréas...)}.
    \item \textbf{Vaccination :} \textbf{HPV} (filles 9-14 ans, 2 doses) $\rightarrow$ prévention 70-90\% cancers col utérus + vulve, vagin, anus, oropharynx. \textbf{Hépatite B} (à la naissance + nourrisson) $\rightarrow$ prévention cirrhose/cancer foie. Intégrées au PEV (gratuites).
    \item \textbf{Dépistage organisé :} Col utérin (test HPV / IVA-VILI, femmes 30-49 ans, tous les 3-5 ans). Sein (palpation clinique, mammographie si $>$40 ans ou risque). Prostate (PSA + toucher rectal, hommes $>$50 ans). \textbf{Diagnostic précoce = guérison possible, traitement moins lourd, moins cher.}
    \item \textbf{Hygiène de vie :} Alimentation variée (fruits/légumes = antioxydants, fibres), limitation alcool, activité physique, protection solaire (albinos particulièrement à risque), sécurité au travail (EPI pesticides, amiante).
    \item \textbf{Lutte contre les infections :} Traitement \emph{H. pylori}, schistosomiase (praziquantel), VIH (ARV restaurent immunité $\rightarrow$ moins cancers associés : Kaposi, lymphomes, col utérin).
\end{enumerate}
\end{aretenir}

\begin{savaistu}[Les cellules HeLa : un héritage éthique et scientifique]
La lignée HeLa utilisée dans l'exercice résolu provient d'une biopsie de tumeur du col de l'utérus prélevée en 1951 sur \textbf{Henrietta Lacks}, une jeune mère afro-américaine, \textbf{sans son consentement ni celui de sa famille}. Ses cellules furent les premières cellules humaines « immortelles » en culture (grâce à l'HPV 18 qui inactive p53/Rb et active la télomérase). Elles ont permis le vaccin antipoliomyélitique, la découverte des télomères, la FIV, la cartographie du génome, les tests de médicaments... Des milliards de cellules HeLa ont été produites, envoyées dans l'espace, utilisées dans 100 000+ publications. La famille Lacks n'a été informée qu'en 1973 et n'a jamais bénéficié financièrement. Depuis 2013, un accord NIH assure un contrôle partagé de l'accès aux données génomiques HeLa. Cette histoire illustre l'importance cruciale de l'\textbf{éthique de la recherche} (consentement éclairé, partage des bénéfices) et la puissance de la prolifération cellulaire dérégulée.
\end{savaistu}

\begin{exercice}[Cas clinique : Adama, 45 ans, cultivateur à Dschang]
Adama consulte pour une masse testiculaire indolore découverte depuis 3 mois, augmentant de volume. Échographie : tumeur testiculaire solide. Marqueurs sanguins : $\alpha$-foetoprotéine (AFP) = 450 ng/mL (N $<$ 10), $\beta$-HCG = 120 UI/L (N $<$ 5), LDH = 450 UI/L (N $<$ 250). Biopsie (orchidectomie) : tumeur maligne à cellules germinales (non-séminome : embryonnaire + choriocarcinome). Scanner thoraco-abdomino-pelvien : ganglions rétropéritonéaux multiples (max 4 cm), nodules pulmonaires bilatéraux.
\begin{enumerate}
    \item Classer cette tumeur (bénigne/maligne) et justifier par trois arguments histologiques/cliniques.
    \item Expliquer l'intérêt des marqueurs AFP et $\beta$-HCG pour le diagnostic et le suivi.
    \item Proposer le traitement standard de première ligne (chimiothérapie) en citant un protocole type et son mécanisme d'action sur le cycle cellulaire.
    \item Citer deux effets secondaires majeurs attendus et leur prise en charge.
    \item Quel message de prévention adresser à la communauté de Dschang concernant les cancers testiculaires ?
\end{enumerate}
\end{exercice}

\begin{corrige}[Exercice : Cas clinique Adama]
\begin{enumerate}
    \item \textbf{Tumeur maligne (cancer).} Arguments : 1) Croissance progressive rapide (3 mois). 2) Marqueurs tumoraux élevés (AFP, $\beta$-HCG) sécrétés par la tumeur. 3) \textbf{Métastases} ganglionnaires (rétropéritoine) et pulmonaires $\rightarrow$ critère de malignité absolu. 4) Histologie : cellules embryonnaires/choriocarcinome = anaplasie, mitoses nombreuses.
    \item \textbf{AFP} (synthétisée par composante endodermique sinus vitellin/embryonnaire) et \textbf{$\beta$-HCG} (synthétisée par composante choriocarcinome = trophoblaste malin) sont des \textbf{marqueurs tumoraux spécifiques} des tumeurs germinales non-séminomateuses. Ils permettent : diagnostic positif, évaluation de l'extension (charge tumorale), \textbf{suivi de la réponse au traitement} (baisse exponentielle normale post-chimio : demi-vie AFP 5-7 j, $\beta$-HCG 1-3 j), détection précoce de rechute (remontée).
    \item \textbf{Chimiothérapie BEP (Bleomycine, Étoposide, Cisplatine) :} 3 à 4 cycles (21 jours).
        \begin{itemize}
            \item \textbf{Cisplatine :} Agent alkylant (ponts ADN intra/inter-brins) $\rightarrow$ blocage réplication/transcription $\rightarrow$ apoptose (phase $G_2/M$ et S).
            \item \textbf{Étoposide :} Inhibiteur topoisomérase II $\rightarrow$ cassures double brin ADN $\rightarrow$ arrêt $G_2/M$, apoptose.
            \item \textbf{Bleomycine :} Glycopeptide $\rightarrow$ radicaux libres $\rightarrow$ cassures ADN simple/double brin (spécifique phase $G_2/M$).
        \end{itemize}
        Taux de guérison $>$ 90\% même métastatique (tumeur très chimiosensible).
    \item \textbf{Effets secondaires majeurs :}
        \begin{itemize}
            \item \textbf{Aplasie médullaire} (neutropénie fébrile risque septicémie) $\rightarrow$ Prophylaxie G-CSF (filgrastim), antibiotiques large spectre si fièvre, transfusion si anémie/plaquettes basses.
            \item \textbf{Neuropathie périphérique} (Cisplatine) + \textbf{Ototoxicité} (surdi haute fréquence) + \textbf{Néphrotoxicité} $\rightarrow$ Hydratation forcée IV (3L/j) avant/après cisplatine, surveillance créatininémie, audiométrie.
            \item \textbf{Alopécie} (réversible), \textbf{Nausées/vomissements} (antiémétiques : ondansétrone, dexaméthasone, aprepitant), \textbf{Infertilité} (conservation sperme AVANT traitement).
        \end{itemize}
    \item \textbf{Message prévention :} « \textit{Le cancer du testicule touche l'homme jeune (15-40 ans). Une masse testiculaire, même indolore, n'est JAMAIS normale. Il faut consulter IMMÉDIATEMENT (centre de santé, hôpital de district). L'autopalpation mensuelle (après douche chaude) permet de la détecter tôt. Guérison quasi certaine si pris à temps. Ne pas attendre, ne pas avoir honte, ne pas recourir au tradithérapeute en première intention pour une "boule" au testicule.} »
\end{enumerate}
\end{corrige}

\coursec{Interprétation de données expérimentales sur la prolifération cellulaire}

Nous venons de voir que la cancérisation résulte d'une perte de contrôle de la division cellulaire : les cellules se multiplient alors de façon anarchique. Mais comment les biologistes et les médecins \emph{mesurent-ils} concrètement cette prolifération ? Comment savoir si un tissu se divise normalement ou de manière excessive ? Comment repérer, dans une population de millions de cellules, celles qui sont en train de répliquer leur ADN ou de réaliser leur mitose ?

Ces questions ne sont pas théoriques. Au Centre Pasteur du Cameroun ou dans les laboratoires d'anatomopathologie des hôpitaux de Yaoundé et de Douala, on analyse chaque jour des prélèvements de tissus pour évaluer leur activité de prolifération : c'est une étape clé du diagnostic des cancers. Cette dernière section te donne les outils pour \cle{interpréter} toi-même les documents expérimentaux que l'on te proposera au Probatoire : courbes de croissance, marquages de l'ADN, histogrammes de cytométrie et indices mitotiques.

\courssub{Suivre la croissance d'une population cellulaire en culture}

\textbf{Situation concrète.} Un chercheur met en culture des cellules humaines dans une boîte de Pétri contenant un milieu nutritif. Chaque jour, il prélève un échantillon et compte les cellules. Que va-t-il observer ?

\begin{experience}[Croissance d'une culture cellulaire]
\textbf{Protocole :} on ensemence \num{10000} cellules dans un milieu de culture renouvelé régulièrement, à \SI{37}{\celsius}. On effectue un comptage quotidien (à l'aide d'une cellule de Malassez ou d'un compteur automatique) pendant 10 jours.

\textbf{Observations :} le nombre de cellules reste d'abord stable pendant environ un jour, puis augmente très rapidement (doublement toutes les 24 h environ), avant de se stabiliser autour d'un plateau.

\textbf{Interprétation :} la courbe obtenue présente trois phases caractéristiques.
\end{experience}

\begin{popfigure}
\begin{center}
\begin{tikzpicture}
\begin{axis}[
    width=0.85\linewidth, height=6.5cm,
    xlabel={Temps (jours)}, ylabel={Nombre de cellules (milliers)},
    xmin=0, xmax=10, ymin=0, ymax=280,
    xtick={0,1,2,3,4,5,6,7,8,9,10},
    ytick={0,50,100,150,200,250},
    grid=major, grid style={popline!40},
    axis line style={popdark},
    label style={popdark},
    clip=false
]
\addplot[popcurve, domain=0:1, samples=20] {10};
\addplot[popcurve, domain=1:6, samples=100] {10*exp(0.65*(x-1))};
\addplot[popcurve, domain=6:10, samples=20] {10*exp(0.65*5)};
\draw[dashed, popmuted] (axis cs:1,0) -- (axis cs:1,10);
\draw[dashed, popmuted] (axis cs:6,0) -- (axis cs:6,256);
\node[popblue, font=\small, align=center] at (axis cs:0.5,45) {Phase de\\latence};
\node[poporange, font=\small, align=center] at (axis cs:3.3,150) {Phase\\exponentielle};
\node[popgreen, font=\small, align=center] at (axis cs:8.2,225) {Phase\\stationnaire};
\end{axis}
\end{tikzpicture}
\end{center}
\legende{Courbe de croissance d'une population de cellules en culture. Phase de latence (adaptation, pas de division), phase exponentielle (divisions actives, le nombre de cellules double à chaque génération), phase stationnaire (épuisement du milieu, encombrement : équilibre entre divisions et morts cellulaires).}
\end{popfigure}

\begin{aretenir}[Les trois phases de la croissance d'une culture]
\begin{itemize}
\item \cle{Phase de latence} : les cellules s'adaptent au milieu, synthétisent des protéines, mais ne se divisent pas encore (elles sont en G1).
\item \cle{Phase exponentielle} (ou logarithmique) : toutes les cellules se divisent à rythme régulier ; le nombre de cellules double à chaque \emph{temps de génération}. C'est la phase où le cycle cellulaire fonctionne à plein régime.
\item \cle{Phase stationnaire} : la croissance ralentit puis s'arrête par épuisement des nutriments, accumulation de déchets et \emph{inhibition par contact} (les cellules s'arrêtent de diviser quand elles se touchent — propriété perdue par les cellules cancéreuses).
\end{itemize}
\end{aretenir}

\begin{attention}[Piège classique sur la phase stationnaire]
Une courbe de croissance qui ralentit ou plafonne \textbf{ne signifie pas} que les mitoses s'arrêtent totalement. En phase stationnaire, il peut exister un \cle{équilibre dynamique} entre divisions cellulaires et morts cellulaires : des mitoses continuent, mais elles compensent exactement les cellules qui meurent. Ne rédige jamais « les cellules ne se divisent plus » sans nuance.
\end{attention}

\courssub{Repérer les cellules en phase S : le marquage de l'ADN}

\textbf{Problème :} comment savoir, à un instant donné, quelles cellules d'un tissu sont en train de répliquer leur ADN ?

L'astuce consiste à fournir aux cellules un \cle{nucléotide marqué}, c'est-à-dire une brique de l'ADN rendue détectable. Seules les cellules en \textbf{phase S} incorporent ce marqueur dans leur ADN en cours de réplication.

\begin{itemize}
\item \textbf{Méthode historique :} la \cle{thymidine tritiée} (radioactive, marquée au tritium $^3$H). Après incubation, on réalise une \emph{autoradiographie} : un film photographique appliqué sur la coupe révèle des grains d'argent au-dessus des noyaux radioactifs, donc des cellules en phase S.
\item \textbf{Méthode actuelle :} la \cle{BrdU} (bromodéoxyuridine), analogue chimique de la thymidine, non radioactive. Elle est incorporée à la place de la thymidine pendant la phase S, puis révélée par un anticorps fluorescent spécifique : les noyaux des cellules en S deviennent fluorescents au microscope.
\end{itemize}

\begin{experience}[Marquage à la thymidine tritiée — une expérience de référence]
\textbf{Protocole :} on place des cellules en culture 30 minutes en présence de thymidine tritiée (« impulsion »), puis on les remet dans un milieu non radioactif et on réalise des autoradiographies à différents temps.

\textbf{Observations :} à $t = 0$, seuls certains noyaux sont marqués (environ 30 \%) ; aucune cellule en mitose n'est marquée. Quelques heures plus tard, apparaissent des \emph{mitoses marquées}, dont la proportion augmente puis diminue.

\textbf{Interprétation :} à $t = 0$, seules les cellules en phase S incorporent le marqueur ; les mitoses (phase M) ne répliquent pas d'ADN, donc ne sont pas marquées. Le délai d'apparition des mitoses marquées correspond à la durée de G2 : les cellules marquées en S doivent traverser G2 avant d'entrer en mitose. Ce type d'expérience a permis de \emph{mesurer la durée des phases du cycle cellulaire}.
\end{experience}

\begin{exemplebox}[Application au diagnostic]
Chez un patient, on administre de la BrdU avant une biopsie tumorale. Si une forte proportion de cellules tumorales incorpore la BrdU, cela signifie que beaucoup de cellules sont en phase S : la tumeur prolifère activement. Cette information guide le choix du traitement (les chimiothérapies ciblent préférentiellement les cellules en division).
\end{exemplebox}

\courssub{Mesurer la quantité d'ADN par noyau : la cytométrie}

\textbf{Principe.} On note \cle{Q} la quantité d'ADN contenue dans un noyau en phase G1 (cellule diploïde à chromosomes simples, à une chromatide). Au cours du cycle :

\begin{center}
\begin{tabularx}{\linewidth}{|l|X|X|}
\hline
\rowcolor{popblueL} \textbf{Phase du cycle} & \textbf{Quantité d'ADN} & \textbf{Explication} \\
\hline
G1 & Q & Chromosomes à 1 chromatide, ADN non répliqué \\
\hline
S & entre Q et 2Q & Réplication en cours : toutes les valeurs intermédiaires existent \\
\hline
G2 et M & 2Q & Réplication achevée : chromosomes à 2 chromatides \\
\hline
\end{tabularx}
\end{center}

La \cle{cytométrie de flux} permet de mesurer automatiquement la quantité d'ADN de dizaines de milliers de cellules : chaque cellule, colorée par un marqueur fluorescent de l'ADN (iodure de propidium, DAPI, etc.), passe devant un faisceau laser ; l'intensité de fluorescence est proportionnelle à la quantité d'ADN. On obtient un \textbf{histogramme} : en abscisse la quantité d'ADN, en ordonnée le nombre de cellules.

\begin{popfigure}
\begin{center}
\begin{tikzpicture}
\begin{axis}[
    width=0.85\linewidth, height=6.5cm,
    xlabel={Quantité d'ADN par noyau (unités arbitraires)},
    ylabel={Nombre de cellules},
    xmin=0, xmax=3.5, ymin=0, ymax=105,
    xtick={1,2}, xticklabels={Q, 2Q},
    ytick=\empty,
    axis line style={popdark},
    label style={popdark},
]
\addplot[popcurve, domain=0.55:1.45, samples=100] {90*exp(-((x-1)/0.13)^2)};
\addplot[popcurve, domain=1.55:2.45, samples=100] {45*exp(-((x-2)/0.13)^2)};
\addplot[popcurve, domain=1.1:1.9, samples=60] {18};
\node[popblue, font=\small, align=center] at (axis cs:1,98) {Pic G1\\(Q)};
\node[poporange, font=\small, align=center] at (axis cs:2,53) {Pic G2 + M\\(2Q)};
\node[poppurple, font=\small, align=center] at (axis cs:1.5,26) {Cellules en S\\(valeurs intermédiaires)};
\end{axis}
\end{tikzpicture}
\end{center}
\legende{Histogramme de cytométrie de flux (lecture simplifiée — complément utile au Probatoire). Le premier pic correspond aux cellules en G1 (quantité Q), le second pic aux cellules en G2 et en mitose (quantité 2Q) ; entre les deux, les cellules en phase S présentent des quantités intermédiaires d'ADN.}
\end{popfigure}

\begin{attention}[Le pic 2Q ne distingue pas G2 et M]
Un pic de cellules à 2Q en cytométrie regroupe les cellules en \textbf{G2 ET en mitose}, car toutes deux possèdent la même quantité d'ADN (2Q). La cytométrie de flux seule \textbf{ne permet pas de les distinguer} : il faut un marqueur supplémentaire (par exemple une protéine spécifique de la mitose comme la phospho-histone H3) ou une observation microscopique. C'est un piège fréquent dans les sujets.
\end{attention}

\begin{exemplebox}[Lecture d'un histogramme]
Sur l'histogramme ci-dessus, le pic G1 est environ deux fois plus haut que le pic G2/M. Cela reflète simplement les \emph{durées relatives} des phases : dans un cycle de 24 h, G1 dure environ 10 h, S 8 h, G2 + M environ 6 h. Plus une phase dure longtemps, plus la proportion de cellules qui s'y trouvent à un instant donné est grande.
\end{exemplebox}

\courssub{L'indice mitotique : mesurer l'activité de prolifération d'un tissu}

Sur une coupe histologique observée au microscope, on peut reconnaître directement les cellules en mitose (chromosomes condensés, figures de division visibles). Leur proportion renseigne sur l'activité de renouvellement du tissu.

\begin{definition}[Indice mitotique]
L'\cle{indice mitotique} est le rapport entre le nombre de cellules en mitose et le nombre total de cellules observées, exprimé en pourcentage :
\[ \text{Indice mitotique} = \frac{\text{nombre de cellules en mitose}}{\text{nombre total de cellules}} \times 100 \]
Il traduit l'\textbf{activité de prolifération} d'un tissu : plus il est élevé, plus le tissu se divise activement.
\end{definition}

\begin{exemplebox}[Calcul d'un indice mitotique]
Sur une coupe de tissu, on observe 500 cellules dont 25 en mitose :
\[ \text{Indice mitotique} = \frac{25}{500} \times 100 = 5\,\% \]
Un indice de 5 \% est le signe d'un tissu en \textbf{renouvellement actif} (comme la moelle osseuse ou l'épithélium intestinal). À titre de comparaison, un tissu peu renouvelé comme le foie normal présente un indice inférieur à 1 \%, tandis qu'une tumeur cancéreuse agressive peut dépasser 10 à 20 \%.
\end{exemplebox}

\begin{savaistu}[L'indice mitotique au service du diagnostic du cancer]
En anatomopathologie, le comptage des mitoses sur une biopsie fait partie des critères de \emph{gradation} des tumeurs : un indice mitotique élevé indique une tumeur qui prolifère vite, donc plus agressive, nécessitant un traitement plus intensif. C'est l'une des applications médicales directes de la connaissance du cycle cellulaire.
\end{savaistu}

\courssub{La démarche d'analyse d'un document expérimental}

Au Probatoire, on te demandera souvent d'exploiter un document (courbe, histogramme, microphotographie, tableau de comptage). Voici la méthode infaillible.

\begin{methode}[Interpréter un document expérimental en 4 temps]
\begin{enumerate}
\item \textbf{Lire} le document : identifier les axes (grandeurs et unités), la légende, les conditions expérimentales (témoin / expérience). Annoncer ce que le document représente.
\item \textbf{Décrire} les résultats objectivement, avec des valeurs chiffrées relevées sur le document (« la courbe passe de ... à ... entre ... et ... »), sans encore interpréter.
\item \textbf{Comparer} : mettre en évidence les différences entre les situations (témoin/expérience, sain/malade, avant/après traitement) et dégager ce qui varie.
\item \textbf{Conclure} en reliant les observations aux connaissances sur le cycle cellulaire : quelle phase est affectée ? Que devient la prolifération ? Répondre explicitement au problème posé.
\end{enumerate}
\end{methode}

\begin{exoresolu}[Exploitation d'un comptage cellulaire]
\textbf{Énoncé.} On cultive deux lots de cellules dans les mêmes conditions. Le lot A reçoit un milieu normal ; le lot B reçoit un milieu additionné d'une substance X. Après 48 h, on observe au microscope 400 cellules par lot et on compte les mitoses : 20 mitoses dans le lot A, 60 dans le lot B. Que peut-on conclure sur l'effet de la substance X ?
\tcblower
\textbf{Solution détaillée.}

\emph{Étape 1 — Lecture :} le document compare deux lots de cellules ; la grandeur mesurée est le nombre de mitoses sur 400 cellules après 48 h ; le lot A est le témoin, le lot B l'expérience (variable : présence de X).

\emph{Étape 2 — Description et calcul :}
\begin{align*}
\text{Lot A (témoin) :} \quad & \frac{20}{400} \times 100 = 5\,\% \\
\text{Lot B (+ X) :} \quad & \frac{60}{400} \times 100 = 15\,\%
\end{align*}
L'indice mitotique du lot B est trois fois plus élevé que celui du lot A.

\emph{Étape 3 — Comparaison :} à conditions identiques, la seule différence est la présence de la substance X ; c'est donc elle qui est responsable de l'augmentation de la proportion de cellules en mitose.

\emph{Étape 4 — Conclusion :} la substance X \textbf{stimule la prolifération cellulaire} (elle accélère le cycle cellulaire ou recrute davantage de cellules en division). Une telle substance pourrait être un facteur de croissance ; si elle agissait de façon incontrôlée dans un organisme, elle pourrait favoriser une cancérisation.
\end{exoresolu}

\begin{exercice}[À toi de jouer]
Une biopsie de tissu tumoral est analysée par cytométrie de flux. L'histogramme montre un pic G1 très réduit et un pic 2Q très développé par rapport à un tissu sain. 1) Que signifie le pic 2Q ? 2) Quelle hypothèse peux-tu formuler sur le déroulement du cycle dans ces cellules tumorales ? 3) Quelle vérification complémentaire proposerais-tu ?
\end{exercice}

\begin{corrige}[Exercice 1]
1) Le pic 2Q regroupe les cellules contenant une quantité double d'ADN, c'est-à-dire les cellules en G2 et en mitose (chromosomes à deux chromatides).

2) Un pic 2Q anormalement développé suggère qu'une grande proportion de cellules se trouve en G2/M : soit le cycle est accéléré et beaucoup de cellules arrivent en fin de cycle, soit il existe un \emph{blocage} en G2 ou en mitose (les cellules s'y accumulent). Compte tenu du contexte tumoral, une prolifération intense est l'hypothèse la plus probable.

3) Pour trancher, on peut réaliser une observation microscopique de la coupe et calculer l'\textbf{indice mitotique} : s'il est élevé, les cellules du pic 2Q sont effectivement nombreuses en mitose (prolifération active) ; s'il est faible, les cellules sont plutôt bloquées en G2. On peut aussi utiliser un marqueur spécifique de la mitose en cytométrie (ex. : anticorps anti-phospho-histone H3).
\end{corrige}

\begin{aretenir}[L'essentiel de la section]
\begin{itemize}
\item La croissance d'une culture cellulaire suit trois phases : \cle{latence}, \cle{exponentielle}, \cle{stationnaire} (équilibre possible entre divisions et morts).
\item Les nucléotides marqués (thymidine tritiée, BrdU) sont incorporés uniquement en \cle{phase S} : ils repèrent les cellules qui répliquent leur ADN.
\item La \cle{cytométrie} mesure la quantité d'ADN par noyau : G1 = Q, S = intermédiaire, G2 + M = 2Q (sans distinguer G2 de M).
\item L'\cle{indice mitotique} $= \dfrac{\text{mitoses}}{\text{cellules totales}} \times 100$ mesure l'activité de prolifération d'un tissu ; il est élevé dans les tissus à renouvellement rapide et dans les tumeurs.
\item Toute interprétation suit la démarche : \textbf{lire → décrire → comparer → conclure} en lien avec le cycle cellulaire.
\end{itemize}
\end{aretenir}

\newpage

\coursec{Activité d'intégration}

\coursec{ND01 : Sensibilisation sur la nécessité de la mitose pour le maintien de l'identité biologique des organismes}

\courssub{Objectifs de l'activité}

À l'issue de cette activité d'intégration, tu seras capable de :

\begin{itemize}
\item \cle{identifier} les différentes phases de la mitose (prophase, métaphase, anaphase, télophase) ainsi que l'interphase sur des schémas de cellules observées au microscope ;
\item \cle{calculer} un indice mitotique à partir d'un comptage de cellules et comparer des échantillons de tissus ;
\item \cle{exploiter} des mesures de quantité d'ADN par noyau pour repérer la réplication de l'ADN et détecter une anomalie ;
\item \cle{interpréter} une courbe de prolifération cellulaire au cours du temps ;
\item \cle{rédiger} une conclusion argumentée reliant la mitose à un phénomène de la vie courante : la cicatrisation d'une plaie.
\end{itemize}

\courssub{Situation}

\textbf{Au centre de santé de quartier.} À la suite d'une coupure au pied, survenue en jouant au football sur le terrain du quartier à Yaoundé, un élève de Première D se rend au centre de santé intégré. L'infirmier nettoie la plaie et lui explique : « Ta peau va se réparer toute seule en quelques jours, car tes cellules se multiplient. » Intrigué, l'élève demande à son professeur de SVT comment de simples divisions cellulaires peuvent refermer une plaie sans modifier la nature des cellules de la peau. Le professeur propose alors à la classe de mener l'enquête, à partir d'un dossier constitué de trois documents.

\textbf{Document 1 : microphotographie schématisée du bord de la plaie.} Une lame histologique du bord de la plaie, prélevée 3 jours après la blessure, est observée au microscope optique (grossissement $\times 400$). Le champ comporte \cle{200 cellules}. Le schéma ci-dessous représente un échantillon représentatif de ce champ.

\begin{popfigure}
\begin{tikzpicture}[scale=0.9]
% Interphase cells
\draw[fill=popblueL] (0,0) circle (0.55); \draw[fill=popblue!50] (0,0) circle (0.3);
\draw[fill=popblueL] (1.6,0.3) circle (0.55); \draw[fill=popblue!50] (1.6,0.3) circle (0.3);
% Prophase
\draw[fill=popgreenL] (3.2,0) circle (0.55);
\draw[popdark,thick] (3.05,0.1) -- (3.35,-0.1); \draw[popdark,thick] (3.05,-0.15) -- (3.35,0.15);
\draw[popdark,thick] (3.2,0.25) -- (3.2,-0.25);
% Metaphase
\draw[fill=poporangeL] (4.8,0) circle (0.55);
\draw[popdark,thick] (4.65,0) -- (4.95,0); \draw[popdark,thick] (4.65,0.08) -- (4.95,0.08); \draw[popdark,thick] (4.65,-0.08) -- (4.95,-0.08);
% Anaphase
\draw[fill=poppinkL] (6.4,0) ellipse (0.75 and 0.5);
\draw[popdark,thick] (6.0,0.15) -- (6.2,0.3); \draw[popdark,thick] (6.0,-0.15) -- (6.2,-0.3);
\draw[popdark,thick] (6.8,0.15) -- (6.6,0.3); \draw[popdark,thick] (6.8,-0.15) -- (6.6,-0.3);
% Telophase
\draw[fill=poppurpleL] (8.3,0.25) circle (0.4); \draw[fill=poppurpleL] (8.3,-0.45) circle (0.4);
\draw[fill=poppurple!50] (8.3,0.25) circle (0.2); \draw[fill=poppurple!50] (8.3,-0.45) circle (0.2);
% Labels
\draw[popmuted,thin] (0,-0.6) -- (0,-1.1) node[below,font=\small]{a};
\draw[popmuted,thin] (3.2,-0.6) -- (3.2,-1.1) node[below,font=\small]{b};
\draw[popmuted,thin] (4.8,-0.6) -- (4.8,-1.1) node[below,font=\small]{c};
\draw[popmuted,thin] (6.4,-0.6) -- (6.4,-1.1) node[below,font=\small]{d};
\draw[popmuted,thin] (8.3,-0.9) -- (8.3,-1.1) node[below,font=\small]{e};
\end{tikzpicture}
\legende{Schéma de cellules observées au bord de la plaie : a : noyau à chromatine diffuse ; b : chromosomes condensés, enveloppe nucléaire en disparition ; c : chromosomes alignés sur la plaque équatoriale ; d : chromosomes à chromatides séparées migrant vers les pôles ; e : deux noyaux fils en reconstitution.}
\end{popfigure}

Le comptage complet du champ (200 cellules) donne les résultats suivants :

\begin{center}
\begin{tabularx}{\linewidth}{|l|X|X|X|X|X|}
\hline
\rowcolor{popblueL} \textbf{État de la cellule} & Interphase & Prophase & Métaphase & Anaphase & Télophase \\
\hline
\textbf{Nombre de cellules} & 176 & 10 & 6 & 5 & 3 \\
\hline
\end{tabularx}
\end{center}

\textbf{Document 2 : quantité d'ADN par noyau.} Par une technique de mesure de l'ADN nucléaire, on détermine la quantité moyenne d'ADN par noyau (en unités arbitraires, ua) dans trois échantillons : peau saine, bord de la plaie (3 jours après la blessure) et une tumeur cutanée prélevée au laboratoire d'anatomopathologie de l'hôpital central.

\begin{center}
\begin{tabularx}{\linewidth}{|l|X|X|X|}
\hline
\rowcolor{popblueL} \textbf{Échantillon} & Peau saine & Bord de plaie & Tumeur \\
\hline
\textbf{Cellules en mitose (sur 200)} & 4 & 24 & 62 \\
\hline
\textbf{Quantité d'ADN par noyau (ua)} & 6,6 & 6,6 & 13,2 (nombreux noyaux) \\
\hline
\end{tabularx}
\end{center}

\textbf{Document 3 : évolution du nombre de cellules au bord de la plaie.} On compte chaque jour le nombre de cellules présentes dans un volume de référence au bord de la plaie.

\begin{center}
\begin{tabularx}{\linewidth}{|l|X|X|X|X|X|X|}
\hline
\rowcolor{popblueL} \textbf{Temps (jours)} & 0 & 1 & 2 & 3 & 4 & 5 \\
\hline
\textbf{Nombre de cellules} & 100 & 110 & 150 & 260 & 420 & 600 \\
\hline
\end{tabularx}
\end{center}

\courssub{Consignes}

\begin{enumerate}
\item À l'aide du document 1, identifie la phase du cycle cellulaire représentée par chacun des schémas a, b, c, d et e. Justifie chaque identification par un indice observable.
\item Définis l'\cle{indice mitotique}, puis calcule sa valeur au bord de la plaie à partir du comptage des 200 cellules.
\item Calcule de même les indices mitotiques de la peau saine et de la tumeur (document 2). Compare les trois valeurs et propose une interprétation.
\item Explique, à l'aide de la quantité d'ADN par noyau, pourquoi la peau saine et le bord de la plaie présentent un caryotype stable, et ce que révèle la valeur anormale mesurée dans la tumeur.
\item Trace la courbe du nombre de cellules au bord de la plaie en fonction du temps (document 3). Décris son allure et explique comment la mitose rend compte de cette évolution.
\item Rédige une conclusion argumentée (10 à 15 lignes) répondant à la question de l'élève : comment la mitose assure-t-elle la réparation de la plaie, et pourquoi le profil de la tumeur est-il anormal ?
\end{enumerate}

\courssub{Ressources à mobiliser}

\begin{aretenir}[Formules et notions utiles]
\begin{itemize}
\item \cle{Indice mitotique} (IM) : $IM = \dfrac{\text{nombre de cellules en mitose}}{\text{nombre total de cellules observées}} \times 100$ (exprimé en \%) .
\item Sont en mitose les cellules en \textbf{prophase, métaphase, anaphase et télophase} ; l'interphase n'est pas une phase de la mitose.
\item Le \cle{cycle cellulaire} comprend l'interphase (G1, S, G2) et la mitose ; la \cle{réplication de l'ADN} a lieu en phase S : la quantité d'ADN passe de $Q$ à $2Q$, puis revient à $Q$ après la division.
\item La mitose conserve le \cle{caryotype} : chaque cellule fille reçoit le même nombre de chromosomes (ici $2n = 46$) que la cellule mère.
\item Une prolifération rapide et incontrôlée, avec des anomalies de la répartition de l'ADN, est un signe de \cle{cancérisation}.
\end{itemize}
\end{aretenir}

\courssub{Démarche et corrigé commenté}

\begin{exoresolu}[Enquête sur la cicatrisation : corrigé complet]
\textbf{Énoncé.} Identifier les phases du document 1, calculer les indices mitotiques des trois échantillons, exploiter les documents 2 et 3, puis conclure sur le rôle de la mitose dans la cicatrisation.

\tcblower
\textbf{1. Identification des phases (document 1).}
\begin{itemize}
\item Schéma a : noyau intact à chromatine diffuse $\Rightarrow$ \textbf{interphase} (la cellule n'est pas en division ; les chromosomes sont décondensés).
\item Schéma b : chromosomes condensés et individualisés, enveloppe nucléaire qui disparaît $\Rightarrow$ \textbf{prophase}.
\item Schéma c : chromosomes alignés sur la plaque équatoriale, au centre de la cellule $\Rightarrow$ \textbf{métaphase}.
\item Schéma d : chromatides sœurs séparées migrant vers les pôles opposés $\Rightarrow$ \textbf{anaphase}.
\item Schéma e : deux lots de chromosomes aux pôles, deux noyaux fils en reconstitution $\Rightarrow$ \textbf{télophase}.
\end{itemize}

\textbf{2. Indice mitotique au bord de la plaie.} Les cellules en mitose sont celles en prophase, métaphase, anaphase et télophase :
\[ N_{\text{mitose}} = 10 + 6 + 5 + 3 = 24 \]
\[ IM_{\text{plaie}} = \frac{24}{200} \times 100 = 12\,\% \]

\textbf{3. Comparaison des trois échantillons.}
\begin{align*}
IM_{\text{peau saine}} &= \frac{4}{200} \times 100 = 2\,\% \\
IM_{\text{plaie}} &= \frac{24}{200} \times 100 = 12\,\% \\
IM_{\text{tumeur}} &= \frac{62}{200} \times 100 = 31\,\%
\end{align*}
Interprétation : le bord de la plaie divise \textbf{6 fois plus} que la peau saine : la mitose est stimulée pour remplacer les cellules détruites. La tumeur présente un indice encore bien supérieur (31 \%), signe d'une prolifération \textbf{excessive et incontrôlée}.

\textbf{4. Quantité d'ADN et stabilité du caryotype.} La peau saine et le bord de la plaie présentent la même quantité d'ADN par noyau (6,6 ua) : la réplication en phase S ($Q \to 2Q$) est suivie d'une répartition équitable des chromosomes entre les deux cellules filles ($2Q \to Q$), ce qui conserve le caryotype ($2n = 46$). Dans la tumeur, de nombreux noyaux contiennent 13,2 ua, soit le double : la répartition de l'ADN est déréglée (divisions anarchiques, anomalies chromosomiques), ce qui traduit un dysfonctionnement de la mitose caractéristique de la cancérisation.

\textbf{5. Courbe de prolifération (document 3).}

\begin{popfigure}
\begin{tikzpicture}
\begin{axis}[width=0.85\linewidth,height=6cm,xlabel={Temps (jours)},ylabel={Nombre de cellules},xmin=0,xmax=5.5,ymin=0,ymax=650,xtick={0,1,2,3,4,5},grid=both,grid style={popline!40}]
\addplot[popcurve,mark=*,mark options={fill=poporange}] coordinates {(0,100) (1,110) (2,150) (3,260) (4,420) (5,600)};
\end{axis}
\end{tikzpicture}
\legende{Évolution du nombre de cellules au bord de la plaie au cours des 5 jours suivant la blessure.}
\end{popfigure}

La courbe est d'abord presque plate (jours 0 à 1 : phase de latence, les cellules se préparent en interphase), puis s'élève de plus en plus vite (croissance quasi exponentielle) : chaque mitose double le nombre de cellules, ce qui explique l'accélération. Au jour 5, le nombre de cellules a été multiplié par 6 : la plaie se referme.

\textbf{6. Conclusion rédigée (modèle).} La réparation de la plaie repose sur la multiplication des cellules du bord de la blessure par mitose. L'indice mitotique élevé (12 \% contre 2 \% dans la peau saine) montre que les cellules se divisent activement. Grâce à la réplication de l'ADN en phase S, puis à la séparation rigoureuse des chromatides en anaphase, chaque cellule fille reçoit un lot complet et identique de chromosomes : la quantité d'ADN par noyau reste de 6,6 ua et le caryotype est conservé. Les nouvelles cellules, génétiquement identiques aux cellules de la peau, comblent progressivement la plaie, comme le montre la courbe de prolifération. La mitose assure ainsi le maintien de l'identité biologique du tissu tout en permettant sa régénération. En revanche, la tumeur présente un indice mitotique très élevé (31 \%) et des noyaux à quantité d'ADN doublée : le contrôle du cycle cellulaire est perdu, les divisions deviennent anarchiques et le patrimoine génétique n'est plus correctement transmis : c'est la cancérisation.
\end{exoresolu}

\courssub{Bilan de l'activité}

Cette enquête t'a permis de relier un phénomène quotidien — la cicatrisation d'une coupure — aux mécanismes cellulaires étudiés dans la leçon :

\begin{itemize}
\item l'\cle{identification des phases} de la mitose repose sur des indices morphologiques précis (condensation des chromosomes, plaque équatoriale, migration des chromatides, reconstitution des noyaux) ;
\item l'\cle{indice mitotique} est un outil quantitatif simple pour comparer l'activité de division de différents tissus : il distingue une prolifération normale et contrôlée (cicatrisation) d'une prolifération pathologique (tumeur) ;
\item la \cle{réplication de l'ADN} suivie d'une répartition équitable garantit la conservation du caryotype et donc le maintien de l'identité biologique de l'organisme ;
\item un \cle{dérèglement du cycle cellulaire} (divisions trop nombreuses, ADN mal réparti) est la signature de la cancérisation, ce qui justifie l'importance du dépistage précoce dans les services d'anatomopathologie des hôpitaux camerounais.
\end{itemize}

\begin{attention}[Erreurs fréquentes à éviter]
\begin{itemize}
\item Ne compte \textbf{jamais} les cellules en interphase dans le calcul de l'indice mitotique : l'interphase n'est pas une phase de la mitose.
\item N'oublie pas de multiplier par 100 pour exprimer l'indice en pourcentage, et précise toujours le nombre total de cellules observées.
\item Une quantité d'ADN de $2Q$ dans un noyau ne signifie pas forcément une anomalie : c'est normal en fin d'interphase (après la phase S). L'anomalie, c'est la \emph{persistance} de noyaux à $2Q$ ou plus dans un tissu où les divisions sont anarchiques.
\item Dans la conclusion, cite toujours des \textbf{données chiffrées} du dossier pour appuyer ton argumentation : c'est ce qui fait la différence entre une affirmation et une démonstration scientifique.
\end{itemize}
\end{attention}

\newpage

\coursec{Méthodes et exercices résolus}

\coursec{Sensibilisation sur la nécessité de la mitose pour le maintien de l'identité biologique des organismes}

\courssub{Le cycle cellulaire : alternance interphase et mitose}

\begin{definition}[Cycle cellulaire]
Le \cle{cycle cellulaire} est l'ensemble ordonné des événements qui conduisent une cellule depuis sa formation par division d'une cellule mère jusqu'à sa propre division en deux cellules filles. Il comprend deux grandes périodes :
\begin{itemize}
\item l'\cle{interphase} (phase de croissance, de synthèse et de préparation), longue et métaboliquement active ;
\item la \cle{mitose} (phase de division nucléaire), brève, suivie de la \cle{cytokinèse} (division cytoplasmique).
\end{itemize}
\end{definition}

\begin{propriete}[Phases de l'interphase]
L'interphase se subdivise en trois phases successives :
\begin{center}
\begin{tabularx}{\linewidth}{|l|X|X|}
\hline
\rowcolor{popblueL} \textbf{Phase} & \textbf{Activités principales} & \textbf{Événement clé pour la division} \\
\hline
G1 (Gap 1) & Croissance cellulaire, synthèse de protéines et d'organites, fonctionnement métabolique normal & Atteinte de la \cle{taille critique} ; décision de diviser (point de restriction R) \\
\hline
S (Synthesis) & \textbf{Réplication de l'ADN} (chaque chromosome passe d'une à deux chromatides sœurs) & Doublement de la quantité d'ADN ($Q \rightarrow 2Q$) ; le nombre de chromosomes reste $2n$ \\
\hline
G2 (Gap 2) & Préparation de la mitose : synthèse des protéines du fuseau (tubuline), vérification de l'intégrité de l'ADN réplicité & Contrôle au \cle{point de contrôle G2/M} (checkpoint) avant d'entrer en mitose \\
\hline
\end{tabularx}
\end{center}
\end{propriete}

\begin{aretenir}[Repères temporels et quantitatifs au cours du cycle]
Pour une cellule humaine typique en culture (durée totale $\approx 24$ h) :
\begin{itemize}
\item G1 $\approx 11$ h ; S $\approx 8$ h ; G2 $\approx 4$ h ; M (mitose + cytokinèse) $\approx 1$ h.
\item Quantité d'ADN : $Q$ en G1 $\rightarrow$ $2Q$ en G2 et au début de la mitose $\rightarrow$ $Q$ dans chaque cellule fille après cytokinèse.
\item Nombre de chromosomes : $2n$ (ex. 46) en G1, S, G2, métaphase ; $4n$ (chromosomes à une chromatide) dans la cellule entière en anaphase ; $2n$ dans chaque noyau fille en télophase.
\end{itemize}
\end{aretenir}

\begin{experience}[Mise en évidence des divisions mitotiques : pointe de racine d'oignon]
\textbf{Objectif :} Observer les différentes phases de la mitose dans un tissu méristématique à forte activité proliférative.
\begin{methode}[Protocole simplifié]
\begin{enumerate}
\item \textbf{Matériel :} Racines d'oignon (\emph{Allium cepa}) germées 2--3 jours, acide chlorhydrique 1 N, colorant (carmin acétique, orcéine ou bleu de toluidine), lames et lamelles, microscope optique.
\item \textbf{Fixation :} Couper l'extrémité de la racine (1--2 mm), plonger dans l'acide chlorhydrique 1 N à 60 °C pendant 1--2 min (ramollissement des parois et arrêt de la division).
\item \textbf{Coloration :} Rincer à l'eau, placer sur lame, ajouter une goutte de colorant, chauffer doucement (sans ébullition) 30 s.
\item \textbf{Ecrasement :} Couvrir d'une lamelle, interposer du papier buvard, appuyer fermement avec le pouce pour obtenir une monocouche cellulaire.
\item \textbf{Observation :} Au microscope (objectif $\times 40$ ou $\times 100$), repérer la zone méristématique (cellules petites, isodiamétriques, noyaux gros) et identifier les phases.
\end{enumerate}
\textbf{Observations attendues :} Majorité de noyaux en interphase (noyau clair, nucléole visible) ; quelques cellules en prophase (chromosomes condensés), métaphase (alignement), anaphase (migration en V), télophase (deux noyaux en formation, plaque cellulaire).
\textbf{Interprétation :} La fréquence des phases observées est proportionnelle à leur durée relative : interphase $\gg$ prophase $>$ métaphase $\approx$ anaphase $\approx$ télophase.
\end{methode}
\end{experience}

\begin{popfigure}
\begin{tikzpicture}[scale=0.7, every node/.style={font=\small}]
% Styles
\tikzset{
    chroma/.style={draw=popdark, thick, fill=poporange!30, rounded corners=2pt},
    chromb/.style={draw=popdark, thick, fill=popblue!30, rounded corners=2pt},
    spindle/.style={draw=poppurple, dashed, thick},
    pole/.style={draw=popgreen, fill=popgreen!20, circle, minimum size=4pt},
    nuclearenv/.style={draw=popdark, thick, dashed},
    cellmem/.style={draw=popdark, very thick},
}
% Interphase
\begin{scope}[xshift=0cm]
\node[cellmem, ellipse, minimum width=3cm, minimum height=2.5cm] (c1) at (0,0) {};
\node[nuclearenv, ellipse, minimum width=2cm, minimum height=1.5cm] (n1) at (0,0) {};
\foreach \i in {1,...,30} {
    \pgfmathsetmacro{\ang}{rand*360}
    \pgfmathsetmacro{\rx}{0.7+rand*0.5}
    \pgfmathsetmacro{\ry}{0.5+rand*0.4}
    \draw[popmuted, thin] (\ang:\rx cm and \ry cm) -- ++(\ang:0.15);
}
\node[below=0.2cm of c1] {Interphase (G1)};
\node[above=0.2cm of n1, popmuted] {Chromatine diffuse};
\node[right=0.2cm of n1, popmuted] {Nucléole};
\end{scope}

% Prophase
\begin{scope}[xshift=5cm]
\node[cellmem, ellipse, minimum width=3cm, minimum height=2.5cm] (c2) at (0,0) {};
\foreach \i/\style in {1/chroma, 2/chroma, 3/chromb, 4/chromb} {
    \pgfmathsetmacro{\ang}{-45+\i*30}
    \pgfmathsetmacro{\rx}{0.8+rand*0.3}
    \pgfmathsetmacro{\ry}{0.5+rand*0.2}
    \node[\style, minimum width=0.6cm, minimum height=0.15cm, rotate=\ang] at (\ang:\rx cm and \ry cm) {};
}
\node[below=0.2cm of c2] {Prophase};
\node[popmuted, align=center] at (0,0.8){Enveloppe nucléaire\\(en disparition)};
\draw[spindle] (-1.2,0) -- (-1.5,0.8) (-1.2,0) -- (-1.5,-0.8);
\draw[spindle] (1.2,0) -- (1.5,0.8) (1.2,0) -- (1.5,-0.8);
\node[pole] at (-1.2,0) {}; \node[pole] at (1.2,0) {};
\node[popmuted, left] at (-1.2,0) {Centrioles};
\end{scope}

% Metaphase
\begin{scope}[xshift=10cm]
\node[cellmem, ellipse, minimum width=3cm, minimum height=2.5cm] (c3) at (0,0) {};
\foreach \i/\style in {1/chroma, 2/chroma, 3/chromb, 4/chromb} {
    \node[\style, minimum width=0.6cm, minimum height=0.15cm] at (0, -0.6+0.4*\i) {};
    \draw[spindle] (0, -0.6+0.4*\i) -- (-1.2, 1.2);
    \draw[spindle] (0, -0.6+0.4*\i) -- (-1.2, -1.2);
    \draw[spindle] (0, -0.6+0.4*\i) -- (1.2, 1.2);
    \draw[spindle] (0, -0.6+0.4*\i) -- (1.2, -1.2);
}
\node[pole] at (-1.2, 1.2) {}; \node[pole] at (-1.2, -1.2) {};
\node[pole] at (1.2, 1.2) {}; \node[pole] at (1.2, -1.2) {};
\node[below=0.2cm of c3] {Métaphase};
\node[popmuted, above] at (0,1.2) {Plaque équatoriale};
\end{scope}

% Anaphase
\begin{scope}[xshift=15cm]
\node[cellmem, ellipse, minimum width=3.5cm, minimum height=3cm] (c4) at (0,0) {};
% Pole gauche
\foreach \i/\style in {1/chroma, 2/chroma, 3/chromb, 4/chromb} {
    \pgfmathsetmacro{\y}{-0.6+0.4*\i}
    \node[\style, minimum width=0.3cm, minimum height=0.15cm, rotate=-30] at (-1.0, \y) {};
    \draw[spindle, ->] (-1.0, \y) -- (-1.5, \y+0.3);
}
% Pole droit
\foreach \i/\style in {1/chroma, 2/chroma, 3/chromb, 4/chromb} {
    \pgfmathsetmacro{\y}{-0.6+0.4*\i}
    \node[\style, minimum width=0.3cm, minimum height=0.15cm, rotate=30] at (1.0, \y) {};
    \draw[spindle, ->] (1.0, \y) -- (1.5, \y+0.3);
}
\node[pole] at (-1.5, 1.5) {}; \node[pole] at (-1.5, -1.5) {};
\node[pole] at (1.5, 1.5) {}; \node[pole] at (1.5, -1.5) {};
\node[below=0.2cm of c4] {Anaphase};
\node[popmuted, align=center] at (0,1.3){Chromatides sœurs\\(séparées)};
\end{scope}

% Telophase
\begin{scope}[xshift=20cm]
\node[cellmem, ellipse, minimum width=4cm, minimum height=3cm] (c5) at (0,0) {};
% Plaque cellulaire
\draw[popgreen, very thick] (-2, -1.5) -- (2, 1.5);
% Noyau gauche
\node[nuclearenv, ellipse, minimum width=1.5cm, minimum height=1.2cm] (n5a) at (-1,0) {};
\foreach \i in {1,...,4} {
    \pgfmathsetmacro{\ang}{rand*360}
    \pgfmathsetmacro{\rx}{0.3+rand*0.3}
    \pgfmathsetmacro{\ry}{0.2+rand*0.2}
    \draw[popmuted, thin] (-1,0) ++(\ang:\rx cm and \ry cm) -- ++(\ang:0.1);
}
% Noyau droit
\node[nuclearenv, ellipse, minimum width=1.5cm, minimum height=1.2cm] (n5b) at (1,0) {};
\foreach \i in {1,...,4} {
    \pgfmathsetmacro{\ang}{rand*360}
    \pgfmathsetmacro{\rx}{0.3+rand*0.3}
    \pgfmathsetmacro{\ry}{0.2+rand*0.2}
    \draw[popmuted, thin] (1,0) ++(\ang:\rx cm and \ry cm) -- ++(\ang:0.1);
}
\node[below=0.2cm of c5] {Télophase + Cytokinèse};
\node[popmuted, align=center] at (-1,0.9){Enveloppe nucléaire\\reconstituée};
\node[popmuted, align=center] at (0,-1.6){Plaque cellulaire\\(végétale)};
\end{scope}
\end{tikzpicture}
\legende{Schémas des principales étapes du cycle cellulaire dans une cellule animale ($2n=4$, deux paires de chromosomes homologues distinguées par la couleur). L'interphase montre la chromatine diffuse ; la prophase la condensation et la disparition de l'enveloppe ; la métaphase l'alignement sur la plaque équatoriale ; l'anaphase la séparation des chromatides sœurs (chromosomes à une chromatide) en forme de V ; la télophase la décondensation, la reconstitution des enveloppes nucléaires et la cytokinèse (ici plaque cellulaire type végétal pour l'exemple oignon).}
\end{popfigure}

\courssub{La réplication de l'ADN et la conservation de l'information génétique}

\begin{definition}[Réplication de l'ADN (phase S)]
La \cle{réplication de l'ADN} est le processus par lequel une molécule d'ADN double brin donne deux molécules filles identiques à la molécule mère. Elle est \cle{semi-conservatrice} : chaque molécule fille conserve un brin parental (matrice) et un brin néo-synthétisé. Elle a lieu exclusivement pendant la \cle{phase S} de l'interphase, dans le noyau.
\end{definition}

\begin{propriete}[Mécanisme et fidélité]
\begin{itemize}
\item \textbf{Origines de réplication :} multiples sur les chromosomes eucaryotes ; formation de \cle{fourches de réplication}.
\item \textbf{Enzymes clés :} \cle{ADN polymérases} (synthèse 5' $\rightarrow$ 3', lecture 3' $\rightarrow$ 5'), \cle{hélicase} (ouverture), \cle{primase} (amorce ARN), \cle{ligase} (jonction fragments d'Okazaki sur brin retardé).
\item \textbf{Fidélité :} activité de \cle{relecture (proofreading)} 3' $\rightarrow$ 5' des ADN polymérases + systèmes de \cle{réparation des mésappariements} (MMR) $\rightarrow$ taux d'erreur $\approx 10^{-9}$ à $10^{-10}$ par nucléotide.
\item \textbf{Conséquence chromosomique :} chaque chromosome (1 molécule d'ADN = 1 chromatide) devient un chromosome à \textbf{deux chromatides sœurs génétiquement identiques}, jointes au \cle{centromère}.
\end{itemize}
\end{propriete}

\begin{aretenir}[Bilan de la réplication pour le caryotype]
\begin{center}
\begin{tabular}{|c|c|c|c|}
\hline
\rowcolor{popblueL} \textbf{Stade} & \textbf{Chromosomes} & \textbf{Chromatides / chr.} & \textbf{Quantité ADN} \\
\hline
Avant réplication (G1) & $2n$ & 1 & $Q$ \\
\hline
Après réplication (G2, Prophase, Métaphase) & $2n$ & 2 & $2Q$ \\
\hline
\end{tabular}
\end{center}
\textbf{Point clé :} Le nombre de chromosomes ($2n$) ne change \textbf{pas} pendant la réplication ; c'est la structure du chromosome (1 vs 2 chromatides) et la quantité d'ADN qui changent.
\end{aretenir}

\courssub{Les phases de la mitose : description détaillée}

\begin{propriete}[Caractères distinctifs des phases mitotiques (cellule animale $2n=4$)]
\begin{center}
\begin{tabularx}{\linewidth}{|l|X|X|X|}
\hline
\rowcolor{popblueL} \textbf{Phase} & \textbf{Chromosomes} & \textbf{Fuseau / Centrioles / Enveloppe} & \textbf{Cytokinèse} \\
\hline
\textbf{Prophase} & Condensation maximale : chromosomes courts, épais, à 2 chromatides sœurs jointes au centromère. & Centrioles migrés aux pôles ; fuseau achromatique en formation (microtubules) ; \textbf{disparition de l'enveloppe nucléaire} et du nucléole. & -- \\
\hline
\textbf{Métaphase} & \textbf{Alignés sur la plaque équatoriale} (plan médian) ; centromères sur le plan, bras orientés aléatoirement. & Fuseau achevé : fibres kinétochoriennes (pôle $\leftrightarrow$ centromère), fibres polaires (pôle $\leftrightarrow$ pôle), fibres astrales. Enveloppe absente. & -- \\
\hline
\textbf{Anaphase} & \textbf{Séparation des chromatides sœurs} (devenues chromosomes à 1 chromatide) ; migration vers les pôles opposés (centromère en tête, forme en V). & Raccourcissement des fibres kinétochoriennes ; élongation de la cellule (fibres polaires). & Début de l'étranglement (cellule animale) / formation du phragmoplaste (végétale). \\
\hline
\textbf{Télophase} & Chromosomes arrivés aux pôles, \textbf{décondensation} (chromatine), réapparition des nucléoles. & Fuseau disparu ; \textbf{reconstitution de deux enveloppes nucléaires}. & \textbf{Achevée} : étranglement (animal) ou plaque cellulaire (végétal) $\rightarrow$ 2 cellules filles. \\
\hline
\end{tabularx}
\end{center}
\end{propriete}

\courssub{Comportement des chromosomes et conservation du caryotype}

\begin{propriete}[Du chromosome à deux chromatides au caryotype conservé]
\begin{enumerate}
\item \textbf{Unité de comptage :} Un chromosome est défini par \textbf{un seul centromère}. Tant que les chromatides sœurs sont jointes, c'est \emph{un} chromosome (à deux chromatides).
\item \textbf{Métaphase :} $2n$ chromosomes à 2 chromatides alignés. L'alignement assure que chaque pôles recevra un exemplaire de chaque chromosome.
\item \textbf{Anaphase :} Séparation des centromères $\rightarrow$ $2 \times 2n = 4n$ chromosomes à 1 chromatide dans la cellule entière (transitoire), répartis en deux lots identiques de $2n$ chromosomes.
\item \textbf{Télophase / Cytokinèse :} Deux noyaux $2n$ (1 chromatide/chromosome) $\rightarrow$ deux cellules filles $2n$, identiques à la mère en G1.
\end{enumerate}
\end{propriete}

\begin{methode}[Raisonner sur les quantités d'ADN (Q) et le nombre de chromosomes (n)]
Pour résoudre tout exercice portant sur l'évolution du stock d'ADN ou de chromosomes au cours du cycle cellulaire :
\begin{enumerate}
\item \textbf{Noter les données de l'espèce} : $2n$ (nombre de chromosomes en diploïdie) et $Q$ (quantité d'ADN d'une cellule en G1).
\item \textbf{Situer la cellule dans le cycle} : G1 $\rightarrow$ S $\rightarrow$ G2 $\rightarrow$ M (mitose) $\rightarrow$ cytokinèse.
\item \textbf{Appliquer les valeurs de référence} :
\end{enumerate}
\begin{center}
\begin{tabularx}{\linewidth}{|l|X|X|X|}
\hline
\rowcolor{popblueL} \textbf{Étape} & \textbf{Nb de chromosomes} & \textbf{Nb de chromatides/chromosome} & \textbf{Quantité d'ADN} \\
\hline
G1 & $2n$ & 1 & $Q$ \\
\hline
Fin de S / G2 & $2n$ & 2 & $2Q$ \\
\hline
Métaphase & $2n$ & 2 & $2Q$ \\
\hline
Anaphase (cellule entière) & $4n$ (chromosomes à 1 chromatide) & 1 & $2Q$ \\
\hline
Fin de télophase (par cellule fille) & $2n$ & 1 & $Q$ \\
\hline
\end{tabularx}
\end{center}
\begin{enumerate}
\setcounter{enumi}{3}
\item \textbf{Vérifier la conservation} : chaque cellule fille reçoit $2n$ chromosomes et une quantité $Q$ d'ADN, identiques à la cellule mère en G1 : c'est le fondement de la \cle{conservation du caryotype}.
\end{enumerate}
\textbf{Réflexe} : la quantité d'ADN double pendant la phase S (réplication semi-conservative) ; le nombre de chromosomes ne double jamais avant l'anaphase, car un chromosome à deux chromatides reste \emph{un} chromosome (un seul centromère).
\end{methode}

\begin{exoresolu}[Évolution de la quantité d'ADN dans des cellules en culture]
Des fibroblastes humains ($2n = 46$) sont mis en culture. Une cellule en phase G1 contient une quantité d'ADN $Q = \SI{6,4}{pg}$ (picogrammes).
\begin{enumerate}
\item Quelle est la quantité d'ADN d'une cellule en phase G2 ? Justifier.
\item Quelle est la quantité d'ADN de chaque cellule fille à l'issue de la mitose ?
\item Combien de chromosomes et de chromatides compte une cellule en métaphase ?
\end{enumerate}
\tcblower
\textbf{1. Quantité d'ADN en G2.}

Entre G1 et G2 se déroule la \textbf{phase S}, au cours de laquelle chaque molécule d'ADN est répliquée selon le mode \cle{semi-conservateur} : chaque chromosome passe de une à deux chromatides. La quantité d'ADN double donc :
\[ Q_{\mathrm{G2}} = 2Q = 2 \times 6{,}4 = \SI{12,8}{pg}. \]

\textbf{2. Quantité d'ADN par cellule fille.}

La mitose répartit \textbf{équitablement} les chromatides sœurs entre les deux cellules filles : chacune reçoit un exemplaire de chaque chromosome. La quantité d'ADN revient donc à sa valeur initiale :
\[ Q_{\text{fille}} = \frac{2Q}{2} = Q = \SI{6,4}{pg}. \]

\textbf{3. Chromosomes et chromatides en métaphase.}

En métaphase, la cellule entière contient $2n = 46$ chromosomes, chacun formé de 2 chromatides sœurs, soit :
\[ 46 \times 2 = 92 \text{ chromatides.} \]

\textbf{Conclusion} : la quantité d'ADN passe de \SI{6,4}{pg} (G1) à \SI{12,8}{pg} (G2) puis revient à \SI{6,4}{pg} dans chaque cellule fille ; en métaphase, la cellule contient 46 chromosomes à deux chromatides, soit 92 chromatides. La réplication suivie d'une répartition équitable assure la conservation de l'information génétique d'une génération cellulaire à l'autre.
\end{exoresolu}

\begin{methode}[Construire un raisonnement « réplication $\rightarrow$ stabilité du caryotype »]
Pour répondre à une question du type « Expliquer comment la mitose assure le maintien de l'identité biologique de l'organisme », suivre une démarche en chaîne causale :
\begin{enumerate}
\item \textbf{Phase S} : chaque molécule d'ADN est répliquée de façon \cle{semi-conservatrice} ; chaque chromosome devient formé de deux chromatides sœurs \textbf{génétiquement identiques} (copies conformes).
\item \textbf{Métaphase} : les chromosomes s'alignent sur la plaque équatoriale, chacun fixé aux deux pôles par les fibres du fuseau : cette disposition prépare une répartition symétrique.
\item \textbf{Anaphase} : les chromatides sœurs se séparent et migrent vers les pôles opposés ; chaque pôle reçoit \textbf{un exemplaire de chaque chromosome}.
\item \textbf{Télophase et cytokinèse} : deux noyaux puis deux cellules filles se forment, chacune avec $2n$ chromosomes.
\item \textbf{Conclure} : les deux cellules filles possèdent le même caryotype ($2n$ chromosomes) et la même information génétique que la cellule mère ; la mitose assure donc la \cle{stabilité du patrimoine génétique}, condition de la croissance et du renouvellement des tissus sans perte d'identité biologique.
\end{enumerate}
\textbf{Réflexe} : toujours citer les deux événements clés — \textbf{réplication de l'ADN} (copie fidèle) et \textbf{répartition équitable des chromatides} (partage exact). L'un sans l'autre ne suffit pas.
\end{methode}

\begin{exoresolu}[La mitose et le maintien de l'identité biologique]
La peau d'un élève se renouvelle en permanence : les cellules basales de l'épiderme se divisent par mitose pour remplacer les cellules mortes éliminées en surface. Expliquer, en vous appuyant sur le comportement des chromosomes, comment ce renouvellement s'effectue sans modification du caryotype ni de l'information génétique des cellules de la peau.
\tcblower
Le maintien de l'identité biologique repose sur deux mécanismes complémentaires du cycle cellulaire.

\textbf{Étape 1 : la copie fidèle de l'information (interphase, phase S).} Avant toute division, en phase S, chaque molécule d'ADN des 46 chromosomes de la cellule basale est répliquée selon le mode semi-conservateur : chaque brin parental sert de matrice à la synthèse d'un brin complémentaire. Chaque chromosome devient ainsi formé de \textbf{deux chromatides sœurs génétiquement identiques}.

\textbf{Étape 2 : le partage équitable (mitose).} En métaphase, les 46 chromosomes à deux chromatides s'alignent sur la plaque équatoriale. En anaphase, les chromatides sœurs de chaque chromosome se séparent et migrent vers les pôles opposés : chaque pôle reçoit donc \textbf{exactement un exemplaire de chacun des 46 chromosomes}. En télophase, deux noyaux à $2n = 46$ chromosomes se reconstituent, puis la cytokinèse forme deux cellules filles.

\textbf{Conclusion} : chaque cellule fille de l'épiderme possède le même nombre de chromosomes ($2n = 46$, caryotype conservé) et la même information génétique que la cellule mère. La réplication fidèle de l'ADN suivie de la répartition équitable des chromatides garantit que le renouvellement de la peau s'effectue sans perte ni modification de l'identité biologique de l'organisme.
\end{exoresolu}

\courssub{Dysfonctionnements de la mitose et cancérisation}

\begin{definition}[Points de contrôle du cycle cellulaire (Checkpoints)]
Ce sont des mécanismes de surveillance moléculaire qui bloquent la progression du cycle si des anomalies sont détectées. Les trois principaux sont :
\begin{itemize}
\item \textbf{Point G1/S (Restriction) :} Vérifie la taille cellulaire, nutriments, facteurs de croissance, \textbf{intégrité de l'ADN}. Décide de l'entrée en phase S.
\item \textbf{Point G2/M :} Vérifie l'\textbf{achèvement de la réplication} et l'\textbf{intégrité de l'ADN réplicité}. Décide de l'entrée en mitose.
\item \textbf{Point M (Métaphase/Anaphase) :} Vérifie la \textbf{fixation correcte de tous les chromosomes} au fuseau (tension aux kinétochores). Déclenche l'anaphase (via APC/C).
\end{itemize}
Si l'anomalie est irréparable $\rightarrow$ \cle{apoptose} (mort cellulaire programmée).
\end{definition}

\begin{propriete}[Bases génétiques de la cancérisation]
La cancérisation résulte de l'accumulation de \cle{mutations} dans des gènes contrôlant la prolifération :
\begin{itemize}
\item \textbf{Proto-oncogènes} $\xrightarrow{\text{mutation gain de fonction}}$ \cle{Oncogènes} : stimulent la division de façon constitutive (ex. facteurs de croissance, récepteurs, kinases de transduction, cyclines).
\item \textbf{Gènes suppresseurs de tumeurs} (ex. \emph{TP53}, \emph{RB}) $\xrightarrow{\text{mutation perte de fonction}}$ : perdent leur rôle de frein (arrêt du cycle, réparation, apoptose).
\item \textbf{Gènes de réparation de l'ADN} (ex. \emph{BRCA1/2}) : mutation $\rightarrow$ instabilité génomique $\rightarrow$ accélération de l'accumulation des mutations.
\end{itemize}
\end{propriete}

\begin{savaistu}[Cancer du col de l'utérus et HPV au Cameroun]
Au Cameroun, le cancer du col de l'utérus est le 2\textsuperscript{e} cancer chez la femme (après le sein). Il est causé par une infection persistante à \textbf{papillomavirus humain (HPV) à haut risque} (types 16, 18). Les protéines virales E6 et E7 inactivent respectivement p53 (arrêt cycle/apoptose) et pRb (contrôle G1/S), supprimant les points de contrôle. La \textbf{vaccination anti-HPV} (campagnes ciblant les jeunes filles de 9--14 ans) et le \textbf{dépistage par test HPV ou VIA/VILI} (inspection visuelle) sont les piliers de la prévention promus par le Ministère de la Santé Publique.
\end{savaistu}

\begin{methode}[Raisonner sur la cancérisation]
Pour expliquer comment un dérèglement de la division cellulaire conduit à un cancer :
\begin{enumerate}
\item \textbf{Rappeler le contrôle normal} : le cycle cellulaire est contrôlé par des points de contrôle (notamment en G1 et en métaphase) qui vérifient l'intégrité de l'ADN et la bonne fixation des chromosomes au fuseau ; une cellule anormale est bloquée ou éliminée (apoptose).
\item \textbf{Identifier le dysfonctionnement} : des \cle{mutations} de gènes contrôlant le cycle (proto-oncogènes devenus oncogènes, gènes suppresseurs de tumeurs inactivés) font perdre ce contrôle.
\item \textbf{Décrire la conséquence cellulaire} : la cellule se divise de façon \textbf{incontrôlée et illimitée}, ignorant l'inhibition par contact ; elle forme une masse de cellules anormales : la \cle{tumeur}.
\item \textbf{Préciser la gravité} : les cellules cancéreuses peuvent acquérir la capacité d'envahir les tissus voisins et de migrer par le sang ou la lymphe pour former des \cle{métastases}.
\item \textbf{Conclure} en reliant au savoir de la leçon : la stabilité du patrimoine génétique et le contrôle de la mitose sont indispensables ; leur rupture menace l'organisme entier.
\end{enumerate}
\textbf{Réflexe} : ne jamais écrire « le cancer est une maladie contagieuse » ; c'est une maladie de la division cellulaire, liée à des anomalies génétiques acquises (tabac, rayons UV, virus comme le papillomavirus...).
\end{methode}

\begin{exoresolu}[De la mutation à la tumeur]
Dans un centre de santé de Yaoundé, on sensibilise les populations au dépistage précoce du cancer du col de l'utérus, lié au papillomavirus humain (HPV). On explique que ce virus peut provoquer l'inactivation de protéines qui bloquent normalement le cycle cellulaire lorsque l'ADN est endommagé.
\begin{enumerate}
\item Quel est le rôle normal des points de contrôle du cycle cellulaire ?
\item Expliquer comment l'inactivation de ces protéines peut conduire à la formation d'une tumeur.
\end{enumerate}
\tcblower
\textbf{1. Rôle des points de contrôle.}

Les points de contrôle (en G1, en G2 et en métaphase) sont des mécanismes de surveillance du cycle cellulaire : ils vérifient que l'ADN est intact et correctement répliqué, et que tous les chromosomes sont bien fixés au fuseau avant la séparation des chromatides. Si une anomalie est détectée, le cycle est \textbf{bloqué} : la cellule répare son ADN ou, si la réparation est impossible, elle est éliminée par \cle{apoptose} (mort cellulaire programmée). Ils empêchent ainsi la multiplication de cellules porteuses d'anomalies génétiques.

\textbf{2. De l'inactivation à la tumeur.}

Si le papillomavirus inactive les protéines de contrôle (p53, pRb), une cellule du col de l'utérus dont l'ADN est endommagé ou muté \textbf{échappe au blocage} : elle poursuit son cycle et se divise par mitose malgré ses anomalies. Ses cellules filles héritent des mêmes anomalies et se divisent à leur tour de façon \textbf{incontrôlée} : les mitoses s'enchaînent sans frein, ignorant l'inhibition par contact. Il se forme alors une masse de cellules anormales : une \cle{tumeur}. Si les cellules tumorales acquièrent la capacité d'envahir les tissus voisins et de migrer par les vaisseaux sanguins ou lymphatiques, elles forment des \cle{métastases} : c'est le cancer.

\textbf{Conclusion} : les points de contrôle garantissent qu'une cellule ne se divise que si son patrimoine génétique est intact ; leur inactivation par le virus HPV supprime ce verrou, déclenchant une prolifération mitotique incontrôlée à l'origine de la tumeur — d'où l'importance de la vaccination et du dépistage précoce promus au Cameroun.
\end{exoresolu}

\courssub{Interprétation de données expérimentales sur la prolifération cellulaire}

\begin{methode}[Exploiter un tableau ou une courbe de prolifération cellulaire]
Face à des données expérimentales (comptage de cellules, courbe de croissance d'une culture, mesure de la quantité d'ADN par cellule) :
\begin{enumerate}
\item \textbf{Décrire les données avec précision} : lire les axes (grandeur, unité), relever les valeurs remarquables (valeur initiale, valeur maximale, temps de doublement).
\item \textbf{Dégager la variation} : la population cellulaire augmente-t-elle ? Selon quel rythme (phase de latence, croissance exponentielle, plateau) ?
\item \textbf{Calculer si nécessaire} : si une population double à chaque cycle de durée $T$, après $k$ cycles le nombre de cellules est :
\[ N = N_0 \times 2^{k} \qquad \text{avec } k = \frac{t}{T}. \]
\item \textbf{Interpréter biologiquement} : relier l'augmentation du nombre de cellules à la succession des mitoses ; un plateau traduit une limitation (nutriments, espace, inhibition par contact).
\item \textbf{Comparer les conditions} si l'énoncé propose plusieurs courbes (témoin / traité) : un traitement qui accélère le doublement stimule la prolifération ; un traitement qui bloque le cycle (ex. colchicine bloquant le fuseau) arrête la courbe.
\end{enumerate}
\textbf{Réflexe} : toute interprétation doit citer une donnée chiffrée du document (« le nombre de cellules passe de ... à ... en ... heures »).
\end{methode}

\begin{exoresolu}[Croissance d'une culture de cellules animales]
On met en culture $N_0 = 10^{4}$ cellules animales dans un milieu nutritif. La durée du cycle cellulaire est $T = \SI{20}{h}$ dans ces conditions. On compte les cellules à différents temps :
\begin{center}
\begin{tabular}{|l|c|c|c|c|c|}
\hline
\rowcolor{popblueL} Temps (h) & 0 & 20 & 40 & 60 & 80 \\
\hline Nombre de cellules & $10^{4}$ & $2 \times 10^{4}$ & $4 \times 10^{4}$ & $8 \times 10^{4}$ & $1{,}6 \times 10^{5}$ \\
\hline
\end{tabular}
\end{center}
\begin{enumerate}
\item Montrer que la population double à chaque cycle cellulaire.
\item Calculer le nombre théorique de cellules après 6 jours de culture si le rythme se maintient.
\item En pratique, on observe un ralentissement puis un plateau. Proposer une interprétation.
\end{enumerate}
\tcblower
\textbf{1. Doublement à chaque cycle.}

Comparons les valeurs successives :
\[ \frac{2 \times 10^{4}}{10^{4}} = 2 \ ; \quad \frac{4 \times 10^{4}}{2 \times 10^{4}} = 2 \ ; \quad \frac{8 \times 10^{4}}{4 \times 10^{4}} = 2 \ ; \quad \frac{1{,}6 \times 10^{5}}{8 \times 10^{4}} = 2. \]
Toutes les 20 h, c'est-à-dire à chaque cycle cellulaire, le nombre de cellules est multiplié par 2 : chaque cellule effectue une mitose donnant deux cellules filles.

\textbf{2. Nombre théorique après 6 jours.}

6 jours $= 6 \times 24 = \SI{144}{h}$. Le nombre de cycles écoulés est :
\[ k = \frac{t}{T} = \frac{144}{20} = 7{,}2 \text{ cycles.} \]
En appliquant la relation $N = N_0 \times 2^{k}$ :
\[ N = 10^{4} \times 2^{7{,}2} \approx 10^{4} \times 147 \approx 1{,}47 \times 10^{6} \text{ cellules.} \]

\textbf{3. Interprétation du plateau.}

Le ralentissement puis le plateau s'expliquent par la \textbf{limitation des ressources} : épuisement des nutriments du milieu, accumulation de déchets toxiques, et pour les cellules animales adhérentes, \cle{inhibition par contact} (les cellules cessent de se diviser lorsqu'elles recouvrent toute la surface du récipient). Le cycle cellulaire est alors bloqué en phase G1.

\textbf{Conclusion} : la population double à chaque cycle de 20 h (croissance exponentielle), ce qui conduirait théoriquement à environ $1{,}5 \times 10^{6}$ cellules en 6 jours ; en réalité, l'épuisement du milieu et l'inhibition par contact limitent la prolifération, ce qui montre que la division cellulaire est un processus étroitement contrôlé.
\end{exoresolu}

\courssub{Méthodes transversales et savoir-faire graphiques}

\courssub{Méthode 1 : Identifier les phases de la mitose à partir d'une photographie ou d'un schéma}

\begin{methode}[Reconnaître les quatre phases de la mitose]
Pour identifier la phase de la mitose représentée sur un document (photographie de cellules végétales en pointe de racine, schéma), procéder par étapes :
\begin{enumerate}
\item \textbf{Observer l'état des chromosomes} : sont-ils individualisés et condensés (bâtonnets visibles) ou décondensés (chromatine filamenteuse, noyau uniforme) ?
\item \textbf{Observer la position des chromosomes} : dispersés dans la cellule, alignés sur un plan équatorial, ou migrés vers les deux pôles ?
\item \textbf{Observer l'enveloppe nucléaire} : présente, en disparition, ou en reconstitution ?
\item \textbf{Observer la membrane cytoplasmique} : y a-t-il un étranglement (cellule animale) ou une plaque cellulaire (cellule végétale) ?
\item \textbf{Conclure} en appliquant les critères récapitulés :
\end{enumerate}
\begin{center}
\begin{tabularx}{\linewidth}{|l|X|}
\hline
\rowcolor{popblueL} \textbf{Phase} & \textbf{Critères d'identification} \\
\hline
Prophase & Chromosomes condensés (2 chromatides), enveloppe nucléaire qui disparaît, fuseau en formation \\
\hline
Métaphase & Chromosomes alignés sur la plaque équatoriale, fuseau achromatique complet \\
\hline
Anaphase & Séparation des chromatides sœurs, migration vers les pôles opposés (figures en V) \\
\hline
Télophase & Chromosomes aux pôles, décondensation, reconstitution des enveloppes nucléaires, cytodiérèse \\
\hline
\end{tabularx}
\end{center}
\textbf{Réflexe} : si le document montre une cellule à noyau intact et chromatine diffuse, il s'agit d'une \cle{interphase}, pas d'une phase de la mitose — ne pas la confondre avec la prophase.
\end{methode}

\begin{exoresolu}[Identification de phases sur des microphotographies]
On réalise une coupe de pointe de racine d'oignon (\emph{Allium cepa}, $2n = 16$) colorée au carmin acétique. Trois cellules A, B et C sont photographiées. Cellule A : les chromosomes, formés chacun de deux chromatides, sont alignés dans le plan équatorial de la cellule. Cellule B : des groupes de chromosomes en forme de V migrent vers les deux pôles de la cellule. Cellule C : les chromosomes condensés sont dispersés dans le cytoplasme et l'enveloppe nucléaire n'est plus visible.
\begin{enumerate}
\item Identifier la phase de la mitose représentée par chaque cellule, en justifiant.
\item Indiquer le nombre de chromosomes et de chromatides présents dans la cellule A.
\end{enumerate}
\tcblower
\textbf{1. Identification des phases.}

\textbf{Cellule A} : les chromosomes à deux chromatides sont \textbf{alignés sur la plaque équatoriale} : c'est le critère caractéristique de la \cle{métaphase}. Le fuseau achromatique est achevé et chaque chromosome est fixé par ses fibres kinétochoriennes.

\textbf{Cellule B} : les chromatides sœurs se sont séparées et \textbf{migrent vers les pôles opposés} en adoptant une forme de V (le centromère en tête) : c'est l'\cle{anaphase}. La migration résulte du raccourcissement des fibres du fuseau.

\textbf{Cellule C} : les chromosomes sont condensés mais encore dispersés, et \textbf{l'enveloppe nucléaire a disparu} : c'est la \cle{prophase} (fin de prophase). La condensation des chromosomes et la disparition de l'enveloppe nucléaire en sont les critères.

\textbf{2. Nombre de chromosomes et de chromatides en métaphase.}

En métaphase, la cellule n'a pas encore été divisée : elle contient le stock complet de chromosomes de l'espèce, soit $2n = 16$ chromosomes. Chaque chromosome est formé de \textbf{deux chromatides sœurs} (l'ADN a été répliqué en phase S de l'interphase), soit :
\[ 16 \times 2 = 32 \text{ chromatides.} \]

\textbf{Conclusion} : la cellule A est en métaphase et contient 16 chromosomes à deux chromatides, soit 32 chromatides ; la cellule B est en anaphase ; la cellule C est en prophase.
\end{exoresolu}

\courssub{Méthode 2 : Construire et lire un schéma des phases de la mitose}

\begin{methode}[Schématiser correctement la mitose (savoir-faire graphique exigible)]
Pour réaliser un schéma légendé d'une phase de la mitose, exigible dans les compositions :
\begin{enumerate}
\item \textbf{Choisir un nombre de chromosomes faible} : représenter $2n = 4$ (deux paires, distinguées par la taille ou la couleur) pour rester lisible.
\item \textbf{Respecter la cohérence entre les étapes} : mêmes chromosomes reconnaissables de la prophase à la télophase ; en anaphase, les chromatides sœurs migrées vers chaque pôle doivent être \textbf{en nombre égal et de mêmes types}.
\item \textbf{Représenter les structures indispensables} : fuseau achromatique (fibres reliant les pôles aux centromères), centrioles aux pôles (cellule animale), enveloppe nucléaire (présente ou non selon la phase).
\item \textbf{Légender chaque élément} par une ligne fine reliant la structure à son nom, sans croisement de lignes de légende.
\item \textbf{Titrer le schéma} : « Cellule animale en métaphase ($2n = 4$) » par exemple.
\end{enumerate}
\textbf{Réflexe} : un schéma sans légende ni titre ne rapporte qu'une fraction des points ; la rigueur de la légende est aussi importante que le dessin.
\end{methode}

\begin{exoresolu}[Analyse critique d'un schéma de mitose]
Un élève schématise une cellule animale ($2n = 4$) en anaphase : il dessine deux chromosomes à deux chromatides migrant vers un pôle et deux chromosomes à une chromatide migrant vers l'autre pôle.
\begin{enumerate}
\item Relever les erreurs commises par cet élève.
\item Décrire ce que devrait montrer un schéma correct de l'anaphase pour cette espèce.
\end{enumerate}
\tcblower
\textbf{1. Erreurs commises.}

\begin{itemize}
\item \textbf{Erreur 1} : en anaphase, les chromosomes migrants sont nécessairement à \textbf{une seule chromatide}, car l'anaphase est précisément définie par la séparation des chromatides sœurs. Dessiner des chromosomes à deux chromatides en migration est une contradiction avec la définition de la phase.
\item \textbf{Erreur 2} : la répartition est \textbf{inéquitable} : un pôle reçoit des chromosomes à deux chromatides et l'autre des chromosomes à une chromatide. Or la mitose répartit de façon strictement symétrique : chaque pôle doit recevoir \textbf{un exemplaire de chacun des 4 chromosomes} (soit 4 chromosomes à une chromatide par pôle).
\end{itemize}

\textbf{2. Schéma correct attendu.}

Un schéma correct de l'anaphase pour $2n = 4$ doit montrer :
\begin{itemize}
\item les \textbf{chromatides sœurs séparées}, devenues des chromosomes à une chromatide, en forme de V (centromère en tête, tiré par les fibres du fuseau qui raccourcissent) ;
\item \textbf{deux lots identiques de 4 chromosomes} migrant vers les deux pôles opposés (les deux paires de chromosomes, reconnaissables par leur taille, représentées dans chaque lot) ;
\item le \textbf{fuseau achromatique} et les centrioles aux pôles ;
\item une légende complète : chromosomes à une chromatide, centromère, fibres du fuseau, pôle cellulaire, centriole.
\end{itemize}

\textbf{Conclusion} : en anaphase, chaque pôle reçoit 4 chromosomes à une chromatide, génétiquement identiques ; cette répartition équitable est le mécanisme qui garantit que chaque cellule fille recevra le caryotype complet $2n = 4$.
\end{exoresolu}

\begin{attention}[Les 5 erreurs qui coûtent des points au Probatoire]
\begin{enumerate}
\item \textbf{Confondre chromosome et chromatide.} Après la phase S, un chromosome est formé de deux chromatides mais reste \emph{un} chromosome (un seul centromère). En métaphase, une cellule humaine possède 46 chromosomes et 92 chromatides — jamais « 92 chromosomes ».
\item \textbf{Affirmer que le nombre de chromosomes double pendant la phase S.} C'est la \textbf{quantité d'ADN} qui double ($Q \rightarrow 2Q$), pas le nombre de chromosomes, qui reste $2n$ jusqu'à l'anaphase.
\item \textbf{Identifier une interphase comme une prophase.} Un noyau intact à chromatine diffuse est une cellule en interphase ; la prophase se reconnaît à la \emph{condensation} des chromosomes et à la disparition de l'enveloppe nucléaire. Toujours justifier l'identification d'une phase par un critère visible sur le document.
\item \textbf{Oublier de citer les données chiffrées dans l'interprétation.} Une phrase comme « le nombre de cellules augmente » est insuffisante : il faut écrire « le nombre de cellules passe de $10^{4}$ à $8 \times 10^{4}$ en 60 h, soit un doublement toutes les 20 h ».
\item \textbf{Présenter le cancer comme une maladie contagieuse ou due à « trop de mitoses normales ».} La cancérisation résulte de \textbf{mutations} des gènes de contrôle du cycle cellulaire, entraînant des mitoses \emph{incontrôlées} ; la mitose elle-même reste un mécanisme normal et indispensable à la croissance et au renouvellement des tissus.
\end{enumerate}
\end{attention}

\courssub{Exercices d'entraînement (non résolus)}

\begin{exercice}[Application : Cycle cellulaire et quantités d'ADN]
Une espèce végétale a un caryotype $2n = 24$. Une cellule de cette espèce en phase G1 contient une quantité d'ADN $Q = \SI{5,0}{pg}$.
\begin{enumerate}
\item Compléter le tableau suivant pour une cellule de cette espèce :
\begin{center}
\begin{tabular}{|c|c|c|c|}
\hline
\rowcolor{popblueL} Stade & Nombre de chromosomes & Nombre de chromatides & Quantité d'ADN (pg) \\
\hline
G1 & & & \\
\hline
G2 & & & \\
\hline
Métaphase & & & \\
\hline
Anaphase (cellule entière) & & & \\
\hline
Cellule fille (après cytokinèse) & & & \\
\hline
\end{tabular}
\end{center}
\item Expliquer pourquoi le nombre de chromosomes est identique en G1 et en métaphase, alors que la quantité d'ADN a doublé.
\end{enumerate}
\end{exercice}

\begin{exercice}[Analyse de document : Effet de la colchicine]
La colchicine est un alcaloïde qui se lie à la tubuline et empêche la polymérisation des microtubules. On traite une culture de cellules de pointe de racine d'oignon avec de la colchicine pendant 24 h, puis on prépare des lames pour observer les chromosomes.
\begin{enumerate}
\item Prédire l'effet de la colchicine sur le fuseau achromatique et sur la progression de la mitose.
\item Quel stade de la mitose sera le plus fréquent dans la préparation ? Justifier.
\item Cette technique est utilisée en cytogénétique pour obtenir des caryotypes. Expliquer pourquoi.
\end{enumerate}
\end{exercice}

\begin{exercice}[Interprétation de courbe : Inhibition par contact]
On met en culture deux lignées cellulaires : des fibroblastes normaux (ligne continue) et des fibroblastes transformés (cancéreux, ligne pointillée). La courbe ci-dessous montre l'évolution du nombre de cellules en fonction du temps.
\begin{center}
\begin{tikzpicture}
\begin{axis}[
    width=\linewidth, height=6cm,
    xlabel={Temps (jours)}, ylabel={Nombre de cellules (échelle log)},
    xmin=0, xmax=10, ymin=3, ymax=7,
    ymode=log, log basis y=10,
    grid=major, legend pos=north west,
    tick label style={font=\small}, label style={font=\small},
]
\addplot[popblue, thick, domain=0:10, samples=50] {3 + 0.3*x - 0.015*x^2}; % Normal: rises then plateaus
\addlegendentry{Fibroblastes normaux}
\addplot[poporange, thick, dashed, domain=0:10, samples=50] {3 + 0.35*x}; % Cancer: exponential
\addlegendentry{Fibroblastes transformés}
\end{axis}
\end{tikzpicture}
\end{center}
\begin{enumerate}
\item Décrire l'évolution des deux populations.
\item Expliquer la différence de comportement au-delà du 5\textsuperscript{e} jour en lien avec les propriétés des cellules cancéreuses.
\item Que se passerait-il si l'on ajoutait du sérum frais au 8\textsuperscript{e} jour pour la culture normale ?
\end{enumerate}
\end{exercice}

\begin{corrige}[Exercice 1]
\begin{enumerate}
\item
\begin{center}
\begin{tabular}{|c|c|c|c|}
\hline
\rowcolor{popblueL} Stade & Nombre de chromosomes & Nombre de chromatides & Quantité d'ADN (pg) \\
\hline
G1 & 24 & 24 & 5,0 \\
\hline
G2 & 24 & 48 & 10,0 \\
\hline
Métaphase & 24 & 48 & 10,0 \\
\hline
Anaphase (cellule entière) & 48 & 48 & 10,0 \\
\hline
Cellule fille & 24 & 24 & 5,0 \\
\hline
\end{tabular}
\end{center}
\item En phase S, chaque chromosome (1 molécule d'ADN) est répliqué en deux chromatides sœurs (2 molécules d'ADN identiques) jointes par un centromère unique. Le comptage des chromosomes se fait par le nombre de centromères, qui ne change pas avant l'anaphase. La quantité d'ADN, elle, double car il y a deux fois plus de molécules d'ADN.
\end{enumerate}
\end{corrige}

\begin{corrige}[Exercice 2]
\begin{enumerate}
\item La colchicine dépolymerise les microtubules $\rightarrow$ pas de fuseau achromatique $\rightarrow$ les chromosomes ne peuvent pas s'aligner (métaphase) ni migrer (anaphase). Le cycle est bloqué en \textbf{métaphase} (point de contrôle M non validé).
\item La \textbf{métaphase} sera surreprésentée. Les cellules s'accumulent à ce stade car elles ne peuvent pas passer en anaphase.
\item En bloquant en métaphase, les chromosomes sont maximalement condensés et individualisés, ce qui facilite leur observation, leur comptage, leur photographie et leur classement par paires homologues pour établir le caryotype.
\end{enumerate}
\end{corrige}

\begin{corrige}[Exercice 3]
\begin{enumerate}
\item \textbf{Normaux :} Phase de latence (0--1 j), croissance exponentielle (1--5 j), puis \textbf{plateau} à partir du 5\textsuperscript{e} jour (confluence, inhibition par contact). \textbf{Transformés :} Croissance exponentielle soutenue sans plateau visible sur la période (pas d'inhibition par contact, perte du contrôle de densité).
\item Les cellules cancéreuses ont perdu les points de contrôle (mutations p53, Rb, etc.) et l'inhibition par contact : elles continuent de proliférer au-delà de la confluence, s'empilent (foyers), et peuvent croître en suspension (ancrage-indépendant).
\item L'ajout de sérum frais (facteurs de croissance) pourrait stimuler une reprise de la prolifération des cellules normales (réentrée en cycle depuis G0/G1), mais le plateau est aussi dû à l'épuisement du substrat et à l'inhibition par contact physique ; la reprise serait partielle et temporaire.
\end{enumerate}
\end{corrige}

\newpage

\coursec{Exercices}

\courssub{Je vérifie mes connaissances}

\begin{exercice}[QCM : le cycle cellulaire]
Pour chaque affirmation, choisir la (ou les) bonne(s) réponse(s) et justifier brièvement.
\begin{enumerate}
\item Au cours de l'interphase, la cellule :
\begin{enumerate}
\item[a)] se divise en deux cellules filles ;
\item[b)] réplique son ADN pendant la phase S ;
\item[c)] voit ses chromosomes se condenser au maximum ;
\item[d)] réalise la majeure partie de sa croissance et de son activité métabolique.
\end{enumerate}
\item La réplication de l'ADN est dite semi-conservative parce que :
\begin{enumerate}
\item[a)] seule la moitié des molécules d'ADN est répliquée ;
\item[b)] chaque molécule fille conserve un brin parental et un brin nouvellement synthétisé ;
\item[c)] les deux brins parentaux restent toujours associés ;
\item[d)] la réplication ne concerne que la moitié du génome.
\end{enumerate}
\item Au cours de l'anaphase :
\begin{enumerate}
\item[a)] les chromosomes s'alignent sur la plaque équatoriale ;
\item[b)] l'enveloppe nucléaire se reforme ;
\item[c)] les chromatides sœurs se séparent et migrent vers les pôles opposés ;
\item[d)] les chromosomes se décondensent.
\end{enumerate}
\item La cancérisation résulte :
\begin{enumerate}
\item[a)] d'un arrêt définitif de toute division cellulaire ;
\item[b)] d'une perte du contrôle du cycle cellulaire entraînant une prolifération anarchique ;
\item[c)] d'une mutation toujours héréditaire ;
\item[d)] exclusivement d'une infection virale.
\end{enumerate}
\end{enumerate}
\end{exercice}

\begin{exercice}[Vrai ou faux, à justifier]
Répondre par VRAI ou FAUX, puis justifier chaque réponse en une ou deux phrases.
\begin{enumerate}
\item L'interphase est une phase de repos pendant laquelle il ne se passe rien dans la cellule.
\item À l'issue de la phase S, chaque chromosome est formé de deux chromatides sœurs identiques.
\item C'est au cours de la métaphase que les chromatides sœurs se séparent.
\item La mitose produit deux cellules filles génétiquement identiques à la cellule mère.
\item Un tissu cancéreux présente un indice mitotique généralement plus faible qu'un tissu sain.
\end{enumerate}
\end{exercice}

\begin{exercice}[Définitions essentielles]
Définir de manière précise et concise les termes suivants : \cle{cycle cellulaire} ; \cle{mitose} ; \cle{chromatide sœur} ; \cle{caryotype} ; \cle{réplication semi-conservative} ; \cle{indice mitotique}.
\end{exercice}

\begin{exercice}[Calculs directs]
\begin{enumerate}
\item Une cellule animale possède $2n = 12$ chromosomes et une quantité d'ADN égale à $Q$ en début de phase G1. Donner, en fin de phase G2, le nombre de chromosomes, le nombre de chromatides et la quantité d'ADN (en fonction de $Q$).
\item Sur une coupe de racine d'oignon, on observe \num{1000} cellules dont $45$ en mitose. Calculer l'indice mitotique de ce tissu.
\item Une culture cellulaire compte \num{20000} cellules à l'instant initial. Sachant que toutes les cellules se divisent en même temps et que la durée du cycle est de \SI{24}{h}, combien de cellules obtient-on après \SI{72}{h} ?
\end{enumerate}
\end{exercice}

\courssub{Je m'entraîne}

\begin{exercice}[Identifier les phases de la mitose]
Quatre photographies de cellules de méristème de racine d'oignon sont décrites ci-dessous.
\begin{itemize}
\item \textbf{Photo A} : les chromosomes, formés chacun de deux chromatides, sont alignés dans le plan équatorial de la cellule ; le fuseau est visible.
\item \textbf{Photo B} : la chromatine se condense en chromosomes individualisés ; l'enveloppe nucléaire disparaît progressivement.
\item \textbf{Photo C} : deux groupes de chromosomes à une chromatide migrent vers les pôles opposés de la cellule.
\item \textbf{Photo D} : deux nouveaux noyaux se reforment aux pôles ; les chromosomes se décondensent et la cellule se divise en deux.
\end{itemize}
\begin{enumerate}
\item Identifier la phase de la mitose représentée par chaque photographie, en justifiant la réponse à partir d'un critère observable.
\item Remettre les quatre photographies dans l'ordre chronologique de la mitose.
\item Citer l'événement du cycle cellulaire qui se déroule avant la photo B et qui est indispensable à la conservation de l'information génétique.
\end{enumerate}
\end{exercice}

\begin{exercice}[Quantité d'ADN au cours du cycle]
On mesure la quantité d'ADN par noyau dans une culture de cellules animales en division. Les résultats sont consignés dans le tableau ci-dessous.
\begin{center}
\begin{tabularx}{\linewidth}{|l|X|X|X|X|X|X|X|X|X|}
\hline
\rowcolor{popblueL} Temps (h) & 0 & 2 & 4 & 6 & 8 & 10 & 12 & 14 & 16 \\
\hline
Quantité d'ADN (u.a.) & 6 & 6 & 6 & 7,5 & 9 & 10,5 & 12 & 12 & 6 \\
\hline
\end{tabularx}
\end{center}
\begin{enumerate}
\item Tracer la courbe de variation de la quantité d'ADN en fonction du temps (échelle : \SI{1}{cm} pour \SI{2}{h} et \SI{1}{cm} pour 2 u.a.).
\item Repérer sur la courbe les phases G1, S, G2 et M du cycle cellulaire. Justifier les limites choisies.
\item Déterminer graphiquement la durée de la phase S.
\item Expliquer, en une phrase, pourquoi la quantité d'ADN double pendant la phase S alors que le nombre de chromosomes reste inchangé.
\end{enumerate}
\end{exercice}

\begin{exercice}[Une cellule à $2n = 8$ chromosomes]
Une cellule mère possède $2n = 8$ chromosomes. On note $Q$ la quantité d'ADN de son noyau en début de phase G1.
\begin{enumerate}
\item Recopier et compléter le tableau suivant, donnant pour chaque étape le nombre de chromosomes, le nombre de chromatides et la quantité d'ADN (en fonction de $Q$).
\begin{center}
\begin{tabularx}{\linewidth}{|l|X|X|X|}
\hline
\rowcolor{popblueL} Étape & Chromosomes & Chromatides & ADN \\
\hline
Fin de G1 & & & \\
\hline
Fin de G2 & & & \\
\hline
Métaphase & & & \\
\hline
Anaphase (dans toute la cellule) & & & \\
\hline
Chaque cellule fille & & & \\
\hline
\end{tabularx}
\end{center}
\item Comparer le caryotype des cellules filles à celui de la cellule mère. Conclure sur le rôle de la mitose dans la conservation de l'identité biologique.
\item Expliquer pourquoi la réplication de l'ADN pendant l'interphase est une condition nécessaire à cette conservation.
\end{enumerate}
\end{exercice}

\begin{exercice}[Indice mitotique et activité proliférative]
Au laboratoire d'anatomopathologie du Centre Hospitalier d'Essos à Yaoundé, un technicien examine deux coupes de tissus.
\begin{itemize}
\item \textbf{Coupe 1 (tissu sain)} : $800$ cellules observées, dont $32$ en mitose.
\item \textbf{Coupe 2 (tissu tumoral)} : indice mitotique mesuré égal à $12\,\%$.
\end{itemize}
\begin{enumerate}
\item Rappeler la définition de l'indice mitotique.
\item Calculer l'indice mitotique de la coupe 1.
\item Comparer les deux indices et conclure sur l'activité proliférative de chaque tissu.
\item Expliquer, en deux ou trois phrases, le lien entre un indice mitotique élevé et le phénomène de cancérisation.
\end{enumerate}
\end{exercice}

\courssub{Je me prépare au Probatoire}

\begin{exercice}[Type Probatoire : la mitose dans le méristème d'oignon]
Au cours d'une séance de travaux pratiques au lycée, un groupe d'élèves de Première D réalise une préparation microscopique de pointes de racines d'oignon (\emph{Allium cepa}, $2n = 16$) préalablement placées dans l'eau pendant quelques jours. Après coloration, ils observent quatre cellules notées 1, 2, 3 et 4, décrites ci-dessous.
\begin{itemize}
\item \textbf{Cellule 1} : les chromosomes, très condensés et formés de deux chromatides, sont dispersés dans le cytoplasme ; on ne distingue plus d'enveloppe nucléaire.
\item \textbf{Cellule 2} : deux lots de chromosomes à une chromatide se rapprochent des pôles ; la distance entre les deux lots augmente.
\item \textbf{Cellule 3} : tous les chromosomes à deux chromatides sont disposés sur une même plaque, à égale distance des deux pôles.
\item \textbf{Cellule 4} : aux deux pôles, la chromatine se décondense ; un sillon de partage apparait au milieu de la cellule.
\end{itemize}
\begin{enumerate}
\item Nommer la phase de la mitose illustrée par chacune des cellules 1, 2, 3 et 4, en justifiant chaque identification par un argument tiré de la description. \hfill (4 pts)
\item Classer ces quatre cellules dans l'ordre chronologique de la mitose. \hfill (1 pt)
\item Pour la cellule 3, indiquer le nombre de chromosomes, le nombre de chromatides et la quantité d'ADN (on notera $Q$ la quantité d'ADN d'un noyau en G1). \hfill (1,5 pt)
\item Schématiser la cellule 2 en y représentant quatre chromosomes et le fuseau, avec une légende correcte. \hfill (2 pts)
\item Expliquer en quelques lignes comment le comportement des chromosomes au cours de la mitose garantit que chaque cellule fille reçoit un caryotype identique à celui de la cellule mère. \hfill (1,5 pt)
\end{enumerate}
\end{exercice}

\begin{exercice}[Type Probatoire : suivi d'une culture cellulaire]
Des chercheurs de l'Université de Yaoundé I étudient la prolifération de cellules en culture. Ils mesurent la quantité d'ADN par noyau (en unités arbitraires) à différents instants après la mise en culture synchronisée des cellules. Résultats :
\begin{center}
\begin{tabularx}{\linewidth}{|l|X|X|X|X|X|X|X|X|X|X|X|}
\hline
\rowcolor{popblueL} Temps (h) & 0 & 2 & 4 & 6 & 8 & 10 & 12 & 14 & 16 & 18 & 20 \\
\hline
Quantité d'ADN (u.a.) & 5 & 5 & 5 & 5,8 & 7 & 8,2 & 9,4 & 10 & 10 & 10 & 5 \\
\hline
\end{tabularx}
\end{center}
\begin{enumerate}
\item Tracer la courbe de variation de la quantité d'ADN par noyau en fonction du temps. \hfill (2 pts)
\item Délimiter sur le graphique les phases G1, S, G2 et M du cycle cellulaire, en justifiant chaque délimitation. \hfill (2 pts)
\item Calculer la durée de la phase S et la durée totale du cycle cellulaire de ces cellules. \hfill (1 pt)
\item Sachant que la culture comptait \num{50000} cellules à l'instant initial et que toutes les cellules se divisent à chaque cycle, calculer l'effectif théorique de la culture après \SI{60}{h}. \hfill (1,5 pt)
\item En réalité, on ne compte que \num{300000} cellules au bout de \SI{60}{h}. Proposer deux hypothèses expliquant cet écart entre l'effectif théorique et l'effectif réel. \hfill (1,5 pt)
\item Expliquer pourquoi la réplication de l'ADN pendant la phase S est indispensable au maintien de l'identité biologique des cellules au fil des générations. \hfill (2 pts)
\end{enumerate}
\end{exercice}

\begin{exercice}[Type Probatoire : tâche complexe — un extrait de plante contre la prolifération cellulaire]
La médecine traditionnelle camerounaise utilise depuis longtemps des extraits de certaines plantes pour traiter des tumeurs. Une équipe de recherche teste l'effet d'un extrait de plante locale sur des cellules cancéreuses en culture. Deux cultures identiques sont réalisées : une culture témoin (sans extrait) et une culture traitée (avec extrait). Le nombre de cellules est compté au cours du temps :
\begin{center}
\begin{tabularx}{\linewidth}{|l|X|X|X|X|X|}
\hline
\rowcolor{popblueL} Temps (h) & 0 & 24 & 48 & 72 & 96 \\
\hline
Culture témoin (milliers de cellules) & 50 & 100 & 200 & 400 & 800 \\
\hline
Culture traitée (milliers de cellules) & 50 & 62 & 70 & 72 & 71 \\
\hline
\end{tabularx}
\end{center}
Par ailleurs, l'examen microscopique de la culture traitée à $t = \SI{48}{h}$ montre de nombreuses cellules bloquées avec des chromosomes alignés sur la plaque équatoriale et un fuseau mitotique mal formé.
\begin{enumerate}
\item Tracer sur un même graphique les courbes d'évolution du nombre de cellules pour les deux cultures. \hfill (2 pts)
\item Comparer les deux courbes et dégager l'effet de l'extrait sur la prolifération cellulaire. \hfill (2 pts)
\item Calculer l'indice de multiplication (rapport effectif final / effectif initial) de chaque culture entre 0 et \SI{96}{h}. \hfill (1 pt)
\item À partir des observations microscopiques, identifier la phase de la mitose bloquée par l'extrait et formuler une hypothèse précise sur son mécanisme d'action au niveau cellulaire. \hfill (2 pts)
\item Proposer une expérience complémentaire permettant de tester cette hypothèse. \hfill (1,5 pt)
\item En t'appuyant sur tes connaissances sur la cancérisation, expliquer l'intérêt médical potentiel d'un tel extrait et préciser une précaution indispensable avant toute application thérapeutique. \hfill (1,5 pt)
\end{enumerate}
\end{exercice}

\newpage

\coursec{Corrigés des exercices}

\begin{corrige}[Exercice 1 — QCM : le cycle cellulaire]
\textbf{Question 1 :} Les bonnes réponses sont \textbf{b)} et \textbf{d)}.
\begin{itemize}
\item \textbf{a)} Faux : la division cellulaire (mitose + cytocinèse) a lieu pendant la phase M, pas pendant l'interphase.
\item \textbf{b)} Vrai : la phase S (synthèse) de l'interphase est consacrée à la réplication de l'ADN.
\item \textbf{c)} Faux : la condensation maximale des chromosomes se produit en prophase (phase M), pas en interphase (où la chromatine est décondensée).
\item \textbf{d)} Vrai : l'interphase (G1, S, G2) est la période où la cellule croît, assure ses fonctions métaboliques et réplique son ADN.
\end{itemize}

\textbf{Question 2 :} La bonne réponse est \textbf{b)}.
\begin{itemize}
\item \textbf{a)} Faux : tout le génome est répliqué.
\item \textbf{b)} Vrai : c'est la définition même de la réplication semi-conservative (Meselson et Stahl, 1958) : chaque double brin fille conserve un brin parental (matrice) et un brin néo-synthétisé.
\item \textbf{c)} Faux : les deux brins parentaux se séparent pour servir de matrices.
\item \textbf{d)} Faux : la réplication concerne l'intégralité du génome.
\end{itemize}

\textbf{Question 3 :} La bonne réponse est \textbf{c)}.
\begin{itemize}
\item \textbf{a)} Faux : c'est l'événement caractéristique de la métaphase.
\item \textbf{b)} Faux : la réformation de l'enveloppe nucléaire marque le début de la télophase.
\item \textbf{c)} Vrai : en anaphase, les centromères se divisent, les chromatides sœurs deviennent des chromosomes filles et migrent vers les pôles sous l'action du fuseau.
\item \textbf{d)} Faux : la décondensation a lieu en télophase.
\end{itemize}

\textbf{Question 4 :} La bonne réponse est \textbf{b)}.
\begin{itemize}
\item \textbf{a)} Faux : la cancérisation se traduit par une prolifération excessive, pas un arrêt.
\item \textbf{b)} Vrai : la cancérisation résulte de l'accumulation de mutations affectant les gènes de contrôle du cycle cellulaire (oncogènes, gènes suppresseurs de tumeurs), entraînant une division anarchique.
\item \textbf{c)} Faux : la grande majorité des cancers sont sporadiques (mutations somatiques) ; seuls 5 à 10 \% sont héréditaires (prédisposition génétique).
\item \textbf{d)} Faux : certains virus sont oncogènes (HPV, HBV, EBV), mais la cancérisation a des causes multiples (chimiques, radiations, erreurs de réplication spontanées).
\end{itemize}
\end{corrige}

\begin{corrige}[Exercice 2 — Vrai ou faux, à justifier]
\begin{enumerate}
\item \textbf{FAUX.} L'interphase n'est pas une phase de repos : la cellule y assure sa croissance (G1), réplique son ADN (S), prépare la division (G2) et maintient une intense activité métabolique (transcription, traduction, respiration cellulaire).
\item \textbf{VRAI.} À l'issue de la phase S, chaque chromosome est constitué de deux chromatides sœurs identiques (issues de la réplication semi-conservative), jointes au niveau du centromère. Le nombre de chromosomes (2n) ne change pas avant l'anaphase.
\item \textbf{FAUX.} C'est au cours de l'\textbf{anaphase} que les chromatides sœurs se séparent. En métaphase, les chromosomes (à deux chromatides) s'alignent sur la plaque équatoriale.
\item \textbf{VRAI.} La mitose est une division équationnelle : elle distribue à chaque cellule fille un jeu chromosomique identique à celui de la cellule mère (même nombre, même forme, même information génétique), assurant la stabilité du caryotype.
\item \textbf{FAUX.} Un tissu cancéreux se caractérise par une perte de contrôle du cycle cellulaire, ce qui se traduit par un \textbf{indice mitotique élevé} (souvent > 5-10 \%), bien supérieur à celui des tissus sains adultes (généralement < 1-2 \%).
\end{enumerate}
\end{corrige}

\begin{corrige}[Exercice 3 — Définitions essentielles]
\textbf{\cle{Cycle cellulaire}} : Suite ordonnée et cyclique des événements qui conduit une cellule, depuis sa naissance par division d'une cellule mère, jusqu'à sa propre division en deux cellules filles. Il comprend l'interphase (G1, S, G2) et la phase mitotique (M).

\textbf{\cle{Mitose}} : Processus de division du noyau au cours duquel les chromosomes condensés se partagent équitablement entre deux noyaux fils, garantissant la conservation du caryotype (nombre et structure des chromosomes) d'une génération cellulaire à l'autre. Elle comprend la prophase, la métaphase, l'anaphase et la télophase.

\textbf{\cle{Chromatide sœur}} : Chacun des deux brins d'ADN identiques résultant de la réplication d'un chromosome unique lors de la phase S. Les deux chromatides sœurs sont jointes par leur centromère et se séparent lors de l'anaphase pour devenir des chromosomes filles indépendants.

\textbf{\cle{Caryotype}} : Ensemble des chromosomes d'une cellule ou d'un individu, caractérisé par leur nombre (formule chromosomique, ex. $2n=46$ chez l'homme), leur morphologie (taille, position du centromère, bandage) et leur disposition par paires homologues.

\textbf{\cle{Réplication semi-conservative}} : Mode de duplication de l'ADN au cours duquel chaque molécule fille conserve l'un des deux brins parentaux (brin matrice) et incorpore un brin complémentaire néo-synthétisé. Démontrée par Meselson et Stahl (1958).

\textbf{\cle{Indice mitotique}} : Pourcentage de cellules observées en cours de division (mitose) par rapport au nombre total de cellules comptées dans un échantillon tissulaire. Formule : $IM = \frac{N_{\text{mitose}}}{N_{\text{total}}} \times 100$. Il reflète l'activité proliférative du tissu.
\end{corrige}

\begin{corrige}[Exercice 4 — Calculs directs]
\begin{enumerate}
\item \textbf{Données :} $2n = 12$, ADN = $Q$ en début de G1 (chaque chromosome = 1 chromatide).
\begin{itemize}
\item Fin de G2 (après phase S) : l'ADN a été répliqué $\rightarrow$ quantité d'ADN = $2Q$.
\item Nombre de chromosomes : toujours $2n = 12$ (chaque chromosome est maintenant formé de 2 chromatides sœurs).
\item Nombre de chromatides : $12 \times 2 = 24$.
\end{itemize}
\textbf{Réponse :} 12 chromosomes ; 24 chromatides ; $2Q$ d'ADN.

\item \textbf{Indice mitotique} $IM = \frac{45}{1000} \times 100 = \mathbf{4,5\,\%}$.

\item \textbf{Durée du cycle} = 24 h. \textbf{Durée totale} = 72 h $\rightarrow$ nombre de cycles = $\frac{72}{24} = 3$.
Chaque cycle double l'effectif : $N = N_0 \times 2^n = 20\,000 \times 2^3 = 20\,000 \times 8 = \mathbf{160\,000 \text{ cellules}}$.
\end{enumerate}
\end{corrige}

\begin{corrige}[Exercice 5 — Identifier les phases de la mitose]
\begin{enumerate}
\item \textbf{Photo A :} \textbf{Métaphase}. Critère : chromosomes (à deux chromatides) alignés sur la plaque équatoriale, fuseau visible.
\item \textbf{Photo B :} \textbf{Prophase}. Critère : condensation de la chromatine en chromosomes individualisés, disparition progressive de l'enveloppe nucléaire.
\item \textbf{Photo C :} \textbf{Anaphase}. Critère : deux groupes de chromosomes (à une chromatide) migrant vers les pôles opposés.
\item \textbf{Photo D :} \textbf{Télophase} (début de cytocinèse). Critère : reformation de deux noyaux aux pôles, décondensation des chromosomes, apparition du sillon de division.

\item \textbf{Ordre chronologique :} B (Prophase) $\rightarrow$ A (Métaphase) $\rightarrow$ C (Anaphase) $\rightarrow$ D (Télophase).

\item L'événement indispensable qui précède la prophase (photo B) est la \textbf{réplication de l'ADN pendant la phase S de l'interphase}. Elle assure que chaque chromosome soit constitué de deux chromatides sœurs identiques, condition nécessaire à la distribution équitable du matériel génétique.
\end{enumerate}
\end{corrige}

\begin{corrige}[Exercice 6 — Quantité d'ADN au cours du cycle]
\begin{enumerate}
\item \textbf{Courbe} (voir figure ci-dessous).
\begin{popfigure}
\begin{tikzpicture}
\begin{axis}[
width=\linewidth,
height=8cm,
xlabel={Temps (h)},
ylabel={Quantité d'ADN (u.a.)},
xmin=0, xmax=16,
ymin=0, ymax=14,
xtick={0,2,4,6,8,10,12,14,16},
ytick={0,2,4,6,8,10,12,14},
grid=major,
grid style={popline},
tick label style={font=\small},
label style={font=\small},
axis line style={popdark},
]
\addplot[popcurve, very thick, mark=*, mark size=2pt, mark options={fill=popblue}] coordinates {
(0,6) (2,6) (4,6) (6,7.5) (8,9) (10,10.5) (12,12) (14,12) (16,6)
};
\end{axis}
\end{tikzpicture}
\legende{Variation de la quantité d'ADN par noyau au cours du cycle cellulaire.}
\end{popfigure}

\item \textbf{Repérage des phases :}
\begin{itemize}
\item \textbf{G1 :} de 0 à 4 h. Quantité d'ADN constante (6 u.a.) : pas de réplication, croissance cellulaire.
\item \textbf{S :} de 4 à 12 h. Augmentation continue et linéaire de 6 à 12 u.a. : réplication de l'ADN.
\item \textbf{G2 :} de 12 à 14 h. Quantité d'ADN constante (12 u.a.) : fin de réplication, préparation mitose.
\item \textbf{M :} de 14 à 16 h. Chute brutale de 12 à 6 u.a. : division nucléaire (mitose) et cytocinèse, répartition de l'ADN dans deux cellules filles.
\end{itemize}

\item \textbf{Durée de la phase S :} début à 4 h, fin à 12 h $\rightarrow$ \textbf{8 h}.

\item La quantité d'ADN double pendant la phase S car chaque chromosome se réplique pour former \textbf{deux chromatides sœurs identiques} ; le nombre de chromosomes (compté par le nombre de centromères) ne change pas tant que les chromatides sœurs ne se sont pas séparées (anaphase).
\end{enumerate}
\end{corrige}

\begin{corrige}[Exercice 7 — Une cellule à $2n = 8$ chromosomes]
\begin{enumerate}
\item \textbf{Tableau complété :}
\begin{center}
\begin{tabularx}{\linewidth}{|l|X|X|X|}
\hline
\rowcolor{popblueL} \textbf{Étape} & \textbf{Chromosomes} & \textbf{Chromatides} & \textbf{ADN} \\
\hline
Fin de G1 & 8 & 8 & $Q$ \\
\hline
Fin de G2 & 8 & 16 & $2Q$ \\
\hline
Métaphase & 8 & 16 & $2Q$ \\
\hline
Anaphase (dans toute la cellule) & 16 & 16 & $2Q$ \\
\hline
Chaque cellule fille & 8 & 8 & $Q$ \\
\hline
\end{tabularx}
\end{center}
\textit{Note :} En anaphase, la séparation des chromatides sœurs double transitoirement le nombre de chromosomes dans la cellule mère (16 chromosomes à une chromatide) avant la cytocinèse.

\item \textbf{Comparaison :} Les deux cellules filles présentent un caryotype \textbf{strictement identique} à celui de la cellule mère : même nombre de chromosomes ($2n=8$), même morphologie, même information génétique.
\textbf{Conclusion :} La mitose est une division \textbf{équationnelle} (conservatrice) qui assure la stabilité du patrimoine génétique (caryotype) d'une génération cellulaire à l'autre, fondement de l'identité biologique de l'organisme.

\item La réplication de l'ADN (phase S) est indispensable car elle \textbf{double la quantité d'information génétique} avant la division. Sans elle, chaque division réduirait de moitié la quantité d'ADN par noyau, entraînant une perte progressive du génome et l'impossibilité de transmettre l'information génétique complète aux cellules filles. La réplication semi-conservative garantit que chaque cellule fille reçoit une copie fidèle de chaque molécule d'ADN parentale.
\end{enumerate}
\end{corrige}

\begin{corrige}[Exercice 8 — Indice mitotique et activité proliférative]
\begin{enumerate}
\item \textbf{Définition :} L'indice mitotique (IM) est le pourcentage de cellules en mitose par rapport au nombre total de cellules observées dans un échantillon : $IM = \frac{\text{Nombre de cellules en mitose}}{\text{Nombre total de cellules observées}} \times 100$.

\item \textbf{Coupe 1 (tissu sain) :} $IM_1 = \frac{32}{800} \times 100 = \mathbf{4\,\%}$.

\item \textbf{Comparaison :} Coupe 1 (sain) : $IM = 4\,\%$ ; Coupe 2 (tumoral) : $IM = 12\,\%$.
L'indice mitotique du tissu tumoral est \textbf{trois fois plus élevé}. Cela traduit une \textbf{activité proliférative intense et anarchique} des cellules cancéreuses, alors que le tissu sain présente une prolifération contrôlée, limitée au renouvellement physiologique.

\item Un indice mitotique élevé témoigne d'une \textbf{perte des points de contrôle du cycle cellulaire} (checkpoints G1/S, G2/M, fuseau). Les cellules cancéreuses échappent aux signaux d'arrêt (absence de facteurs de croissance, inhibition par contact, dommages ADN) et entrent en division de manière continue et non régulée. Cette prolifération anarchique est la marque distinctive de la cancérisation.
\end{enumerate}
\end{corrige}

\begin{corrige}[Exercice 9 — Type Probatoire : la mitose dans le méristème d'oignon]
\begin{enumerate}
\item \textbf{Cellule 1 : Prophase.} Justification : chromosomes très condensés (à deux chromatides), dispersés dans le cytoplasme, enveloppe nucléaire invisible (disparue).
\item \textbf{Cellule 2 : Anaphase.} Justification : deux lots de chromosomes à une chromatide migrent vers les pôles opposés, distance inter-polaire croissante.
\item \textbf{Cellule 3 : Métaphase.} Justification : tous les chromosomes (à deux chromatides) alignés sur la plaque équatoriale (plan médian).
\item \textbf{Cellule 4 : Télophase} (avec cytocinèse). Justification : décondensation de la chromatine aux pôles (reformation des noyaux), apparition du sillon de partage.

\item \textbf{Ordre chronologique :} 1 (Prophase) $\rightarrow$ 3 (Métaphase) $\rightarrow$ 2 (Anaphase) $\rightarrow$ 4 (Télophase).

\item \textbf{Cellule 3 (Métaphase) :} $2n = 16$ chromosomes (chaque chromosome = 2 chromatides) $\rightarrow$ \textbf{32 chromatides}. Quantité d'ADN = \textbf{$2Q$} (réplication effectuée en phase S).

\item \textbf{Schéma de la cellule 2 (Anaphase) :}
\begin{popfigure}
\begin{tikzpicture}[scale=0.9, every node/.style={font=\small}]
% Cell outline
\draw[popdark, thick] (-3,-2) rectangle (3,2);
% Poles
\draw[popblue, thick] (-3,0) -- (-3.5,0) node[left] {Pôle};
\draw[popblue, thick] (3,0) -- (3.5,0) node[right] {Pôle};
% Spindle fibers
\foreach \y in {-1.2,-0.6,0,0.6,1.2} {
  \draw[poporange, dashed] (-3,\y) -- (-1.5,\y);
  \draw[poporange, dashed] (3,\y) -- (1.5,\y);
}
% Chromosomes (4 chromosomes = 4 pairs of chromatids separating)
% Actually in anaphase, sister chromatids have separated. 4 chromosomes total in haploid set? 2n=16, but question says "représentant quatre chromosomes". So draw 4 chromosomes (single chromatids) moving to each pole.
% Let's draw 4 chromosomes going left, 4 going right.
\foreach \i/\y in {1/-1.2, 2/-0.4, 3/0.4, 4/1.2} {
  % Left group
  \draw[poppurple, very thick] (-1.5,\y) -- (-2.5,\y);
  \node[poppurple] at (-2.8,\y) {Chr};
  % Right group
  \draw[poppurple, very thick] (1.5,\y) -- (2.5,\y);
  \node[poppurple] at (2.8,\y) {Chr};
}
% Centromeres/kinetochores
\foreach \y in {-1.2,-0.4,0.4,1.2} {
  \fill[poppink] (-2,\y) circle (2pt);
  \fill[poppink] (2,\y) circle (2pt);
}
% Legend
\node[align=left, anchor=north west, fill=popcream, draw=popdark, rounded corners, inner sep=4pt] at (-3.8,2.2) {
\textbf{Légende :}\\
\tikz\draw[popdark, thick] (0,0) -- (0.5,0); Membrane cellulaire\\
\tikz\draw[poporange, dashed] (0,0) -- (0.5,0); Fibres du fuseau (microtubules)\\
\tikz\draw[poppurple, very thick] (0,0) -- (0.5,0); Chromosome (1 chromatide)\\
\tikz\fill[poppink] (0,0) circle (2pt); Centromère / Kinétochore\\
\tikz\draw[popblue, thick] (0,0) -- (0.5,0); Pôle cellulaire (centrosome)
};
\end{tikzpicture}
\legende{Cellule en anaphase (2n=16 simplifié à 4 chromosomes pour la clarté) : chromosomes à une chromatide séparés migrant vers les pôles.}
\end{popfigure}

\item \textbf{Conservation du caryotype :} Au cours de la mitose, la réplication de l'ADN (phase S) produit deux chromatides sœurs \textbf{identiques} pour chaque chromosome. En métaphase, les chromosomes s'alignent sur la plaque équatoriale. En anaphase, les chromatides sœurs se séparent et \textbf{chacune migre vers un pôle opposé}. Ainsi, chaque pôle reçoit \textbf{un exemplaire de chaque chromosome} (donc un jeu complet de $2n$ chromosomes). Comme les chromatides sœurs sont des copies conformes, les deux noyaux fils reconstitués en télophase possèdent un caryotype \textbf{identique} à celui de la cellule mère (même nombre, même structure, mêmes allèles). C'est ce mécanisme précis de ségrégation équationnelle qui garantit le maintien de l'identité biologique.
\end{enumerate}

\textbf{Barème indicatif (10 pts) :}
\begin{itemize}
\item Q1 : 4 pts (1 pt par identification correcte + justification)
\item Q2 : 1 pt (ordre correct)
\item Q3 : 1,5 pt (0,5 pt par valeur : chr, chromatides, ADN)
\item Q4 : 2 pts (schéma correct + légende complète)
\item Q5 : 1,5 pt (explication claire : réplication $\rightarrow$ chromatides identiques $\rightarrow$ séparation équatoriale $\rightarrow$ distribution équitable)
\end{itemize}

\textbf{Points de vigilance :} Préciser « chromatides sœurs identiques » pour Q5. Pour le schéma, respecter la consigne « quatre chromosomes » (donc 4 vers chaque pôle) et légender fuseau, chromosomes, pôles.
\end{corrige}

\begin{corrige}[Exercice 10 — Type Probatoire : suivi d'une culture cellulaire]
\begin{enumerate}
\item \textbf{Courbe} (voir figure).
\begin{popfigure}
\begin{tikzpicture}
\begin{axis}[
width=\linewidth,
height=8cm,
xlabel={Temps (h)},
ylabel={Quantité d'ADN par noyau (u.a.)},
xmin=0, xmax=20,
ymin=0, ymax=12,
xtick={0,2,4,6,8,10,12,14,16,18,20},
ytick={0,2,4,6,8,10,12},
grid=major,
grid style={popline},
tick label style={font=\small},
label style={font=\small},
axis line style={popdark},
]
\addplot[popcurve, very thick, mark=*, mark size=2pt, mark options={fill=popblue}] coordinates {
(0,5) (2,5) (4,5) (6,5.8) (8,7) (10,8.2) (12,9.4) (14,10) (16,10) (18,10) (20,5)
};
\end{axis}
\end{tikzpicture}
\legende{Variation de la quantité d'ADN par noyau au cours du temps (culture synchronisée).}
\end{popfigure}

\item \textbf{Délimitation des phases :}
\begin{itemize}
\item \textbf{G1 :} 0 -- 4 h. ADN constant à 5 u.a. (pas de synthèse).
\item \textbf{S :} 4 -- 14 h. ADN augmente de 5 à 10 u.a. (réplication).
\item \textbf{G2 :} 14 -- 18 h. ADN constant à 10 u.a. (pré-mitose).
\item \textbf{M :} 18 -- 20 h. Chute de 10 à 5 u.a. (division nucléaire + cytocinèse).
\end{itemize}

\item \textbf{Durée phase S :} $14 - 4 = \mathbf{10 \text{ h}}$.
\textbf{Durée totale du cycle :} $20 - 0 = \mathbf{20 \text{ h}}$ (retour à la valeur initiale G1 au temps 20 h).

\item \textbf{Calcul effectif théorique à 60 h :}
Nombre de cycles = $\frac{60}{20} = 3$ cycles complets.
$N_{\text{théo}} = N_0 \times 2^3 = 50\,000 \times 8 = \mathbf{400\,000 \text{ cellules}}$.

\item \textbf{Hypothèses pour l'écart (300 000 observés vs 400 000 théoriques) :}
\begin{enumerate}
\item \textbf{Mortalité cellulaire / apoptose :} une fraction des cellules meurt au cours de la culture (épuisement nutriments, accumulation déchets, densité).
\item \textbf{Asynchronie / sortie de cycle :} toutes les cellules ne se divisent pas à chaque cycle ; certaines entrent en phase G0 (quiescence) ou ont une durée de cycle plus longue.
\end{enumerate}
(Autres acceptables : limitation par contact, épuisement facteurs de croissance, contamination.)

\item La réplication de l'ADN en phase S est indispensable car elle assure la \textbf{duplication fidèle du génome entier} avant la division. Grâce au mode semi-conservatif, chaque cellule fille hérite d'une copie complète et identique de chaque molécule d'ADN parentale. Cela garantit la transmission intégrale de l'information génétique (allèles, gènes, régulation) et donc le maintien du phénotype et de l'identité biologique de la lignée cellulaire au fil des générations. Sans réplication, le génome serait dilué par deux à chaque division, entraînant la perte de l'identité cellulaire.
\end{enumerate}

\textbf{Barème indicatif (10 pts) :}
\begin{itemize}
\item Q1 : 2 pts (courbe correcte : axes, échelles, points, trait)
\item Q2 : 2 pts (4 phases bien délimitées + justifications)
\item Q3 : 1 pt (durée S = 10 h, cycle = 20 h)
\item Q4 : 1,5 pt (calcul $50\,000 \times 2^3 = 400\,000$)
\item Q5 : 1,5 pt (2 hypothèses pertinentes distinctes)
\item Q6 : 2 pts (réplication $\rightarrow$ copie complète $\rightarrow$ transmission fidèle $\rightarrow$ identité conservée)
\end{itemize}

\textbf{Points de vigilance :} Pour Q2, justifier chaque frontière par la variation de l'ADN (constant = pas réplication ; croissant = réplication ; chute = division). Pour Q4, vérifier que 60 h est bien un multiple de la durée du cycle (3 cycles). Pour Q6, utiliser les termes « réplication semi-conservative », « génome complet », « fidélité ».
\end{corrige}

\begin{corrige}[Exercice 11 — Type Probatoire : tâche complexe — extrait de plante]
\begin{enumerate}
\item \textbf{Courbes comparées} (voir figure).
\begin{popfigure}
\begin{tikzpicture}
\begin{axis}[
width=\linewidth,
height=8cm,
xlabel={Temps (h)},
ylabel={Nombre de cellules ($\times 10^3$)},
xmin=0, xmax=96,
ymin=0, ymax=900,
xtick={0,24,48,72,96},
ytick={0,100,200,300,400,500,600,700,800,900},
grid=major,
grid style={popline},
tick label style={font=\small},
label style={font=\small},
axis line style={popdark},
legend pos=north west,
legend style={font=\small, fill=popcream, draw=popdark},
]
\addplot[popblue, very thick, mark=square*, mark size=3pt] coordinates {
(0,50) (24,100) (48,200) (72,400) (96,800)
};
\addlegendentry{Culture témoin}
\addplot[poporange, very thick, mark=triangle*, mark size=3pt] coordinates {
(0,50) (24,62) (48,70) (72,72) (96,71)
};
\addlegendentry{Culture traitée}
\end{axis}
\end{tikzpicture}
\legende{Évolution du nombre de cellules en culture témoin (croissance exponentielle) et traitée par l'extrait de plante (blocage de la prolifération).}
\end{popfigure}

\item \textbf{Comparaison :} La culture témoin présente une \textbf{croissance exponentielle} (doublement tous les 24 h), caractéristique d'une prolifération non limitée. La culture traitée montre une \textbf{forte inhibition de la prolifération} : légère augmentation initiale (peut-être fin de cycle en cours), puis plateau rapide et légère décroissance finale. L'extrait exerce un effet \textbf{cytostatique} (arrêt de la division) voire cytotoxique (mort cellulaire).

\item \textbf{Indice de multiplication (96 h / 0 h) :}
\begin{itemize}
\item Témoin : $\frac{800}{50} = \mathbf{16}$.
\item Traitée : $\frac{71}{50} = \mathbf{1,42}$.
\end{itemize}

\item \textbf{Observation microscopique (48 h) :} Chromosomes alignés sur la plaque équatoriale + fuseau mal formé $\rightarrow$ blocage en \textbf{Métaphase}.
\textbf{Hypothèse de mécanisme :} L'extrait contient une substance qui \textbf{perturbe la polymérisation des microtubules} (comme la colchicine ou la vinblastine) ou leur dépolymérisation (comme le taxol), empêchant la formation d'un fuseau mitotique fonctionnel. Cela active le \textbf{point de contrôle du fuseau (SAC - Spindle Assembly Checkpoint)} qui bloque la cellule en métaphase en inhibant l'APC/C, empêchant ainsi la séparation des chromatides sœurs et la progression en anaphase.

\item \textbf{Expérience complémentaire :} Réaliser une \textbf{immunofluorescence anti-tubuline} sur cultures témoin et traitées pour visualiser l'organisation du fuseau mitotique (microtubules). Comparer la morphologie du fuseau (bipolaire normal vs monopolaire, absent, ou multipolaire). Alternativement : test de réversibilité (lavage de l'extrait à 48 h et suivi de la reprise des divisions).

\item \textbf{Intérêt médical et précaution :}
\begin{itemize}
\item \textbf{Intérêt :} Cet extrait pourrait être une source de \textbf{nouveaux agents chimiothérapeutiques anti-mitotiques} (comme le paclitaxel/Taxol\footnote{Issu de l'if, \emph{Taxus brevifolia}} ou la vincristine), capables d'arrêter spécifiquement la prolifération des cellules cancéreuses qui se divisent plus vite que les cellules saines.
\item \textbf{Précaution indispensable :} Évaluer la \textbf{toxicité sélective} : tester l'effet sur des cellules saines en division rapide (moelle osseuse, épithélium intestinal, follicules pileux) et sur des modèles animaux (pharmacocinétique, toxicité aiguë/chronique, tératogénicité) avant tout essai clinique. Un bon anticancéreux doit avoir un \textbf{indice thérapeutique} (dose toxique / dose efficace) suffisant.
\end{itemize}
\end{enumerate}

\textbf{Barème indicatif (10 pts) :}
\begin{itemize}
\item Q1 : 2 pts (deux courbes sur même graphique, axes légendés, légende, points corrects)
\item Q2 : 2 pts (description correcte : exponentiel vs plateau/inhibition + terme cytostatique)
\item Q3 : 1 pt (calculs exacts : 16 et 1,42)
\item Q4 : 2 pts (identification métaphase + hypothèse précise sur microtubules/fuseau/SAC)
\item Q5 : 1,5 pt (expérience pertinente : immunofluorescence tubuline ou réversibilité)
\item Q6 : 1,5 pt (intérêt : agent anti-mitotique ; précaution : toxicité sur cellules saines / modèle animal)
\end{itemize}

\textbf{Points de vigilance :} Pour Q4, ne pas dire « bloque l'ADN » mais « bloque le fuseau / les microtubules / le checkpoint SAC ». Pour Q6, distinguer « intérêt » (potentiel thérapeutique) et « précaution » (évaluation toxicologique rigoureuse, sélectivité).
\end{corrige}

\newpage

\coursec{Fiche bilan}

\courssub{Carte mentale de la leçon}
\begin{popfigure}
\begin{tikzpicture}[
  mindmap,
  grow cyclic,
  every node/.style={concept, execute at begin node=\hskip0pt},
  concept color=popblue,
  level 1/.append style={level distance=4.5cm, sibling angle=360/6},
  level 2/.append style={level distance=3cm, sibling angle=360/6},
  texte/.style={text width=2.8cm, align=center, font=\small\sffamily},
  mitose/.style={concept color=poporange},
  adn/.style={concept color=popgreen},
  cycle/.style={concept color=poppurple},
  caryo/.style={concept color=poppink},
  cancer/.style={concept color=popdark},
  exp/.style={concept color=popgold}
]
\node[texte, minimum size=2.2cm, align=center]{Le cycle cellulaire \& la mitose\\Maintien de l'identité biologique}
  child[mitose] { node[texte, align=center]{Phases de la mitose\\Prophase, Métaphase, Anaphase, Télophase} }
  child[adn] { node[texte, align=center]{Réplication de l'ADN\\Semi-conservative, fidélité} }
  child[cycle] { node[texte, align=center]{Cycle cellulaire\\Interphase (G1, S, G2) \& Phase M} }
  child[caryo] { node[texte, align=center]{Comportement des chromosomes\\Conservation du caryotype} }
  child[cancer] { node[texte, align=center]{Dysfonctionnements\\Cancérisation, perte de contrôle} }
  child[exp] { node[texte, align=center]{Données expérimentales\\Cinétique de prolifération, index mitotique} };
\end{tikzpicture}
\legende{Carte mentale résumant les 6 piliers de la leçon ND01 : cycle cellulaire, réplication ADN, phases mitotiques, conservation du caryotype, pathologies (cancer) et analyse de données.}
\end{popfigure}

\courssub{Bilan synthétique}

\begin{bilan}

\textbf{Définitions clés à connaître par cœur}
\begin{itemize}
  \item \cle{Cycle cellulaire} : Enchaînement ordonné des événements conduisant une cellule à se diviser en deux cellules filles. Comprend l'\cle{interphase} (G1, S, G2) et la \cle{phase M} (mitose + cytodiérèse).
  \item \cle{Mitose} : Division nucléaire conservant le nombre de chromosomes (2n $\to$ 2n) et l'information génétique. Quatre phases continues : prophase, métaphase, anaphase, télophase.
  \item \cle{Réplication de l'ADN} : Processus \emph{semi-conservatif} (Meselson-Stahl) durant la phase S. Chaque brin sert de matrice ; la nouvelle double hélice contient un brin parental et un brin néo-synthétisé.
  \item \cle{Caryotype} : Formule chromosomique d'une espèce (nombre, forme, taille des chromosomes). Chez l'humain : 2n = 46 (22 paires d'autosomes + 1 paire de gonosomes XX ou XY).
  \item \cle{Index mitotique (IM)} : Pourcentage de cellules en division (phase M) dans une population cellulaire. $IM = \frac{N_{\text{mitoses}}}{N_{\text{total}}} \times 100$.
  \item \cle{Cancérisation} : Processus multi-étapes transformant une cellule saine en cellule cancéreuse (prolifération anarchique, perte de contact inhibition, immortalisation, angiogenèse, métastases).
\end{itemize}

\textbf{Relations fondamentales \& repères chiffrés}
\begin{center}
\begin{tabularx}{\linewidth}{|l|X|X|}\hline
\rowcolor{popblueL}
\textbf{Notion} & \textbf{Énoncé / Formule} & \textbf{Point de vigilance Probatoire} \\ \hline
Durée du cycle & $T_{\text{cycle}} = T_{\text{G1}} + T_{\text{S}} + T_{\text{G2}} + T_{\text{M}}$ & $T_{\text{M}} \ll T_{\text{interphase}}$ (ex. 1 h vs 20 h). \\ \hline
Réplication ADN & \ce{2 ADN (1 brin parental + 1 brin neuf)} & Semi-conservative $\neq$ conservative ni dispersive. \\ \hline
Conservation caryotype & $2n \xrightarrow{\text{mitose}} 2n + 2n$ & Vérifier : nombre de chromosomes \textbf{après} télophase = nombre \textbf{avant} prophase. \\ \hline
ADN / Chromosome & 1 chromosome = 1 molécule d'ADN (condensée) & En prophase/métaphase : 1 chr = 2 chromatides = 2 molécules d'ADN. \\ \hline
Index mitotique & $IM = \frac{N_M}{N_{\text{tot}}} \times 100$ & Compter \textbf{toutes} les phases mitotiques (P+M+A+T), pas seulement métaphase. \\ \hline
\end{tabularx}
\end{center}

\textbf{Méthodes express}
\begin{methode}[Identifier une phase mitotique sur lame / photo]
\begin{enumerate}
  \item \textbf{Prophase :} Chromosomes condensés (bâtonnets en X), noyau disparu, fuseau en formation.
  \item \textbf{Métaphase :} Chromosomes alignés sur le \cle{plan équatorial}, centromères sur la plaque, chromatides bien distinctes.
  \item \textbf{Anaphase :} Séparation des chromatides (devenues chromosomes filles) migrent vers les pôles ; zone équatoriale vide.
  \item \textbf{Télophase :} Chromosomes décondensés aux pôles, réapparition des enveloppes nucléaires, début cytodiérèse.
\end{enumerate}
\end{methode}

\begin{methode}[Calculer l'index mitotique et interpréter]
\begin{enumerate}
  \item Compter $N_{\text{total}}$ noyaux sur plusieurs champs représentatifs (min 500-1000 cellules).
  \item Compter $N_M$ noyaux en mitose (toutes phases confondues).
  \item Appliquer la formule. IM élevé $\Rightarrow$ tissu en forte croissance (méristème, embryon, tumeur). IM $\approx$ 0 $\Rightarrow$ tissu adulte stable (nerf, muscle strié).
\end{enumerate}
\end{methode}

\begin{methode}[Expliquer la conservation de l'information génétique]
\begin{enumerate}
  \item Phase S : Réplication fidélisée par ADN polymérase (lecture 5'$\to$3', activité exonucléasique 3'$\to$5' = \cle{proofreading}).
  \item Mitose : Ségrégation équitable des chromatides sœurs (identiques) vers les deux pôles $\to$ deux noyaux fils \textbf{génétiquement identiques} au noyau mère et entre eux.
\end{enumerate}
\end{methode}

\end{bilan}

\courssub{Checklist avant le Probatoire}
\begin{aretenir}[Checklist avant le Probatoire]
  $\square$ Je définis sans hésiter : cycle cellulaire, interphase (G1, S, G2), mitose, cytodiérèse, caryotype, chromatide, centromère, fuseau mitotique, index mitotique.\par
  $\square$ Je place la réplication de l'ADN dans la \textbf{phase S} de l'interphase et j'explique le mode \textbf{semi-conservatif} (expérience de Meselson \& Stahl).\par
  $\square$ Je reconnais et je décris les \textbf{4 phases de la mitose} (P, M, A, T) sur schéma ou photo microscopique (critères morphologiques précis).\par
  $\square$ Je relie le comportement des chromosomes (condensation, alignement, séparation chromatides) à la \textbf{conservation du caryotype} (2n $\to$ 2n).\par
  $\square$ Je distingue \textbf{chromosome} (unité structurale) et \textbf{molécule d'ADN} (unité chimique) : 1 chr = 1 ADN (G1) ; 1 chr = 2 ADN (après réplication, jusqu'à anaphase).\par
  $\square$ Je calcule l'\textbf{index mitotique} à partir d'un dénombrement et j'interprète sa valeur (croissance, régénération, pathologie).\par
  $\square$ J'explique le lien \textbf{dysfonctionnement mitose / cancérisation} : points de contrôle (checkpoints) G1/S, G2/M, métaphase/anaphase ; rôle de p53, cyclines, CDK ; caractères des cellules cancéreuses.\par
  $\square$ Je sais utiliser un \textbf{tableau de variation} ou un \textbf{graphique} (cinétique de prolifération, courbe de croissance tumorale) pour extraire des informations (durée des phases, effet d'un inhibiteur).\par
  $\square$ Je cite au moins \textbf{2 exemples camerounais} : culture de méristèmes (bananier, manioc) pour l'agriculture ; dépistage cervical (frottis, index mitotique) ou lutte contre le cancer au Centre Hospitalier de Yaoundé.\par
  $\square$ Je rédige une \textbf{argumentation structurée} (Observation $\to$ Hypothèse $\to$ Expérience/Données $\to$ Interprétation $\to$ Conclusion) pour toute question de type « expliquer » ou « montrer que ».\par
  $\square$ Je vérifie mes \textbf{unités} (durées en heures/minutes, index en \%), mes \textbf{notations} (2n, n, c pour quantité d'ADN) et la \textbf{précision du vocabulaire} (chromatide $\neq$ chromosome, réplication $\neq$ transcription).
\end{aretenir}

\findecours

\end{document}
